% French translation of paper.tex (best effort, made with AI help; the English
% version is authoritative). Same structure, numbering and labels as paper.tex.
% Build: `make paper NOTE=turns-of-a-t-mesh` from the root of the repository,
% which compiles paper.tex and this file and first writes version.tex; or by
% hand: pdflatex paper.fr && bibtex paper.fr && pdflatex paper.fr && pdflatex paper.fr
%
% French typography by hand, as in closing-loops/paper.fr.tex: babel loads
% French from its .ini locale, the same with or without the babel-french
% package, and that locale does not space the punctuation under pdflatex. So
% the source writes the spaces itself: `~:` before a colon, `\,;` `\,?` `\,!`
% before the others, `«~...~»`, and `\,\%`. Decimal comma in math through
% icomma: `$0,5$`, but `$1, \dots$` stays a list.
%
% Terms: "ring" is « anneau » and "annulus" « couronne »; "positive" (> 0) is
% « strictement positif » and "non-negative" « positif ou nul »; "loop" (a
% closed curve) is « lacet », "edge loop" « boucle d'arêtes ».
\documentclass[11pt,a4paper]{article}

\usepackage[T1]{fontenc}
\usepackage[utf8]{inputenc}
\usepackage{lmodern}
\usepackage{babel}
\babelprovide[main,import]{french}
\usepackage{icomma}
\usepackage[margin=2.5cm]{geometry}
\usepackage{amsmath,amssymb,amsthm}
\usepackage{booktabs}
\usepackage{array}
\usepackage[font=small]{caption}
\usepackage{tikz}
\usetikzlibrary{arrows.meta,decorations.pathreplacing,patterns}
\usepackage[numbers,sort&compress]{natbib}
\usepackage[hidelinks]{hyperref}
\setlength{\emergencystretch}{2em}
\renewcommand{\topfraction}{0.9}
\renewcommand{\floatpagefraction}{0.8}
% The references in a smaller type.
\renewcommand{\bibfont}{\small}
\setlength{\bibsep}{2pt plus 1pt}
\setlocalecaption{french}{table}{Tableau}

\newtheorem{theorem}{Théorème}
\newtheorem{definition}{Définition}
\newtheorem{proposition}{Proposition}
\newtheorem{lemma}{Lemme}
\newtheorem{corollary}{Corollaire}
\newtheorem{remark}{Remarque}
\newtheorem{conjecture}{Conjecture}
\newtheorem{question}{Question}

\newcommand{\Z}{\mathbb{Z}}
\newcommand{\R}{\mathbb{R}}
\newcommand{\shift}{s}
\newcommand{\RC}{\textup{\textsc{Ring-Closing}}}

% The commit of the Lean proofs described, written by `make paper` into
% version.tex (ignored by git).
\IfFileExists{version.tex}{\input{version.tex}}{}
\providecommand{\enginecommit}{unknown}

% Levels of evidence: each statement and table is framed with a tag naming
% what checks it (Lean, the engine's data, a measure, the proof on paper, a
% citation), read from lean/README.md. evidence.py, beside this file, checks
% that every statement has a level and that every Lean name cited is
% declared in lean/. \def\plain{} before % Research note: turns of a T-mesh, obstructions to closing edge loops (the
% theory paper; the applied companion is docs/research/closing-loops). The Lean
% project that checks its algebraic cores is in lean/, beside this file.
% Build: `make paper NOTE=turns-of-a-t-mesh` from the root of the repository,
% which first writes version.tex (the commit described, \enginecommit);
% without it, or by hand: pdflatex paper && bibtex paper && pdflatex paper && pdflatex paper
\documentclass[11pt,a4paper]{article}

\usepackage[T1]{fontenc}
\usepackage[utf8]{inputenc}
\usepackage{lmodern}
\usepackage[margin=2.5cm]{geometry}
\usepackage{amsmath,amssymb,amsthm}
\usepackage{booktabs}
\usepackage{array}
\usepackage[font=small]{caption}
\usepackage{tikz}
\usetikzlibrary{arrows.meta,decorations.pathreplacing,patterns}
\usepackage[numbers,sort&compress]{natbib}
\usepackage[hidelinks]{hyperref}
\setlength{\emergencystretch}{2em}
\renewcommand{\topfraction}{0.9}
\renewcommand{\floatpagefraction}{0.8}
% The references in a smaller type.
\renewcommand{\bibfont}{\small}
\setlength{\bibsep}{2pt plus 1pt}

\newtheorem{theorem}{Theorem}
\newtheorem{definition}{Definition}
\newtheorem{proposition}{Proposition}
\newtheorem{lemma}{Lemma}
\newtheorem{corollary}{Corollary}
\newtheorem{remark}{Remark}
\newtheorem{conjecture}{Conjecture}
\newtheorem{question}{Question}

\newcommand{\Z}{\mathbb{Z}}
\newcommand{\R}{\mathbb{R}}
\newcommand{\shift}{s}
\newcommand{\RC}{\textup{\textsc{Ring-Closing}}}

% The commit of the Lean proofs described, written by `make paper` into
% version.tex (ignored by git).
\IfFileExists{version.tex}{\input{version.tex}}{}
\providecommand{\enginecommit}{unknown}

\title{Turns of a T-Mesh:\\
Obstructions to Closing Edge Loops\\[0.5em]
\large Research note --- draft, not peer reviewed}
\author{Marc Gardent%
\thanks{This note is distributed under the Creative Commons Attribution
4.0 International licence (CC BY 4.0). The code of openclare and the Lean
proofs are under the GPL-3.0-or-later of the repository. The note describes
the Lean proofs at commit \texttt{\enginecommit} of
\url{https://github.com/marcgardent/openclare}.
A French translation is in \texttt{paper.fr.pdf}, beside this note.}\\
\normalsize Independent researcher\\
\normalsize ORCID: \href{https://orcid.org/0009-0000-6466-3098}{0009-0000-6466-3098}\\
\normalsize Contact: through the repository,
\url{https://github.com/marcgardent/openclare}}
\date{October 2026}

\begin{document}
\maketitle

\begin{abstract}
In the quantization of a T-mesh, an edge loop closes into a ring when a
\emph{turn}, a straight walk across the patches that returns to its first
side, has zero \emph{shift}, an integer linear form in the arc lengths;
otherwise it winds as a helix. We study why turns cannot all be closed.
Reading the shift as translational holonomy of the flat metric, we show that
for a simple turn it is the across component of the Burgers vector of the
singular vertices on one side; that a 3--5 pair or an exchange of two
vertices, joined by a staircase of arcs, forces one of two turns to be
shifted, and that a separatrix winding once round a turn and stopping on
another separatrix of its vertex forces that turn alone to be shifted; and
that on a
closed sphere the shift and the length of a turn have the same parity. Fox calculus extends the Burgers sum to
any combination of turns, and \emph{combs}, local dislocations carried by
single arcs, write every certificate of a conflict found on a human body. The
conflicts bound from below the number of turns that every quantization within
given bounds leaves shifted: on that body, within the bounds of separation of
a remeshing engine, exactly one to nine of the turns it holds, where it leaves
13 to 123 shifted. Choosing the turns to close is NP-hard for bidirected
balances with $0$--$1$ bounds, the algebraic form of quantization; whether
actual T-meshes give hard instances is open. The algebraic cores
of the proofs are checked in Lean; the topology of the surface, that the walks
and regions are those of a T-mesh, is a hypothesis there.
\end{abstract}

\section{Introduction}

Automatic quad remeshing often produces edge loops that wind round a part of
the shape as helices instead of closing into
rings~\cite{bommes2011structure}. Whether a loop closes is decided when a
quantizer assigns whole lengths to the arcs of a T-mesh traced in a seamless
parametrization~\cite{campen2015qgp,lyon2021constrained,heistermann2023bimdf}:
a loop follows a \emph{turn}, a straight walk across the patches that returns
to its first side, and closes when the \emph{shift} of the turn, an integer
linear form in the arc lengths, vanishes. A companion
paper~\cite{gardent2026loops} adds these equations to a quantizer; on real
models they cannot all hold together with the bounds that keep singular
vertices apart. This paper asks why, and what follows.

Sections~\ref{sec:setting} to~\ref{sec:beyond} go from the holonomy of one
turn to conflicts between turns, any combination of turns, and borders and
handles; Sections~\ref{sec:forced} and~\ref{sec:choosing} bound the turns
every quantization leaves shifted and show that choosing them is NP-hard;
Sections~\ref{sec:lean} and~\ref{sec:open} say what is machine-checked and
what is open. The ingredients are classical (holonomy of flat cone metrics, Burgers vectors,
the parity of quadrangulations, Fox calculus, covers of infeasible
subsystems); their statements for turns, the lower bound on helices and the
hardness for the instances of quantization are ours, as far as we found.
Statements are \emph{proved} when they come with a proof, \emph{measured}
when they come from exact arithmetic or solver output on the systems the
remesher openclare exports, and \emph{conjectured} otherwise.

This work began as a programming project, openclare, an open-source quad
remesher, and became research when engineering alone could no longer tell
which results could be guaranteed and which computations saved. The results
here answer from the side of what cannot be done: a repair that keeps
the length of a turn of odd shift cannot close it on a closed sphere; a turn
whose shift is a sign-definite form with an arc held open stays a helix and
need not be tried; the turns every quantization leaves shifted measure what a
heuristic can still gain; and the hardness of choosing them is why an engine
settles for a heuristic.

\section{Related work}\label{sec:related}

The valences of a closed quad mesh of genus $g$ satisfy $\sum_v (4 - k(v)) = 8(1-g)$,
which says how many singular vertices are needed, not where. A closed quad mesh induces a quartic differential
whose zeros and poles are its singular vertices; Lei et
al.~\cite{lei2020quartic} list as open the guarantee that all its
trajectories are finite, and a trajectory that is not finite is an edge loop
that does not close. A quadratic
differential whose trajectories all close is
Jenkins--Strebel~\cite{strebel1984quadratic}, a union of flat cylinders; a
helix is a cylinder glued back with a shift of whole rows. Chen et
al.~\cite{chen2019metric} characterize the flat metrics of quad meshes,
assuming that the lines of the field are finite loops. Izmestiev et al.~\cite{izmestiev2013torus} prove non-existence results for quad meshes,
among them the torus with one vertex of valence three and one of valence
five~\cite{barnette1971toroidal}, with the holonomy of the developing map,
whose translation, when the rotation is trivial, is the Burgers vector of
crystallography; a pair of opposite disclinations is a
dislocation~\cite{kleman2008disclinations}. The periods of a flat cone metric
are local coordinates of it~\cite{thurston1998shapes,veech1993flat}. Bommes et
al.~\cite{bommes2011structure} remove helices from finished meshes; stripe
parameterizations of tubes~\cite{kalberer2011stripe} round the gap round a
tube, our shift, to a whole number. We found no lower bound on the number of helices for a given
singularity configuration; ours is a minimum cover of irreducible infeasible
subsystems~\cite{parker1996iis,amaldi2003maxfs}.

\section{Turns, shifts and holonomy}\label{sec:setting}

A T-mesh is a partition of the surface into rectangles, the \emph{patches},
whose sides are made of \emph{arcs}; nodes where a trace stops on another are
T-junctions, and the singular vertices of the field are nodes. Each patch $P$
has local coordinates with its first corner at the origin and its sides along
the axes. A quantization assigns every arc $a$ a length $\ell_a \in \Z_{\ge
0}$ such that opposite sides of every patch have equal length, the
\emph{balances} $B\ell = 0$; the quad mesh is then the integer grid of each
patch.

\begin{definition}[Turn]
A \emph{turn} $\gamma$ is a straight walk across patches, from an offset of
its first side that meets no node: it enters a patch through a side, leaves
through the opposite side at the same distance from the corresponding corner,
enters the neighbouring patch across the arc it reaches, and so on, until it
comes back into its first patch through its first side. The rotation of the
composed crossings is then the identity, since the walk re-enters its first
side inward.
\end{definition}

Crossing from patch $P$ to patch $Q$ across an arc composes a quarter-turn
rotation of the local frame with a translation whose components are signed
sums of arc lengths along the rims of $P$ and $Q$.

\begin{proposition}[The shift of a turn is linear]\label{prop:linear}
For a turn $\gamma$ there are integer coefficients $c_a$, independent of the
lengths, such that after a full turn a point of the first side has moved
across the walking direction by
\[
  \shift_\gamma(\ell) = \sum_a c_a\, \ell_a .
\]
For integer lengths $\ell$, let $o$ be a half-integer offset of the first side
whose walk at the lengths $\ell$ crosses the same patches as $\gamma$ (we call
the set of such offsets the \emph{corridor} of $\gamma$ at $\ell$). Then the
dual line of the quad mesh at $o$ comes back at $o + \shift_\gamma(\ell)$: it
closes after one round if and only if $\shift_\gamma(\ell) = 0$. If moreover
$o + j\,\shift_\gamma(\ell)$ lies in the corridor for $j = 1, \dots, t - 1$, it
winds $t$ times round as a helix of pitch $|\shift_\gamma(\ell)|$ in the sense
of~\cite{bommes2011structure}. If the corridor holds no half-integer, no strip
of the mesh follows $\gamma$.
\end{proposition}

\begin{proof}[Sketch]
Each crossing is a map $x \mapsto R_k x + t$ with $R_k$ a rotation by $k$
quarter-turns and $t$ a vector whose entries are sums of $\pm\ell_a$. The
composition along $\gamma$ has trivial rotation, hence is a translation
$T_\gamma(\ell)$, linear in $\ell$ with integer coefficients; its component
across the walking direction is $\shift_\gamma(\ell)$. Singularities are
nodes, so the quad mesh is a regular grid along the walk, and a dual line at
distance $o$ from the first corner that crosses the same patches comes back at
$o + \shift_\gamma(\ell)$. The corridor, two inequalities per crossing, may
contain $o$ but not $o + \shift_\gamma(\ell)$; the next round then follows
other patches.
\end{proof}

In the frame of the first side, with $e_1$ the walking direction and $e_2$
the direction of increasing offsets, $T_\gamma(\ell) = -\alpha_\gamma(\ell)\,
e_1 + \shift_\gamma(\ell)\, e_2$: it brings the end of the walk back onto the
first side. $\alpha_\gamma(\ell)$ is the \emph{length} of the turn, which at whole lengths is the number of
quads of the strip in one turn. At the real lengths of the parametrization the
shift is the \emph{real shift}, and its share of the length the \emph{lean}.
At whole balanced lengths with positive patch sides every strip of the quad
mesh is such a walk~\cite[Prop.~2]{gardent2026loops}, and every shift is
whole; closing a ring asks it to be zero.

\paragraph{Flat metric and holonomy.}
Unless stated otherwise the surface $\Sigma$ is closed, and \emph{genus zero}
means a closed sphere. Balanced lengths $\ell > 0$ make each patch a
Euclidean rectangle, and $\Sigma$ a flat surface with cone points at the
singular vertices: a vertex of valence $k$ has angle $k\pi/2$ and curvature
$(4-k)\pi/2$, and T-junctions are not cone points. Let $\Sigma^\circ$ be $\Sigma$ without its singular vertices,
$x_0$ a point inside a patch, and $R$ the positive quarter-turn. Developing a
loop $\gamma$ of $\Sigma^\circ$ based at $x_0$ patch after patch gives its
\emph{holonomy}, a motion $\hat h(\gamma)\colon x \mapsto R^j x + t$ of the
coordinates of the first patch; it depends only on the homotopy class of
$\gamma$, and $\hat h(\alpha\beta) = \hat h(\alpha) \circ \hat
h(\beta)$~\cite{izmestiev2013torus}. A small loop round a singular vertex $c$
of valence $k$, run with $c$ on its right and reached along a path that
develops it at $p$, has holonomy $x \mapsto R^{4-k}(x - p) + p$. When the
rotation of $\hat h(\gamma)$ is the identity, its translation is the
\emph{Burgers vector} of $\gamma$. For a turn $\gamma$, let $\bar\gamma$ be
the closed curve that follows the walk and, back in its first patch, returns
to its start along a segment parallel to the first side; $\hat
h(\bar\gamma)$ is a translation. It continues the chart along the loop,
while the crossings of Proposition~\ref{prop:linear} develop the last patch
into the first: $\hat h(\bar\gamma)$ is the inverse of $T_\gamma(\ell)$, the
translation by $\alpha_\gamma(\ell)\, e_1 - \shift_\gamma(\ell)\, e_2$, and
$\shift_\gamma(\ell) = -\langle e_2, \hat h(\bar\gamma) \rangle$. Every shift
in this paper is that of Proposition~\ref{prop:linear}, the convention of the
shifts the engine exports and of the figures. Shifts, developed positions and Burgers vectors are
integer linear forms in $\ell$ as long as the paths cross the same patches, so
an identity between them that holds on an open set of balanced positive
lengths holds for all balanced lengths; for whole lengths the nodes develop
onto $\Z^2$. A form $f$ equals $g$ \emph{modulo the balances} if $f - g$ is a
combination of the rows of $B$.

\paragraph{Data.}
The measured statements refer to the T-mesh of the male body of Blender's
Essentials library remeshed by openclare at $10{,}590$ quads, exported with
every turn (branch \texttt{rd-exact}, engine before commit \texttt{c8bf2be})
and studied in the companion paper~\cite[\S7]{gardent2026loops}: $261$
singular vertices ($133$ of valence three, $127$ of valence five, one of
valence two) among $1{,}295$ nodes, $2{,}069$ arcs, $776$ patches, $1{,}552$
balances. Turns are found in eight \emph{ways}, way $m + 2j + 4i$ walking from
the shares $\{0.5\}$ ($m = 0$) or $\{0.2, 0.5, 0.8\}$ ($m = 1$) of each side
and \emph{holding} (trying to close) a turn if its real shift is above $0.05$,
its lean below $0.05$ ($j = 0$) or $0.12$ ($j = 1$) and its length below $2$
($i = 0$) or $6$ ($i = 1$) sides of the T-mesh; the engine keeps way~2. The
ways find $1{,}299$ turns and hold $765$ (summed over the ways). The integer
search holds $28$ to $36$ arcs at length at least one, the \emph{bounds of
separation}; $P_{\mathrm{eng}}$, the
balanced lengths meeting them, is neither a subset nor a superset of the valid
quantizations, which only ask a chain between two singular vertices to keep a
positive length. Identities are checked in floating point unless said exact;
programs were solved with HiGHS 1.8.0~\cite{huangfu2018highs} through SciPy
1.16.3~\cite{virtanen2020scipy}, on an Intel Core i7-7700K (4 cores, 2017,
62\,GB) shared with other jobs.

\section{Turns and singularities}\label{sec:holonomy}

\begin{lemma}[Gauss--Bonnet for turns]\label{lem:gauss}
Let $\gamma$ be a turn whose curve $\bar\gamma$ is simple. If $\bar\gamma$
bounds a disc $D \subset \Sigma$, the singular vertices inside $D$ satisfy
$\sum_{c \in D} (4 - k(c)) = 4$. If the curves of two turns are disjoint and
bound an annulus $A$, then $\sum_{c \in A} (4 - k(c)) = 0$.
\end{lemma}

\begin{proof}
For a region of a flat surface with cone points bounded by straight segments,
Gauss--Bonnet reads $\sum_c \kappa_c + \sum_i \theta_i = 2\pi\chi$, with
$\kappa_c = (4 - k(c))\pi/2$ the curvatures inside and $\theta_i$ the turning
angles of the border~\cite{bommes2013survey}. The curve $\bar\gamma$ turns
only at the ends of its closing segment, by $+\pi/2$ and $-\pi/2$, so the
turning angles sum to zero on each curve, seen from either side; $\chi(D) = 1$
and $\chi(A) = 0$.
\end{proof}

On a closed sphere a simple turn thus splits the singular vertices into two
groups of total index four, each of at least two vertices: a fingertip beyond
a ring holds the equivalent of four vertices of valence three.

\begin{lemma}[One self-crossing]\label{lem:eight}
Let the surface be closed of genus zero, and let the curve $\bar\gamma$ of a
turn cross itself exactly once, transversally, with no singular vertex on it.
It cuts the surface into three discs: two \emph{lobes}, each bounded by one of
the two loops of $\bar\gamma$, and a third region bounded by both. If the
crossing lies on the walk, the totals of $4 - k$ over the singular vertices of
the lobes and the third region are $3$, $3$ and $2$; if it lies on the closing
segment, they are $2$, $2$ and $4$.
\end{lemma}

\begin{proof}
The walk crosses itself at a right angle, so each loop leaves the crossing
$X$ along one passage and comes back along the other, its two rays a
quarter-turn apart. Each loop is simple and bounds a disc containing the
quadrant between its rays, a lobe with a corner of angle $\pi/2$ at $X$; the
lobes take opposite quadrants and the third region the other two, and the
three faces of this one-vertex graph on the sphere are discs. By
Gauss--Bonnet a disc whose border turns by $\theta_i$ holds total index $4 -
\sum_i \theta_i/(\pi/2)$, which gives $3$, $3$ and $2$ before the corners of
the closing segment, of turnings $+\pi/2$ and $-\pi/2$ in succession. If the
crossing lies on the walk, both are on one loop, every region sees both, and
they cancel. Otherwise let the walk start at $A$ along $e_1$, come back at $Z$
and cross the closing segment between them; it crosses along $-e_1$, since
along $e_1$ the turn would have ended there. The loop through $A$ then turns
the same way at $A$ as at $X$, and so does the loop through $Z$ at $Z$ and at $X$ (a
check of the four directions): each lobe sees a second convex corner and holds
$2$, and the third region the rest of the total $8$.
\end{proof}

\begin{proposition}[The shift is set by the singularities]\label{prop:burgers}
Let $\gamma$, $\gamma_1$, $\gamma_2$ be turns with simple curves.
\begin{enumerate}
\item[(a)] Let $\bar\gamma$ bound a disc $D$ holding the singular vertices
  $c_1, \dots, c_m$, and let $x_0$ be a point of $\bar\gamma$, the start of the
  walk for instance. Let $T \subset D$ be a tree whose vertices include $c_1,
  \dots, c_m$ and which meets $\bar\gamma$ only at $x_0$. Cut along $T$, the
  disc $D$ becomes a disc $\Omega$ without cone point, which develops into the
  plane, with the frame $(e_1, e_2)$ at $x_0$. Number the $c_i$ in the order in
  which a walk round $T$ from $x_0$, keeping $T$ on its right, first reaches
  them, and let $p_i$ be the developed position of $c_i$ at that first visit.
  Then $\shift_\gamma(\ell) = \langle e_2, B_D \rangle$ if $D$ lies on the
  left of $\bar\gamma$, and $\shift_\gamma(\ell) = -\langle e_2, B_D
  \rangle$ if it lies on its right, with
  \[
    B_D = \sum_{i=1}^m Q_i\,(I - R_i)\,p_i ,
    \qquad R_i = R^{4 - k(c_i)},
    \qquad Q_i = R_m R_{m-1} \cdots R_{i+1},
  \]
  the translation of $\rho_m \circ \dots \circ \rho_1$, where $\rho_i(x) =
  R_i(x - p_i) + p_i$. With $D$ on its left, $B_D = T_\gamma(\ell) =
  -\alpha_\gamma(\ell)\, e_1 + \shift_\gamma(\ell)\, e_2$; with $D$ on its
  right, $B_D$ is the inverse translation.
\item[(b)] Let $\bar\gamma_1$ and $\bar\gamma_2$ be disjoint and bound an
  annulus $A$, round which the two turns go the same way, and let a path
  $\kappa$ inside $A$ from $\bar\gamma_1$ to $\bar\gamma_2$ carry the walking
  direction of $\gamma_2$ onto that of $\gamma_1$ (the turns are
  \emph{parallel}); measure both shifts along the direction $e_2$ carried by
  $\kappa$. Then
  \(
    \shift_{\gamma_1}(\ell) - \shift_{\gamma_2}(\ell) = -\langle e_2, B_A \rangle,
  \)
  where $B_A$ is the Burgers vector of the border of the disc $A$ cut along
  $\kappa$, given by the same sum over the singular vertices inside $A$,
  developed inside the cut disc.
\item[(c)] In particular, parallel turns bounding an annulus without singular
  vertex have the same shift for every $\ell$: a helix is a property of a
  whole cylinder of parallel turns, not of one of them.
\end{enumerate}
\end{proposition}

\begin{figure}[t]
\centering
% The shift of a turn is set by the singular vertices on one side of it
% (Proposition 2, "The shift is set by the singularities").
% \input inside a figure; needs \usepackage{tikz} and
% \usetikzlibrary{arrows.meta,decorations.pathreplacing}.
%
% A cap D of a T-mesh (a fingertip, say), cut along a tree T into a disc
% Omega that develops into the plane; every node is at a whole position and the
% faint grid is the quad mesh at these lengths. Omega is the band [0,11]x[.,4]
% (it goes on below), the block [3,11]x[4,5] and the top, the L-shaped
% [0,3]x[4,6] + [0,2]x[6,7]. T runs from x_0 up the left side to c_1 and then
% round the border of the top: c_1, c_2, ..., c_6, which are the six singular
% vertices of D. The cut sides are glued in pairs: the right side x = 11 to the
% left side one unit lower, (11,y) ~ (0,y-1); and the top of the block, from
% c_6 = (3,5) rightwards, onto the border of the top, turning round it:
% [3,4] onto [(3,5),(3,6)], [4,5] onto [(3,6),(2,6)], [5,6] onto [(2,6),(2,7)],
% [6,8] onto [(2,7),(0,7)], [8,11] onto [(0,7),(0,4)]. So c_1 = (11,5) = (0,4),
% c_2 = (8,5) = (0,7), c_3 = (6,5) = (2,7), c_4 = (5,5) = (2,6), c_5 = (4,5) =
% (3,6), c_6 = (3,5). Angles: 90 + 90 + 90 at c_1; 90 + 180 at c_2, c_3, c_5;
% 180 + 90 at c_6; 270 + 180 at c_4. Valences 3, 3, 3, 5, 3, 3; total index
% 5 - 1 = 4 (Lemma 1, a disc).
%
% Traces besides the cut: c_1 to the right (stops at (2,4)), c_4 to the left
% (stops on the cut at (0,6) = (9,5)) and downwards (x = 2), c_6 to the left
% (stops at (2,5)) and downwards (x = 3), and the separatrices downwards from
% c_5, c_4, c_3, c_2 at x = 4, 5, 6, 8.
%
% The turn gamma from x_0 = (0, 3.5) crosses x = 2, 3, 4, 5, 6, 8, 11 and comes
% back at (0, 2.5): length 11, s = -1 (one row lower; offsets measured
% upwards). The closing segment of its curve runs up the first side from 2.5
% to 3.5; continued along the walk, the chart brings x_0 back at x_0' =
% (11, 4.5), where the walk round T (keeping T on its right) starts.
% First visits on that walk, with R the quarter-turn x -> (-y, x), R_i = R for
% valence 3 and R^{-1} for valence 5, Q_i = R_6 ... R_{i+1}:
%   p_1 = (11,5) v3  (I-R)p = (16,-6)  Q_1 = R^3  -> ( -6,-16)
%   p_2 = ( 8,5) v3  (I-R)p = (13,-3)  Q_2 = R^2  -> (-13,  3)
%   p_3 = ( 6,5) v3  (I-R)p = (11,-1)  Q_3 = R    -> (  1, 11)
%   p_4 = ( 5,5) v5  (I-R^-1)p = (0,10) Q_4 = R^2 -> (  0,-10)
%   p_5 = ( 4,5) v3  (I-R)p = ( 9, 1)  Q_5 = R    -> ( -1,  9)
%   p_6 = ( 3,5) v3  (I-R)p = ( 8, 2)  Q_6 = I    -> (  8,  2)
%   B_D = sum = (-11, -1) = x_0 - x_0'.
% D lies on the left of the walk, so B_D = T_gamma = (-alpha, s): minus the
% length 11 along the walk, the shift -1 across it, s = <e_2, B_D>
% (Proposition 2(a)). The other order of visit gives
% (1,1), which is not a translation of x_0: the order matters.
% gamma' from (0, 1.5) comes back at (0, 0.5): s = -1, length 11. The curves of
% gamma and gamma' bound an annulus without singular vertex (shaded), and the
% shifts agree (Proposition 2(c)).
% Parity on a sphere (Proposition 5): 11 = -1 mod 2.
% check.py builds this surface and checks these numbers by walking the strips
% of the quad mesh: the valences above, every other vertex of valence 4,
% s = -1 and length 11 from o = 1.5, 2.5, 3.5, and the sum B_D.
\begin{tikzpicture}[x=0.55cm, y=0.55cm, font=\footnotesize,
    side/.style={draw=black, line width=0.5pt},
    quadgrid/.style={draw=black!18, line width=0.25pt, step={(1,1)}},
    surface/.style={fill=black!3},
    node dot/.style={circle, fill=black, inner sep=0pt, minimum size=2.6pt},
    turn/.style={line width=1pt, -{Stealth[length=5pt]}},
    closing/.style={line width=0.7pt, densely dotted},
    cut/.style={draw=black, line width=1.1pt, line join=round},
    brace/.style={decorate, decoration={brace, amplitude=3pt}, line width=0.5pt},
    burgers/.style={draw=measured, line width=1.1pt, -{Stealth[length=6pt]}},
    val3/.pic={\fill[valthree, draw=black, line width=0.3pt]
      (90:3.4pt) -- (210:3.4pt) -- (330:3.4pt) -- cycle;},
    val5/.pic={\fill[valfive, draw=black, line width=0.3pt]
      (90:3.4pt) -- (162:3.4pt) -- (234:3.4pt) -- (306:3.4pt) -- (18:3.4pt) -- cycle;}]
  \definecolor{loopspiral}{HTML}{D55E00}
  \definecolor{annulus}{HTML}{CFE0F3}
  \definecolor{valthree}{HTML}{E8890C}
  \definecolor{valfive}{HTML}{7C3AAD}
  \definecolor{measured}{HTML}{14934A}

  % The disc Omega: band, block, top.
  \fill[surface] (0,-0.6) rectangle (11,4) (3,4) rectangle (11,5)
    (0,4) rectangle (3,6) (0,6) rectangle (2,7);
  % The annulus between the curves of gamma' and gamma.
  \fill[annulus] (0,0.5) -- (0.15,0.5) -- (0.15,1.5) -- (11,1.5) -- (11,3.5)
    -- (0.15,3.5) -- (0.15,2.5) -- (0,2.5) -- cycle;
  \draw[quadgrid] (0,-0.6) grid (11,4);
  \draw[quadgrid] (3,4) grid (11,5);
  \draw[quadgrid] (0,4) grid (3,6);
  \draw[quadgrid] (0,6) grid (2,7);
  % Traces. The cut T: the left side up to c_1 and the border of the top; its
  % other side: the right side and the top of the block.
  \draw[side] (0,-0.6) -- (0,3.5) (11,-0.6) -- (11,4.5);
  \draw[cut] (0,3.5) -- (0,7) -- (2,7) -- (2,6) -- (3,6) -- (3,5) -- (11,5) -- (11,4.5);
  \draw[side] (0,4) -- (2,4) (2,6) -- (2,-0.6) (3,5) -- (2,5) (3,5) -- (3,-0.6) (2,6) -- (0,6);
  \foreach \x in {4,5,6,8} \draw[side] (\x,5) -- (\x,-0.6);
  \foreach \p in {(2,4), (2,5), (0,6), (9,5)} \node[node dot] at \p {};
  % The turns, their closing segments and where they come back.
  \foreach \o/\name in {3.5/{\gamma}, 1.5/{\gamma'}} {
    \draw[turn, loopspiral] (0,\o) -- (11.1,\o);
    \draw[closing, loopspiral] (0.15,\o-1) -- (0.15,\o);
    \draw[loopspiral, line width=0.5pt, fill=white] (0,\o-1) circle (1.8pt);
    \draw[loopspiral, line width=0.5pt, fill=white] (0,\o) circle (1.8pt);
    \node[above, loopspiral] at (9.6,\o) {$\name$};
    \node[right, loopspiral] at (11.2,\o) {$s = -1$};
  }
  \draw[loopspiral, line width=0.5pt, fill=white] (11,4.5) circle (1.8pt);
  \node[left=1pt] at (0,3.5) {$x_0$};
  \node[right=1pt] at (11,4.5) {$x_0'$};
  % The singular vertices: first visits p_i on the top of the block, and
  % the other copies on the border of the top.
  \pic at (3,5) {val3};
  \foreach \p in {(0,4), (11,5), (0,7), (8,5), (2,7), (6,5), (3,6), (4,5)} \pic at \p {val3};
  \foreach \p in {(2,6), (5,5)} \pic at \p {val5};
  \begin{scope}[font=\scriptsize, inner sep=1pt]
    \foreach \i/\x in {1/11, 2/8, 3/6, 4/5, 5/4} \node[above=2pt] at (\x,5) {$p_{\i}$};
    \node[above right=1pt] at (3,5) {$p_6$};
    \node[left=2pt] at (0,4) {$c_1$};
    \node[left=2pt] at (0,7) {$c_2$};
    \node[above=2pt] at (2,7) {$c_3$};
    \node[above left=0pt] at (2,6) {$c_4$};
    \node[right=2pt] at (3,6) {$c_5$};
  \end{scope}
  \draw[brace] (-1.5,3.5) -- (-1.5,7) node[midway, left=3pt] {$D$};
  \node[above] at (1,7) {$T$};
  % The Burgers vector of D, from x_0' to x_0, and its two components.
  \draw[burgers] (11,-1.6) -- (0,-2.6);
  \draw[measured, dashed, line width=0.5pt] (11,-1.6) -- (0,-1.6) -- (0,-2.6);
  \node[measured, above] at (5.5,-1.6) {along: $\alpha_\gamma = 11$};
  \node[measured, left] at (0,-2.1) {across: $s_\gamma = -1$};
  \node[measured, below] at (6.5,-2.3) {$B_D = \sum_i Q_i (I - R_i)\, p_i = (-11, -1)$};
\end{tikzpicture}

\caption{The shift read on the singularities
(Proposition~\ref{prop:burgers}). A cap $D$ with six singular vertices
(valences $3, 3, 3, 5, 3, 3$, total index $4$), cut along a tree $T$ and
developed; each $c_i$ is drawn at its first visit $p_i$ on the walk round $T$.
The turn $\gamma$, of length $11$, comes back one row lower, and $D$ lies on
its left: $B_D = \sum_i Q_i (I - R_i) p_i = (-11, -1) = T_\gamma(\ell)$,
minus the length along the walk and the shift $-1$ across it. The turn $\gamma'$ bounds with $\gamma$ an
annulus without singular vertex (shaded) and has the same shift (part~(c));
$11 \equiv -1 \pmod 2$ (Proposition~\ref{prop:parity}). In all figures
offsets are measured upwards.}
\label{fig:holonomy}
\end{figure}

The order matters: on Figure~\ref{fig:holonomy}, walking round $T$ the other
way gives $(1, 1)$ in place of $B_D = (-11, -1)$.

\begin{proof}
(a) Orient each edge $f$ of $T$ away from $x_0$; its two sides on the border of
$\Omega$ are glued by a motion $G_f$, from the side the walk round $T$ follows
first to the other. Let $f$ end at $v$, from which $f_1, \dots, f_r$ leave in
the order of the walk. A small arc turning clockwise round $v$ from the first
side of $f$ to the second crosses $f_1, \dots, f_r$ and turns through the
whole angle at $v$, so $G_f = G_{f_r} \circ \dots \circ G_{f_1} \circ
\rho_v$, with $\rho_v(x) = R^{4-k(v)}(x - p_v) + p_v$, $p_v$ the position at
the first visit ($\rho_v = I$ if $v$ is not singular). At $x_0$, with
$\bar\gamma$ run with $D$ on its right in place of the arc, $\hat
h(\bar\gamma) = G_{f_r} \circ \dots \circ G_{f_1}$. Unfolding the
recursion composes the $\rho_{c_i}$ in the order of first visit; the rotation
is $R^{\sum_i(4 - k(c_i))} = I$ by Lemma~\ref{lem:gauss}, expanding gives
$B_D = \hat h(\bar\gamma) = T_\gamma(\ell)^{-1}$, and with $D$ on its left
$\bar\gamma$ has the inverse holonomy, $B_D = T_\gamma(\ell)$. (b) With $x_0$ on $\bar\gamma_1$ at the start of $\kappa$, the loop
$\lambda = \bar\gamma_1\, \kappa\, \bar\gamma_2^{-1} \kappa^{-1}$ borders the
disc $A$ cut along $\kappa$. Developing along $\kappa$ reads $\hat
h(\bar\gamma_2)$ in the coordinates at $x_0$, where the walk of $\gamma_2$
runs along $e_1$; so $\hat h(\lambda)$ is the translation by $\hat
h(\bar\gamma_1) - \hat h(\bar\gamma_2)$, whose component along $e_2$ is
$\shift_{\gamma_2}(\ell) - \shift_{\gamma_1}(\ell)$, and part~(a) applies to the cut disc,
whose singular vertices have total index zero. (c) The sum is empty.
\end{proof}

In the T-meshes of openclare, whose traces start at singular vertices, an
annulus without singular vertex holds no node unless a trace winding round
enters it through a closing segment, so two turns as in (c) usually lie in
one corridor.

The simplest configuration between two parallel turns is a pair of valences
three and five, which lets the number of edges round a limb change. A pair of
opposite disclinations of rotation $\Omega$ is a dislocation with a Burgers
vector of $2\sin(\Omega/2)$ times their
separation~\cite[Eq.~(5)]{kleman2008disclinations}. We read it
on the shifts of two turns (Figures~\ref{fig:holonomy}
and~\ref{fig:dislocation}).

\begin{proposition}[A dislocation between two turns]\label{prop:dislocation}
In the setting of Proposition~\ref{prop:burgers}(b), let the annulus $A$ hold
exactly two singular vertices, $c_3$ of valence three and $c_5$ of valence
five, and let $d = p_5 - p_3 = d_1 e_1 + d_2 e_2$ be their separation,
developed inside the cut disc. Then:
\begin{enumerate}
\item[(i)] $\shift_{\gamma_1}(\ell) - \shift_{\gamma_2}(\ell) = \varepsilon_1
  d_1 + \varepsilon_2 d_2$ with signs $\varepsilon_1, \varepsilon_2 \in \{+1,
  -1\}$ that depend on the configuration but not on $\ell$, and $|B_A| =
  \sqrt2\,|d|$. The two turns can close together only if $c_5$ lies, as seen
  from $c_3$, on one particular diagonal of the field.
\item[(ii)] If $c_3$ and $c_5$ lie on one straight line of the field inside
  $A$, of length $\delta$ (in a T-mesh: they are joined by a chain of arcs
  along one trace, and $\delta$ is the sum of their lengths), and some path
  $\kappa$ as in Proposition~\ref{prop:burgers}(b) does not cross that line,
  then
  $|\shift_{\gamma_1}(\ell) - \shift_{\gamma_2}(\ell)| = \delta$. The two turns
  can both be rings only if $\delta = 0$, that is if the two singularities
  merge; as long as the bounds of separation keep $\delta \ge 1$, one of them at
  least is a helix.
\item[(iii)] For whole lengths, $\shift_{\gamma_1}(\ell) -
  \shift_{\gamma_2}(\ell) \equiv d_1 + d_2 \pmod 2$: a pair at an odd number of
  edges from one another, counted along the two directions of the field, never
  lies between two rings.
\end{enumerate}
\end{proposition}

\begin{proof}
With $R_3 = R$ and $R_5 = R^{-1}$, the two orders give $B_A = R^{-1}(I -
R)p_3 + (I - R^{-1})p_5 = (I - R^{-1})d$ or $B_A = R(I - R^{-1})p_5 + (I -
R)p_3 = (I - R)(-d)$. With $R(a, b) = (-b, a)$, $(I - R^{-1})(d_1, d_2) = (d_1 -
d_2,\ d_1 + d_2)$ and $(I - R)(-d_1, -d_2) = (-d_1 - d_2,\ d_1 - d_2)$;
reversing the orientation of the border inverts the rotations and gives the
two other sign patterns (checked symbolically). In all four both components
are $\pm d_1 \pm d_2$ and $|B_A|^2 = 2|d|^2$, which is (i). For (ii), take
$\kappa$ not crossing the segment joining $c_3$ to $c_5$ (crossing it would
turn the frame by a quarter-turn round $c_3$) and a tree containing that
segment, and develop $c_5$ from
$c_3$ along it: $d$ is $\delta$ times one of $\pm e_1$, $\pm e_2$. For (iii), $d
\in \Z^2$ and $\pm d_1 \pm d_2 \equiv d_1 + d_2$.
\end{proof}

\begin{figure}[t]
\centering
% A 3-5 dislocation between two parallel turns (Propositions 3 and 6).
% \input inside a figure; needs \usepackage{tikz} and
% \usetikzlibrary{arrows.meta,decorations.pathreplacing}.
%
% Each panel is a band of a T-mesh on a cylinder, cut open along a trace of
% the vertex c_5 of valence five: the right side is glued to the left side, the
% lower part (below c_5) at the same height, the upper part two units higher.
% Every node is at a whole position and the faint grid is the quad mesh at
% these lengths. c_5 is drawn three times (its angle 5 x 90 degrees is split
% between the three copies: 90 at the bottom left, 180 at the right, 180 at
% the left), c_3 once (270 degrees). The empty corner next to c_3 is not part of
% the surface: its two thick sides are glued to each other by a quarter-turn
% about c_3. The turn gamma_1 (below the pair) and gamma_2 (above it) bound an
% annulus that holds exactly c_3 and c_5 (Proposition 2(b)).
%
% (a) Straight chain (Proposition 3(ii)). c_3 = (2,3), c_5 = (0,3) = (10,3) =
% (2,5); the gluings are (10,y) ~ (0,y) for y < 3, (10,y) ~ (2,y+2) for y > 3,
% and the quarter-turn about c_3 taking [c_3, (0,3)] onto [c_3, (2,5)]. c_3 and
% c_5 are joined along the trace y = 3 by the chain of arcs b = [c_3, (1,3)]
% and c = [(1,3), c_5]; a trace from below (its vertex is out of the picture)
% stops at (1,3). delta = l_b + l_c = 1 + 1 = 2.
%   gamma_1 at o = 1.5: crosses x = 1, 2, 10; comes back at 1.5: s = 0,
%   length 10 (a ring).
%   gamma_2 at o = 5.5 on the side x = 2: crosses x = 10 at height 5.5, which
%   is height 7.5 on the side x = 2: s = 2, length 10 - 2 = 8 (a helix).
%   |s_2 - s_1| = 2 = delta (Proposition 3(ii)); the lengths differ by 2 too:
%   the Burgers vector of the pair is (2, 2), of norm sqrt(2) |d| with d = (2, 0)
%   from c_3 to c_5 (Proposition 3(i)); s_2 - s_1 = 2 = d_1 + d_2 mod 2 (iii).
%   The separatrix y = 5 of c_5 goes once round and stops on the trace x = 10
%   two units above c_5 (Proposition 8 gives s_2 = 2 again).
%   Both turns can be rings only if delta = 0: c_3 and c_5 merge.
%   Parity on a sphere (Proposition 5): 10 = 0 and 8 = 2 mod 2.
%
% (b) Staircase (Proposition 6). c_3 = (1,4), c_5 = (0,3) = (10,3) = (0,5);
% gluings (10,y) ~ (0,y) for y < 3, (10,y) ~ (0,y+2) for y > 3, and the
% quarter-turn about c_3 taking [c_3, (1,3)] onto [c_3, (0,4)] and
% [(1,3), (0,3)] onto [(0,4), (0,5)]. c_3 and c_5 are on no common trace; the
% staircase from c_3 to c_5 is b = [c_3, (1,3)], down, then c = [(1,3), c_5],
% left: d = p_5 - p_3 = (-1, -1), and with epsilon = (-1, -1) each step has
% epsilon . v = 1 (b: v = (0,-1); c: v = (-1,0)).
%   gamma_1 at o = 1.5: s = 0, length 10. gamma_2 at o = 5.5: crosses x = 10
%   at height 5.5 = height 7.5 on the side x = 0: s = 2, length 10.
%   s_2 - s_1 = 2 = l_b + l_c = epsilon . d (Proposition 6); here the lengths
%   are equal, the Burgers vector of the pair is (0, 2). Both are rings only
%   if l_b = l_c = 0. Parity: 10 = 0 = 2 mod 2.
% check.py builds both bands and checks these numbers by walking the strips
% of the quad mesh: valences 5 and 3 at the copies listed above, every other
% vertex of valence 4, and the turns above.
\begin{tikzpicture}[x=0.5cm, y=0.5cm, font=\footnotesize,
    side/.style={draw=black, line width=0.5pt},
    quadgrid/.style={draw=black!18, line width=0.25pt, step={(1,1)}},
    surface/.style={fill=black!3},
    node dot/.style={circle, fill=black, inner sep=0pt, minimum size=2.6pt},
    arc mark/.style={line width=1.8pt, draw=black},
    turn/.style={line width=1pt, -{Stealth[length=5pt]}},
    brace/.style={decorate, decoration={brace, amplitude=3pt}, line width=0.5pt},
    val3/.pic={\fill[valthree, draw=black, line width=0.3pt]
      (90:3.4pt) -- (210:3.4pt) -- (330:3.4pt) -- cycle;},
    val5/.pic={\fill[valfive, draw=black, line width=0.3pt]
      (90:3.4pt) -- (162:3.4pt) -- (234:3.4pt) -- (306:3.4pt) -- (18:3.4pt) -- cycle;}]
  \definecolor{loopspiral}{HTML}{D55E00}
  \definecolor{loopclosed}{HTML}{0B3C8C}
  \definecolor{valthree}{HTML}{E8890C}
  \definecolor{valfive}{HTML}{7C3AAD}

  % ---------------------------------------------------------------- (a)
  \begin{scope}
    % The surface: the lower band [0,10]x[0,3] and the upper part [2,10]x[3,8].
    \fill[surface] (0,-0.6) rectangle (10,3) (2,3) rectangle (10,8.6);
    \draw[quadgrid] (0,-0.6) grid (10,3);
    \draw[quadgrid] (2,3) grid (10,8.6);
    % Traces: the cut (a trace of c_5, both copies), the trace from c_3 round
    % to c_5, c_3 downwards, a trace from below, the winding separatrix.
    \draw[side] (0,-0.6) -- (0,3) (10,-0.6) -- (10,8.6) (2,5) -- (2,8.6);
    \draw[side] (2,3) -- (10,3) (2,-0.6) -- (2,3) (1,-0.6) -- (1,3) (2,5) -- (10,5);
    % The chain from c_3 to c_5, and its glued copy.
    \draw[arc mark] (0,3) -- (2,3) (2,3) -- (2,5);
    \node[below] at (1.5,3) {$b$};
    \node[below] at (0.5,3) {$c$};
    \foreach \p in {(1,3), (2,4), (10,5), (2,7)} \node[node dot] at \p {};
    % The turns.
    \draw[turn, loopclosed] (0,1.5) -- (10.1,1.5);
    \draw[loopclosed, line width=0.5pt, fill=white] (0,1.5) circle (1.8pt);
    \node[above, loopclosed] at (6,1.5) {$\gamma_1$};
    \draw[turn, loopspiral] (2,5.5) -- (10.1,5.5);
    \draw[loopspiral, line width=0.5pt, fill=white] (2,5.5) circle (1.8pt);
    \draw[loopspiral, dashed, line width=0.8pt, -{Stealth[length=4pt]}] (2,7.5) -- (4,7.5);
    \draw[loopspiral, line width=0.5pt, fill=white] (2,7.5) circle (1.8pt);
    \node[above, loopspiral] at (6,5.5) {$\gamma_2$};
    \draw[brace] (1.7,5.5) -- (1.7,7.5) node[midway, left=3pt] {$2$};
    % The singular vertices.
    \pic at (2,3) {val3};
    \foreach \p in {(0,3), (10,3), (2,5)} \pic at \p {val5};
    \node[above right=-1pt] at (2,3) {$c_3$};
    \node[left=1pt] at (0,3) {$c_5$};
    % Right: shifts and lengths.
    \node[right, align=left, loopspiral] at (10.3,5.5) {$s = 2$\\ 8 quads};
    \node[right, align=left, loopclosed] at (10.3,1.5) {$s = 0$\\ 10 quads};
    \node[anchor=north, align=center] at (6,-1.1)
      {$s_{\gamma_2} - s_{\gamma_1} = \ell_b + \ell_c = \delta = 2$};
    \node[anchor=north] at (6,-2.4) {\small (a) straight chain};
  \end{scope}

  % ---------------------------------------------------------------- (b)
  \begin{scope}[shift={(16.8,0)}]
    % The surface: [0,10]x[0,3], [1,10]x[3,4], [0,10]x[4,8].
    \fill[surface] (0,-0.6) rectangle (10,3) (1,3) rectangle (10,4) (0,4) rectangle (10,8.6);
    \draw[quadgrid] (0,-0.6) grid (10,3);
    \draw[quadgrid] (1,3) grid (10,4);
    \draw[quadgrid] (0,4) grid (10,8.6);
    % Traces: the cut (a trace of c_5), the loop of c_5 at y = 3, the three
    % separatrices of c_3, the winding separatrix of c_5.
    \draw[side] (0,-0.6) -- (0,3) (10,-0.6) -- (10,8.6) (0,5) -- (0,8.6);
    \draw[side] (0,3) -- (10,3) (1,4) -- (10,4) (1,4) -- (1,5) (0,5) -- (10,5);
    % The staircase from c_3 to c_5, and its glued copy.
    \draw[arc mark] (1,4) -- (1,3) -- (0,3) (1,4) -- (0,4) -- (0,5);
    \node[right] at (1,3.5) {$b$};
    \node[below] at (0.5,3) {$c$};
    \foreach \p in {(1,3), (0,4), (10,4), (0,6), (1,5), (10,5), (0,7)}
      \node[node dot] at \p {};
    % The turns.
    \draw[turn, loopclosed] (0,1.5) -- (10.1,1.5);
    \draw[loopclosed, line width=0.5pt, fill=white] (0,1.5) circle (1.8pt);
    \node[above, loopclosed] at (6,1.5) {$\gamma_1$};
    \draw[turn, loopspiral] (0,5.5) -- (10.1,5.5);
    \draw[loopspiral, line width=0.5pt, fill=white] (0,5.5) circle (1.8pt);
    \draw[loopspiral, dashed, line width=0.8pt, -{Stealth[length=4pt]}] (0,7.5) -- (2,7.5);
    \draw[loopspiral, line width=0.5pt, fill=white] (0,7.5) circle (1.8pt);
    \node[above, loopspiral] at (6,5.5) {$\gamma_2$};
    \draw[brace] (-0.3,5.5) -- (-0.3,7.5) node[midway, left=3pt] {$2$};
    % The singular vertices.
    \pic at (1,4) {val3};
    \foreach \p in {(0,3), (10,3), (0,5)} \pic at \p {val5};
    \node[above right=-1pt] at (1,4) {$c_3$};
    \node[left=1pt] at (0,3) {$c_5$};
    % Right: shifts and lengths.
    \node[right, align=left, loopspiral] at (10.3,5.5) {$s = 2$\\ 10 quads};
    \node[right, align=left, loopclosed] at (10.3,1.5) {$s = 0$\\ 10 quads};
    \node[anchor=north, align=center] at (6,-1.1)
      {$s_{\gamma_2} - s_{\gamma_1} = \ell_b + \ell_c = 2$};
    \node[anchor=north] at (6,-2.4) {\small (b) staircase};
  \end{scope}
\end{tikzpicture}

\caption{A 3--5 dislocation between a ring $\gamma_1$ and a helix $\gamma_2$
of a band, cut open along a trace of the vertex of valence five, which is
drawn three times. (a)~The pair is joined along one trace by two arcs, $\delta
= \ell_b + \ell_c = 2$: $\gamma_2$ comes back $2$ rows higher, with $8$
quads per turn, and the Burgers vector of the pair is $(2, 2)$
(Proposition~\ref{prop:dislocation}). (b)~The pair is joined by a staircase of
two arcs, $d = \varepsilon = (-1, -1)$: $\gamma_2$ comes back $2$ rows higher,
with $10$ quads (Proposition~\ref{prop:staircase}). In both, the shifts differ
by $\ell_b + \ell_c = 2$, and both turns are rings only if the two vertices
merge.}
\label{fig:dislocation}
\end{figure}

\begin{proposition}[Checkerboard rule]\label{prop:checkerboard}
In the setting of Proposition~\ref{prop:burgers}(b), with whole lengths,
colour the points $p \in \Z^2$ of the cut disc black or white by the parity of
$p_1 + p_2$ (either colour may be called white). Then $\shift_{\gamma_1}(\ell)
- \shift_{\gamma_2}(\ell)$ has the parity of the number of singular vertices
of odd valence inside $A$ that develop onto white points. In particular, two
rings bound an annulus only if that number is even.
\end{proposition}

\begin{proof}
In the sum of Proposition~\ref{prop:burgers}, a vertex of even valence has
$R_i = \pm I$ and $(I - R_i)p_i \in \{0, 2p_i\}$; for odd valence both
components of $(I - R^{\pm1})(a, b) = (a \pm b,\ b \mp a)$ have the parity of
$a + b$. Quarter-turns keep the property that both components have a given
parity, and sums add the parities, which gives the count. A quarter-turn keeps
colours; a whole translation with $t_1 + t_2$ odd swaps them, which keeps the
parity of the count because the vertices of odd valence inside $A$ are even in
number, their odd indices summing to zero with the even ones
(Lemma~\ref{lem:gauss}). Developing a vertex along another path inside the cut
disc rotates it about another singular vertex, which keeps its colour.
\end{proof}

The rule also holds on $13{,}000$ random configurations checked with SymPy.

\paragraph{Parity on a sphere.}
Each turn of a closed sphere carries its own parity condition, which needs
neither a simple curve nor the holonomy.

\begin{proposition}[Parity of a turn on a sphere]\label{prop:parity}
Let the surface be closed of genus zero and the T-mesh admit lengths that are
positive and balance every patch (its real lengths do). For every turn
$\gamma$ and all whole lengths $\ell$ balancing every patch, of any sign,
\[
  \alpha_\gamma(\ell) \equiv \shift_\gamma(\ell) \pmod 2 .
\]
\end{proposition}

\begin{proof}
First let $\ell$ be whole, balanced and positive: the patches form a quad mesh
of the sphere, whose graph is bipartite, since $H_1(S^2; \Z/2) = 0$ makes
every cycle a sum of faces of four edges. In a 2-colouring the colour of the
point $(i, j)$ of a patch $P$ is $i + j + \kappa_P \bmod 2$. For $M(x) = R^k
x + t$, $t \in \Z^2$, the coordinate sum of $M(x)$ has the parity of that of
$x$ plus $\pi(M) = t_1 + t_2$. Let $P_0, \dots, P_N = P_0$ be the patches
crossed by $\gamma$ and $D_k$ the motion developing $P_k$ into $P_0$, with
$D_N = T_\gamma(\ell)$. A node of the arc from $P_k$ to $P_{k+1}$, of
coordinates $y$, $y'$ with $D_k(y) = D_{k+1}(y')$, has the colour $(D_k y)_1 +
(D_k y)_2 + \pi(D_k) + \kappa_{P_k}$, and likewise from $P_{k+1}$; so
$\kappa_{P_k} + \pi(D_k)$ is constant, and at $k = N$, $\pi(T_\gamma(\ell))
\equiv 0$. In the frame of the first side, that of $P_0$ up to a quarter-turn
and signs, $T_\gamma(\ell) = (\alpha_\gamma(\ell),
\pm\shift_\gamma(\ell))$. For general whole balanced $\ell$: the balanced
lengths form a rational subspace with positive real points, so there are
positive whole balanced lengths $p$, and the linear form $\alpha_\gamma +
\shift_\gamma$ takes at $\ell = (\ell + Np) - Np$, $N$ large, the difference
of two even values.
\end{proof}

The proof uses only that the quad graph is bipartite; it fails in higher genus
(a $3 \times 3$ grid on a torus has rings of three quads) and with borders
(Theorem~\ref{thm:parity-any}). On a closed sphere it gives at once:

\begin{corollary}\label{cor:parity}
On a closed surface of genus zero:
\begin{enumerate}
\item[(a)] every ring of quads has an even number of quads;
\item[(b)] a helix of odd pitch winds round an odd number of quads per turn;
\item[(c)] closing a helix of odd pitch changes the length of its walk by an
  odd number of edges: a repair that keeps the length of a turn cannot close it
  if its pitch is odd, and changing the pitch by an odd amount always changes
  the length.
\end{enumerate}
\end{corollary}

\begin{proof}
(a) A ring is the strip of a walk with $\shift =
0$~\cite[Prop.~2]{gardent2026loops}, so $\alpha$ is even. (b) and (c) are the
proposition read before and after the change.
\end{proof}

Part (a) is the classical bipartiteness of quadrangulations of the
sphere~\cite[Lemma~3.1]{kaiser2015colouring}. For a turn bounding a disc, the
argument of the checkerboard rule gives both components of $B_D =
(\alpha_\gamma, \shift_\gamma)$ the parity of the number of vertices of odd
valence on white points, which is the proposition for that turn.

\begin{remark}[Borders; a second proof]\label{rem:borders}
The proposition and part (a) of the corollary need the surface to be closed:
an open cylinder of three quads round is a ring of three quads. A second proof,
from the formalization, needs neither the bipartite graph nor positive
lengths: every loop of $\Sigma^\circ$ is a product of loops round the
singular vertices (Section~\ref{sec:fox}), quarter-turns about whole points
that keep the parity of $p_1 + p_2$, so the translation of every loop has an
even coordinate sum.
\end{remark}

\paragraph{How many independent rings.}
Rings that do not cross are disjoint simple closed curves of $\Sigma^\circ$;
on a closed sphere with $n \ge 3$ singular vertices at most $n - 3$ are
pairwise non-homotopic, as for the cylinders of a Jenkins--Strebel
differential~\cite{strebel1984quadratic}, and homotopic ones share their shift,
so their closing conditions have rank at most $n - 3$ ($3g - 3 + n + b$ plus at
most $b$ classes of collars with genus $g$ and $b$ borders).

\paragraph{Certificates.}
Closing a set of turns within bounds is a question of linear feasibility.

\begin{remark}[Sign-definite turns; Farkas certificates]\label{rem:definite}
If the shift of a turn equals, modulo the balances, a form $\sum_a c_a \ell_a$
with every $c_a \ge 0$, then for all balanced $\ell \ge 0$ the shift
$\shift_\gamma(\ell)$ is never negative, and it is zero only if every arc with
$c_a > 0$ has length zero. If one of those arcs is held open by a bound of
separation, the turn is a helix under every quantization meeting the bounds;
it should never be held, and one linear program per turn finds such turns.
More generally, by Farkas' lemma (with $x = b + z$, $z \ge 0$), the system
$B\ell = 0$, $\shift_\gamma(\ell) = 0$ for $\gamma \in H$, $\ell \ge b$ has no
solution exactly when some combination $\sum_{\gamma \in H} u_\gamma
\shift_\gamma$ equals, modulo the balances, a form $\sum_a c_a \ell_a$ with $c
\ge 0$ and $\sum_a c_a b_a > 0$; and an arc $a$ has length zero in every
solution with $\ell \ge 0$ exactly when some such combination has $c_a > 0$.
The data being integral, bounds included, the multipliers of the shifts and of
the balances can be taken whole.
\end{remark}

\begin{remark}[Whole solutions exist at some size]\label{rem:exist}
The balance and closing constraints are homogeneous and the bounds of
separation are lower bounds $\ell \ge b$ with $b \ge 0$. If they have a real
solution they have a rational one, a vertex of a rational polyhedron, and
multiplying it by the common denominator $N$ gives whole lengths $N\ell \ge Nb
\ge b$ that satisfy them all. Closing in whole numbers is thus never
impossible in itself; what can fail is closing \emph{at a given size} (number
of quads) or within a given deviation from the targets.
\end{remark}

\paragraph{On the male body.}
The statements of this section hold wherever they apply (measured). The
crossing maps give the exported shift of all $1{,}299$ turns. On the
$1{,}436$ sides of the $718$ simple ones, of index $4$ and $4$ to $257$
singular vertices, the holonomy equals the sum of
Proposition~\ref{prop:burgers}(a) in both components, while the order of
last visit fails on $150$ of the $158$ discs of way~2; the checkerboard rule
holds there and on the $30{,}148$ annuli between disjoint simple turns, of
index $0$; the $294$ pairs with no singular vertex between them have equal
shifts; on the ten annuli holding one 3--5 pair,
Proposition~\ref{prop:dislocation}(i) holds, the two never on one trace.
Proposition~\ref{prop:parity} holds on all turns and on the $1{,}838$ ways
round of the strips of the eight meshes. Of the $101$ turns crossing themselves
once, $97$ split as $(3, 3, 2)$ and $4$, crossing their closing segment, as
$(2, 2, 4)$; $480$ cross themselves up to $107$ times, and only $356$ of the
$765$ held turns are simple: Section~\ref{sec:fox} handles any curve.

\section{Conflicts between turns}\label{sec:conflicts}

The held turns of the body close together if lengths are only asked to be
non-negative, not within $P_{\mathrm{eng}}$. Table~\ref{tab:iis} lists
\emph{irreducible} infeasible sets, which cannot be closed within the bounds
while every proper subset can; the Farkas multipliers were checked in exact
arithmetic, and for each turn left out, whole lengths closing the others were
found by HiGHS and checked in exact integer arithmetic (measured).

\begin{table}[t]
\centering
\small
\begin{tabular}{llllll}
\toprule
ways & size & turns & multipliers & bounded arcs & exact check \\
\midrule
0, 1, 2, 4 & 1 & 42 (way 0) & $-1$ & 1937 & way 0 \\
1, 3, 5, 7 & 2 & 54, 61 (way 1) & $1, -1$ & 1587 & --- \\
1, 3 & 2 & 45, 49 (way 1) & $1, -1$ & 1069 & --- \\
1, 3, 7 & 2 & 26, 37 (way 1) & $-1, 1$ & 619 & --- \\
5, 7 & 2 & 58, 94 (way 5) & $1, -1$ & 858 & --- \\
5 & 2 & 49, 71 & $-1, 1$ & 619, 625 & --- \\
6 & 2 & 57, 58 & $1, 1$ & 1241 & --- \\
7 & 2 & 8, 15 & $1, 1$ & 191 & --- \\
1 & 3 & 36, 37, 39 & $1, 1, -1$ & 619 & yes \\
1 & 4 & 40, 56, 59, 61 & $1, 1, -1, -1$ & 1587 & yes \\
1 & 4 & 66, 71, 75, 76 & $-\tfrac76, -\tfrac56, -\tfrac13, -\tfrac12$ & 1967, 1980, 1998 & yes \\
3 & 4 & 73, 78, 82, 83 & $-\tfrac74, -\tfrac54, -\tfrac12, -\tfrac34$ & 1980, 1998 & yes \\
5 & 7 & 59, 91, 96, 98, 99, 157, 158 & $2, 4, 1, -4, -1, 1, -1$ & 1036 & yes \\
7 & 7 & 112--116, 145, 147 & $8, \tfrac12, 1, \tfrac12, \tfrac12, -\tfrac12, \tfrac12$ & 1157 & yes \\
5 & 15 & 0, 2--9, 11, 14, 16--18, 64 & whole, up to 23 & 17 & 14 of 15 \\
\bottomrule
\end{tabular}
\caption{Irreducible infeasible sets of held turns on the male body, bounds of
separation kept (measured). Turn numbers are those of the first way listed;
``bounded arcs'': arcs held at length at least one in the support of the
non-negative form. Rows without an exact check come from the floating-point
linear program; their turns are each closable alone, so a pair is
irreducible. In the last row, the whole lengths closing the set without turn 3
were not found within the time limit.}
\label{tab:iis}
\end{table}

In four ways the turn round the left wrist cannot even be closed alone: its
shift is $-(\ell_{1921} + \ell_{1937})$ modulo the balances, and arc $1937$,
$0.08$ edges long, is held at length at least one (exact for way 0; whether
every valid quantization meets this bound we did not check). One may hope that every other conflict contains two parallel
turns whose annulus forbids it. The smallest counterexample has three turns
near the right eye, every two of which close together:
$\shift_{36} + \shift_{37} - \shift_{39} = \ell_{604} + \ell_{619} +
\ell_{620} + \ell_{637}$ modulo the balances, with $\ell_{619} \ge 1$. Every
set found is infeasible in real numbers already, so none is an obstruction of
parity. The conflicts between two turns lie in clusters of singularities at the eyes,
hands and feet, the sum or difference of the shifts being a sum of two to
eight arc lengths, one bounded; yet in none of the $13$ we examined (the pairs of
Table~\ref{tab:iis} but $49$, $71$ of way~5) are the turns separated by a
3--5 pair on one trace, as in Proposition~\ref{prop:dislocation}(ii). What
separates them is a staircase.

\begin{proposition}[Staircases]\label{prop:staircase}
Let $\gamma_1$, $\gamma_2$ be turns and $c$, $c'$ singular vertices such that
\[
  \shift_{\gamma_1}(\ell) - \shift_{\gamma_2}(\ell) = \varepsilon_1 d_1 +
  \varepsilon_2 d_2
\]
for all balanced $\ell$, with $\varepsilon_1, \varepsilon_2 \in \{+1, -1\}$
and $d = p_{c'} - p_c = d_1 e_1 + d_2 e_2$ the separation of $c'$ from $c$
developed along a path $\pi$ from $c$ to $c'$
(Propositions~\ref{prop:dislocation}(i) and~\ref{prop:swap} give such
identities). Let $\sigma$ be a path from $c$ to $c'$ along arcs $a_1, \dots,
a_r$, through no other singular vertex, and homotopic to $\pi$ with its ends
fixed in $\Sigma$ without the other singular vertices. Developed along
$\sigma$, each arc $a_i$ runs along a direction $v_i \in \{\pm e_1, \pm
e_2\}$. If $\varepsilon_1 v_{i,1} + \varepsilon_2 v_{i,2} = 1$ for every $i$,
so that $\sigma$ is a \emph{staircase}, using two directions a quarter-turn
apart and climbing across the diagonal on which the two turns could close,
then
\[
  \shift_{\gamma_1}(\ell) - \shift_{\gamma_2}(\ell) = \sum_{i=1}^r \ell_{a_i}
\]
for all balanced $\ell$. Both turns are rings only if every arc of $\sigma$
has length zero. If this sum equals, modulo the balances, a form with
non-negative coefficients containing an arc held at length at least one, one
of the two turns is a helix under every quantization meeting the bounds.
\end{proposition}

\begin{proof}
Developing along $\sigma$ adds $\ell_{a_i} v_i$ arc after arc, the nodes passed
being T-junctions. As $\sigma$ and $\pi$ are homotopic with fixed ends away
from the other cone points, $d = \sum_i \ell_{a_i} v_i$, so $\varepsilon_1 d_1
+ \varepsilon_2 d_2 = \sum_i \ell_{a_i}$, first for positive balanced lengths,
then for all by linearity.
\end{proof}

The proof applies to any path on whose steps the form of the shift difference
is one; Proposition~\ref{prop:dislocation}(ii) is the straight case.
Exchanges give the other identities.

\begin{proposition}[Exchanging two vertices]\label{prop:swap}
Let $\gamma_1$, $\gamma_2$ be turns, and $\kappa$ a path from the start of
$\bar\gamma_1$ to that of $\bar\gamma_2$ that carries the walking direction of
$\gamma_2$ onto that of $\gamma_1$, as in Proposition~\ref{prop:burgers}(b);
let $\lambda = \bar\gamma_1\, \kappa\, \bar\gamma_2^{-1} \kappa^{-1}$.
Suppose that $\lambda$ is homotopic in $\Sigma^\circ$ to $\lambda_c
\lambda_{c'}^{-1}$, where $\lambda_c$ and $\lambda_{c'}$ are small loops run
with $c$, $c'$ on their right round two singular vertices of the same valence
$k$, reached along paths that develop them at $p_c$ and $p_{c'}$. Then
\[
  \shift_{\gamma_1}(\ell) - \shift_{\gamma_2}(\ell) = \bigl\langle e_2,\
  (I - R^{4-k})(p_{c'} - p_c) \bigr\rangle .
\]
For valence three or five this is $\pm d_1 \pm d_2$ with $d = p_c - p_{c'}$,
as for a dislocation, and Proposition~\ref{prop:staircase} applies.
\end{proposition}

\begin{proof}
As in Proposition~\ref{prop:burgers}(b), $\hat h(\lambda)$ is the translation by
$\hat h(\bar\gamma_1) - \hat h(\bar\gamma_2)$, whose component along $e_2$
is $\shift_{\gamma_2}(\ell) - \shift_{\gamma_1}(\ell)$. With $j = 4 - k$, $\hat
h(\lambda_c) \circ \hat h(\lambda_{c'})^{-1}(x) = R^j\bigl(R^{-j}(x - p_{c'}) +
p_{c'} - p_c\bigr) + p_c = x + (I - R^j)(p_c - p_{c'})$; for $j = \pm1$ both
components are $\pm d_1 \pm d_2$.
\end{proof}

The $13$ conflicts fall into six configurations (Table~\ref{tab:pairs},
measured). In A the annulus between two simple turns holds a 3--5 pair, and the
staircase joining them, $\ell_{611} + \ell_{624}$, $0.990$ long, is the
difference of the real shifts; modulo the balances it equals the certificate
$\ell_{611} + \ell_{616} + \ell_{619} + \ell_{635}$, which holds arc $619$
(exact);
at the engine's whole lengths the pair is one edge apart on one grid line,
Proposition~\ref{prop:dislocation}(ii) with $\delta = 1$ (its hypothesis on
$\kappa$ not checked). In B, C and D
the turns cross and their combination winds $+1$ and $-1$ round two vertices
of valence three, an exchange in homology (the homotopy of
Proposition~\ref{prop:swap} was not checked, the staircase identity was). In E
and F the combination winds once round vertices of total index four, like a
turn round a cap, and no staircase, nor sum of two, was found equal to it.

\begin{table}[t]
\centering
\small
\setlength{\tabcolsep}{4pt}
\newcolumntype{L}[1]{>{\raggedright\arraybackslash}p{#1}}
\begin{tabular}{@{}lL{4.1cm}L{3.9cm}L{4.4cm}L{2.0cm}@{}}
\toprule
& conflict & between the turns & staircase & mechanism \\
\midrule
A & $\shift_{37} - \shift_{26}$, ways 1, 3, 7; right eye; disjoint, simple
  & annulus: 78 (val.\ 3), 76 (val.\ 5)
  & 2 arcs, $0.990$; held arc 619 by the balances
  & dislocation \\
B & $\shift_{45} - \shift_{49}$, ways 1, 3; right hand; simple, crossing
  & lenses: 102, 103 (val.\ 3)
  & 4 arcs, $1.731$; held arc 1069 by the balances
  & exchange \\
C & $\shift_{58} - \shift_{94}$, ways 5, 7; right hand
  & 109, 118 (val.\ 3)
  & 3 arcs, $2.203$; held arc 858 by the balances
  & exchange \\
D & $\shift_{57} + \shift_{58}$, way 6; left foot; figure-eights
  & 155, 156 (val.\ 3)
  & 2 arcs, $1.452$, one of them held
  & exchange \\
E & $\shift_{54} - \shift_{61}$, ways 1, 3, 5, 7; left eye
  & 11 vertices of index 3, and 207 (val.\ 3)
  & none; a closed chain of 6 combs
  & combs \\
F & $\shift_{8} + \shift_{15}$, way 7; right foot
  & 14, 23, 24, and 15 (val.\ 3)
  & none; a closed chain of 4 combs round a cap
  & combs \\
\bottomrule
\end{tabular}
\caption{The conflicts between two turns on the male body (measured). Turn
numbers are those of the first way listed, signs those of the exported
shifts; ``between the turns'': the singular vertices, numbered as nodes,
round which the combination of the two curves winds; ``staircase'': one
between them whose length equals the combination modulo the balances, and
where the held arc of the certificate lies. A:
Propositions~\ref{prop:dislocation}(i) and~\ref{prop:staircase}; B--D:
Propositions~\ref{prop:swap} and~\ref{prop:staircase}; E, F:
Section~\ref{sec:combs}.}
\label{tab:pairs}
\end{table}

\paragraph{A separatrix that winds.}
The turn that cannot be closed alone is of another kind
(Figure~\ref{fig:separatrix}).

\begin{proposition}[A separatrix that winds]\label{prop:winding}
Let $\sigma$ and $\tau$ be separatrices of a singular vertex $v$ such that
$\sigma$, followed from $v$, stops on $\tau$ at a point $q$ at distance
$\delta(\ell)$ from $v$ along $\tau$, a sum of arc lengths. Let $\lambda$ be
the loop that runs along $\sigma$ from $v$ to $q$ and back along $\tau$ to $v$,
pushed off $v$ into the sector it turns through. If $\lambda$ is freely
homotopic in $\Sigma^\circ$ to the curve $\bar\gamma$ of a turn $\gamma$, the
walking direction of $\gamma$ going to that of $\sigma$, then
$\shift_\gamma(\ell) = \pm\delta(\ell)$ for all balanced $\ell$, the sign
depending on the side of $\sigma$ on which $\tau$ lies ($+$ in
Figure~\ref{fig:separatrix}(b), $-$ at the left wrist of the male body). The turn is a
ring only if $\delta = 0$, that is if $\sigma$ comes back onto its own vertex;
if $\delta$ equals, modulo the balances, a non-negative form containing an arc
held at length at least one, $\gamma$ is a helix under every quantization
meeting the bounds.
\end{proposition}

\begin{proof}
Develop $\lambda$ from a point of $\sigma$ next to $v$, with $v$ at the origin
and $\sigma$ along $e_1$. The separatrix stops on $\tau$ at a right angle, so
the return is a segment of length $\delta$ along $\pm e_2$. Homotopic to the
curve of a turn, $\lambda$ has a translation as holonomy, which maps $v$
developed at the start onto $v$ developed at the end, at distance $\delta$
from $q$ across $\sigma$. A free homotopy conjugates the holonomy by a motion
whose rotation carries the walking direction of $\gamma$ onto that of
$\sigma$, here the identity.
\end{proof}

\begin{figure}[t]
\centering
% A separatrix that winds round a cylinder, and the turn it keeps shifted
% (Proposition 8, the mechanism of turn 54 of way 0 on the male body).
% \input inside a figure; needs \usepackage{tikz} and
% \usetikzlibrary{arrows.meta,decorations.pathreplacing}.
%
% (b) is the upper part of the band of dislocation.tex (b), the same T-mesh
% (check.py checks its numbers):
% cut open along the trace tau of the vertex v of valence five, the right side
% glued to the left side two units higher, (10,y) ~ (0,y+2) for y > 3; v is at
% (10,3) and (0,5) (and at (0,3), out of the picture with the lower band),
% the vertex c_3 of valence three at (1,4), its corner glued by a quarter-turn
% (see dislocation.tex). Traces: tau = the side x = 10 from v upwards (and its
% copy x = 0 from v upwards); sigma = the separatrix of v that leaves it to the
% right along y = 5; the separatrix of c_3 along y = 4, which stops on tau at
% n = (10,4); the separatrix of c_3 upwards, which stops on sigma at (1,5).
%   sigma goes once round, 10 units, and stops on tau at q = (10,5), at the
%   distance delta = l_e + l_f = 1 + 1 = 2 from v along tau, with
%   e = [v, n] and f = [n, q].
%   The turn gamma at o = 5.5, just above sigma, crosses tau at (10,5.5) =
%   (0,7.5): it comes back 2 higher, s_gamma = 2 = delta, length 10 = the
%   length of sigma (the along component of the holonomy of the loop lambda:
%   sigma from v to q, then tau back to v). gamma is a ring only if delta = 0,
%   that is if sigma comes back onto v.
%   By the balances of the two patches between y = 3, 4 and 5 (right of
%   x = 1), l_e equals the arc of c_3 downwards and l_f its arc upwards: if a
%   bound of separation holds either at length at least one, gamma is a helix
%   under every quantization meeting the bounds.
%   Parity (Proposition 5): 10 = 2 = 0 mod 2.
% (a) is a sketch of the same band on the cylinder, not to scale: sigma and
% gamma rise by delta = 2 rows in one round (the twist of the band above v).
\begin{tikzpicture}[x=0.46cm, y=0.46cm, font=\footnotesize,
    side/.style={draw=black, line width=0.5pt},
    quadgrid/.style={draw=black!18, line width=0.25pt, step={(1,1)}},
    surface/.style={fill=black!3},
    node dot/.style={circle, fill=black, inner sep=0pt, minimum size=2.6pt},
    separatrix/.style={draw=valfive, line width=1.1pt},
    gap/.style={draw=measured, line width=2pt},
    turn/.style={line width=1pt, -{Stealth[length=5pt]}},
    brace/.style={decorate, decoration={brace, amplitude=3pt}, line width=0.5pt},
    val3/.pic={\fill[valthree, draw=black, line width=0.3pt]
      (90:3.4pt) -- (210:3.4pt) -- (330:3.4pt) -- cycle;},
    val5/.pic={\fill[valfive, draw=black, line width=0.3pt]
      (90:3.4pt) -- (162:3.4pt) -- (234:3.4pt) -- (306:3.4pt) -- (18:3.4pt) -- cycle;}]
  \definecolor{loopspiral}{HTML}{D55E00}
  \definecolor{valthree}{HTML}{E8890C}
  \definecolor{valfive}{HTML}{7C3AAD}
  \definecolor{measured}{HTML}{14934A}

  % ---------------------------------------------------------------- (a)
  % The cylinder; a point (x, y) of the band above v sits at the angle
  % 90 - 36 x degrees (v and tau in front) and the height y + 2x/10 - 5.
  \begin{scope}[shift={(-7.5,1.4)}]
    \pgfmathsetmacro{\R}{2.4}
    \pgfmathsetmacro{\k}{0.22}
    \pgfmathsetmacro{\H}{6.6}
    % Back halves of sigma and gamma (x from 2.5 to 7.5), through the cylinder.
    \draw[valfive!70, dashed, line width=0.7pt]
      plot[domain=2.5:7.5, samples=40, variable=\x]
        ({\R*cos(90-36*\x)}, {1.6 + 0.2*\x - \k*\R*sin(90-36*\x)});
    \draw[loopspiral!70, dashed, line width=0.7pt]
      plot[domain=2.5:7.5, samples=40, variable=\x]
        ({\R*cos(90-36*\x)}, {2.1 + 0.2*\x - \k*\R*sin(90-36*\x)});
    % The cylinder, translucent.
    \fill[black!25, fill opacity=0.18] (-\R,0) -- (-\R,\H) arc[start angle=180, end angle=360,
      x radius=\R, y radius=\k*\R] -- (\R,0) arc[start angle=0, end angle=-180,
      x radius=\R, y radius=\k*\R] -- cycle;
    \draw[line width=0.5pt] (-\R,0) -- (-\R,\H) (\R,0) -- (\R,\H);
    \draw[line width=0.5pt] (0,\H) ellipse[x radius=\R, y radius=\k*\R];
    \draw[line width=0.5pt] (\R,0) arc[start angle=0, end angle=-180,
      x radius=\R, y radius=\k*\R];
    \draw[line width=0.35pt, black!40, dashed] (\R,0) arc[start angle=0, end angle=180,
      x radius=\R, y radius=\k*\R];
    % tau, in front, from v upwards.
    \draw[separatrix] (0,1.6-\k*\R) -- (0,\H-\k*\R);
    % Front halves of sigma (from v) and gamma.
    \draw[separatrix]
      plot[domain=0:2.5, samples=20, variable=\x]
        ({\R*cos(90-36*\x)}, {1.6 + 0.2*\x - \k*\R*sin(90-36*\x)});
    \draw[separatrix]
      plot[domain=7.5:10, samples=20, variable=\x]
        ({\R*cos(90-36*\x)}, {1.6 + 0.2*\x - \k*\R*sin(90-36*\x)});
    \draw[loopspiral, line width=1pt]
      plot[domain=0:2.5, samples=20, variable=\x]
        ({\R*cos(90-36*\x)}, {2.1 + 0.2*\x - \k*\R*sin(90-36*\x)});
    \draw[loopspiral, line width=1pt, -{Stealth[length=5pt]}]
      plot[domain=7.5:9.9, samples=20, variable=\x]
        ({\R*cos(90-36*\x)}, {2.1 + 0.2*\x - \k*\R*sin(90-36*\x)});
    % The other separatrices of v, started, and c_3 next to it.
    \foreach \a in {162, 234, 306}
      \draw[valfive!55, line width=0.8pt] (0,1.6-\k*\R) -- ++(\a:0.55);
    \pic at ({\R*cos(54)}, {0.8 - \k*\R*sin(54)}) {val3};
    % The gap delta on tau, from v to q.
    \draw[gap] (0,1.6-\k*\R) -- (0,3.6-\k*\R);
    \pic at (0,1.6-\k*\R) {val5};
    \node[node dot] at (0,3.6-\k*\R) {};
    \node[left=3pt] at (0,1.6-\k*\R) {$v$};
    \node[right=1pt] at (0,3.6-\k*\R) {$q$};
    \node[valfive, left] at (0,4.8-\k*\R) {$\tau$};
    \node[valfive, right] at (\R,1.95) {$\sigma$};
    \node[loopspiral, right] at (\R,2.75) {$\gamma$};
    \node[anchor=north] at (0,-2.0) {\small (a) the cylinder};
  \end{scope}

  % ---------------------------------------------------------------- (b)
  \begin{scope}
    \fill[surface] (0,1.6) rectangle (10,3) (1,3) rectangle (10,4) (0,4) rectangle (10,8.6);
    \draw[quadgrid] (0,1.6) grid (10,3);
    \draw[quadgrid] (1,3) grid (10,4);
    \draw[quadgrid] (0,4) grid (10,8.6);
    % Traces: the loop of v at y = 3, the separatrices of c_3, and the cut
    % edges of the corner of c_3; then tau and sigma.
    \draw[side] (0,1.6) -- (0,3) (10,1.6) -- (10,3);
    \draw[side] (0,3) -- (10,3) (1,4) -- (10,4) (1,4) -- (1,5);
    \draw[side] (1,4) -- (1,3) (1,4) -- (0,4) -- (0,5);
    \draw[separatrix] (10,3) -- (10,8.6) (0,5) -- (0,8.6);
    \draw[separatrix] (0,5) -- (10,5);
    \draw[gap] (10,3) -- (10,5);
    \foreach \p in {(1,3), (0,4), (10,4), (0,6), (1,5), (10,5), (0,7)}
      \node[node dot] at \p {};
    % The turn gamma, and where it comes back.
    \draw[turn, loopspiral] (0,5.5) -- (10.1,5.5);
    \draw[loopspiral, line width=0.5pt, fill=white] (0,5.5) circle (1.8pt);
    \draw[loopspiral, dashed, line width=0.8pt, -{Stealth[length=4pt]}] (0,7.5) -- (2,7.5);
    \draw[loopspiral, line width=0.5pt, fill=white] (0,7.5) circle (1.8pt);
    \node[above, loopspiral] at (5,5.5) {$\gamma$};
    \draw[brace] (-0.3,5.5) -- (-0.3,7.5) node[midway, left=3pt] {$s_\gamma = 2$};
    % Singular vertices and names.
    \pic at (1,4) {val3};
    \foreach \p in {(0,3), (10,3), (0,5)} \pic at \p {val5};
    \node[above right=-1pt] at (1,4) {$c_3$};
    \node[left=1pt] at (0,4.9) {$v$};
    \node[right=1pt] at (10,3) {$v$};
    \node[right=1pt] at (10,4) {$n$};
    \node[right=1pt] at (10,5) {$q$};
    \node[valfive, below] at (5,5) {$\sigma$};
    \node[valfive, right] at (10,8) {$\tau$};
    \node[measured, left] at (10,3.5) {$e$};
    \node[measured, left] at (10,4.5) {$f$};
    \draw[brace, decoration={mirror}] (11.3,3) -- (11.3,5)
      node[midway, right=3pt] {$\delta = \ell_e + \ell_f = 2$};
  \end{scope}
  \node[anchor=north] at (5,0.9) {$s_\gamma = \delta = 2$: a ring only if $\sigma$ comes back onto $v$};
  \node[anchor=north] at (5,-0.6) {\small (b) cut open along $\tau$};
\end{tikzpicture}

\caption{A separatrix that winds (Proposition~\ref{prop:winding}). (a)~On
a cylinder such as a wrist, the separatrix $\sigma$ of the vertex $v$ of
valence five goes once round and stops on another separatrix $\tau$ of $v$.
(b)~Cut open along $\tau$: $\sigma$, of length $10$, stops at $q$, at
$\delta = \ell_e + \ell_f = 2$ from $v$, and the turn $\gamma$ beside
$\sigma$ comes back $2 = \delta$ rows higher; it is a ring only if $\sigma$
comes back onto $v$. At the left wrist of the male body $\delta = \ell_{1920}
+ \ell_{1921}$, kept positive by a bound of separation.}
\label{fig:separatrix}
\end{figure}

At the left wrist the turn ($42$ of way 0, length $6.99$, real shift
$-0.305$) bounds a cap of four vertices of valence three, whose Burgers sum
gives its shift as a form without sign. The sign comes from outside: the
separatrix $964$ of the vertex $246$, of valence five, runs $6.994$ edges once
round the wrist and stops on another of its separatrices at $\delta =
\ell_{1920} + \ell_{1921} = 0.305$ from $246$; the shift is $-\delta$, and
$\ell_{1920}$ equals the length of the held arc $1937$, both modulo the
balances: whence the certificate and $\shift \le -1$ at every admissible whole
length. We checked the identity, not the homotopy (measured).

\paragraph{Larger conflicts.}
The three turns at the right eye wind round one 3--5 pair, and their
combination is the length of a staircase of three arcs through the held arc
$619$; the four at the eyes wind round a 3--5 pair and two vertices of valence
three, a dislocation and an exchange in homology. The fractional multipliers of Table~\ref{tab:iis} come
from asking the bounded arcs to weigh one in total; at the left hand, where
the turns cross each other $10$ to $46$ times, the primitive ones give
\[
  -(7\shift_{66} + 5\shift_{71} + 2\shift_{75} + 3\shift_{76}) = 2F =
  10\,\ell_{1805} + 2\ell_{1817} + 4\ell_{1853} + 2(\ell_{1870} + \ell_{1879}
  + \ell_{1969} + \ell_{1998} + \ell_{2029})
\]
modulo the balances. The coefficient $5$ is a dislocation: a separatrix of the
vertex $231$ (valence three) stops on one of $228$ (valence five), and arc
$1805$, $0.063$ edges long, joins that junction to $228$, so $\ell_{1805}$ is
the offset of the pair across the line of the field that nearly joins them,
Proposition~\ref{prop:dislocation}(ii) with $\delta = \ell_{1805}$ (hypothesis
on $\kappa$ not checked); their
twisted windings (Section~\ref{sec:fox}) contribute exactly $10\,\ell_{1805}$
(measured).

\paragraph{Local dislocations.}
The curves of an irreducible set need not bound one planar region round a
cluster of singularities: configurations B to F and the sets of four and seven
turns do not. What the conflicts share is local.

\begin{question}[Are conflicts sums of local dislocations?]\label{conj:local}
On a closed surface of genus zero, let $H$ be an irreducible set of turns that
cannot be closed together within the bounds of separation. Does some
combination $\sum_{\gamma \in H} u_\gamma \shift_\gamma$ with whole
multipliers that certifies the conflict (Remark~\ref{rem:definite}) equal,
modulo the balances, a sum with non-negative whole weights of
\begin{enumerate}
\item[(i)] lengths of staircases (Proposition~\ref{prop:staircase}) between a
  vertex of valence three and one of valence five, or between two vertices of
  equal valence exchanged as in Proposition~\ref{prop:swap}, and
\item[(ii)] shifts of caps: sign-definite forms carried by curves that wind
  once round singular vertices of total index four, as the wrist turn and its
  separatrix (Proposition~\ref{prop:winding}),
\end{enumerate}
one term at least containing an arc held open?
\end{question}

It holds with one term, up to the unchecked homotopies, for the wrist,
configurations A to D and the right eye; Section~\ref{sec:combs} settles it on
the body: no as stated, yes and empty when read loosely.

\section{Any combination of turns}\label{sec:fox}

Winding numbers are the homology class of a weighted cycle $\sum_\gamma
u_\gamma \bar\gamma$, which is commutative, while the holonomy is not; the
free differential calculus of Fox~\cite{fox1953free} keeps the order. Let
$\Sigma$ be closed of genus zero, with singular vertices $c_1, \dots, c_n$
(Proposition~\ref{prop:burgers-any} adds borders and handles). Fix $x_0$ inside
a patch, with its frame, and disjoint paths $\pi_i$ from $x_0$ to $c_i$; let
$\lambda_i$ follow $\pi_i$, go round $c_i$ with $c_i$ on its right, and come
back. The fundamental group of $\Sigma^\circ$ is generated by the $\lambda_i$
with a single relation, their product in the order in which the paths leave
$x_0$. Let $p_i$ be the developed position of $c_i$ at the end of $\pi_i$ and
$R_i = R^{4 - k(c_i)}$, so that $\hat h(\lambda_i) = \rho_i$, $\rho_i(x) =
R_i(x - p_i) + p_i$. Write $\hat h(\gamma)(x) = \rho(\gamma)x + \tau(\gamma)$;
then $\tau(\alpha\beta) = \tau(\alpha) + \rho(\alpha)\,\tau(\beta)$. Let $F$
be the free group on $x_1, \dots, x_n$, mapped onto the fundamental group by
$x_i \mapsto \lambda_i$, with the Fox derivatives $\partial x_j/\partial x_i =
\delta_{ij}$, $\partial(uv)/\partial x_i = \partial u/\partial x_i +
u\,\partial v/\partial x_i$. Extended linearly to the group ring, $\rho$ takes
the values $a_0 I + a_1 R + a_2 R^2 + a_3 R^3$, and the augmentation $\epsilon$
maps such an element to $a_0 + a_1 + a_2 + a_3$. The first identity below is
the defining property of Fox derivatives for a crossed
homomorphism~\cite{fox1953free}, and the $p_i$ are, up to a common motion,
twisted-cohomology coordinates of the flat
metric~\cite{thurston1998shapes,veech1993flat}; what is new is the reading on
combinations of turns.

\begin{proposition}[Fox--Burgers formula]\label{prop:fox}
For every word $w$ of $F$,
\[
  \tau(w) = \sum_{i=1}^n \rho\Bigl(\frac{\partial w}{\partial x_i}\Bigr)
  (I - R_i)\,p_i .
\]
Let $\gamma_1, \dots, \gamma_m$ be turns, $\kappa_r$ a path from $x_0$ to the
start of $\bar\gamma_r$, $f_r$ the direction $e_2$ of $\gamma_r$ carried back
to $x_0$ along $\kappa_r$ (one of $\pm e_1$, $\pm e_2$), and $w_r$ a word for
$\kappa_r \bar\gamma_r \kappa_r^{-1}$. For multipliers $u_r$ define the
\emph{twisted winding} of the weighted cycle round $c_i$,
\[
  W_i = -\sum_{r=1}^m u_r\, \rho\Bigl(\frac{\partial w_r}{\partial x_i}\Bigr)^{\!\top} f_r
  \in \Z^2 \text{ (for whole } u_r) .
\]
Then, for all balanced lengths,
\[
  \sum_{r=1}^m u_r\, \shift_{\gamma_r}(\ell) = \sum_{i=1}^n \bigl\langle W_i,\
  (I - R_i)\,p_i \bigr\rangle ,
  \qquad
  \sum_{i=1}^n (I - R_i)^\top W_i = 0 ,
\]
so the sum does not change when all the $p_i$ are moved by one translation: it
is a combination of separations $p_i - p_j$. The ordinary winding number of
$\sum_r u_r \bar\gamma_r$ round $c_i$ is, up to a constant common to all $i$,
$w_i = \sum_r u_r\, \epsilon(\partial w_r/\partial x_i)$. The words $w_r$ are
defined up to the relation, and the $W_i$ accordingly: inserting a conjugate
$g\varrho g^{-1}$ of the relation $\varrho$ in $w_r$, after a prefix $a$ of
$w_r$, subtracts $u_r \rho(a g\, \partial \varrho/\partial x_i)^\top f_r$ from
every $W_i$ (with $a$ empty, $u_r \rho(g\, \partial \varrho/\partial x_i)^\top f_r$),
and leaves the sum unchanged.
\end{proposition}

\begin{proof}
Both $w \mapsto \tau(w)$ and the right-hand side satisfy $\phi(uv) = \phi(u) +
\rho(u)\phi(v)$, and they agree on the letters and the empty word, so they
agree everywhere by induction on the length (for $x_i^{-1}$, use $0 =
\phi(x_i x_i^{-1})$). For a turn, $\hat h(\kappa_r \bar\gamma_r
\kappa_r^{-1})$ is the translation by $M_r T_r$, with $T_r$ the translation of
$\bar\gamma_r$ in its own frame and $M_r$ the rotation of $\hat h(\kappa_r)$;
so $\shift_{\gamma_r} = -\langle M_r e_2, M_r T_r \rangle = -\langle f_r,
\tau(w_r) \rangle$, and $\langle f, \rho(D) b\rangle = \langle \rho(D)^\top f,
b \rangle$ gives the formula. The fundamental formula $\sum_i (\partial
w/\partial x_i)(x_i - 1) = w - 1$ maps under $\rho$ to $\sum_i \rho(\partial
w_r/\partial x_i)(R_i - I) = \rho(w_r) - I = 0$, the rotation of a turn being
trivial; transposed and summed with the $u_r f_r$, it is the second identity.
Finally $\epsilon(\partial w/\partial x_i)$ is the exponent sum of $x_i$ in
$w$, the homology class of $w$ on the basis of the $\lambda_i$, defined up to
the relation, in which each $\lambda_i$ appears once.
\end{proof}

For a simple turn $\rho(\partial w/\partial x_i) = Q_i$ and the formula is
Proposition~\ref{prop:burgers}(a). In general a cycle may go round $c_i$
several times after different rotations; modulo $2$ these disappear.

\begin{corollary}[Checkerboard rule for any combination]\label{cor:checker}
On a closed surface of genus zero, let $\ell$ be whole and balanced and the
multipliers $u_r$ whole. Colour each singular vertex $c_i$ by the parity
$\chi_i$ of $p_{i,1} + p_{i,2}$; it does not depend on the paths $\pi_i$, and
for positive lengths it is the colour of $c_i$ in the bipartition of the quad
mesh. Then
\[
  \sum_r u_r\, \shift_{\gamma_r}(\ell) \equiv \sum_r u_r\, \alpha_{\gamma_r}(\ell)
  \equiv \sum_{i:\ k(c_i) \text{ odd}} w_i\, \chi_i \pmod 2 ,
\]
with $w_i$ the winding numbers of $\sum_r u_r \bar\gamma_r$, whatever the
common constant and whatever the orientation of each $\bar\gamma_r$. In
particular, a combination whose cycle winds an even number of times round every
vertex of odd valence takes even values at every whole length.
\end{corollary}

\begin{proof}
Since $(R - I)^2 = -2R$, $R^k(I - R^{\pm1}) \equiv I - R^{\pm1}$ modulo $2$, so
$\rho(D)(I - R_i) \equiv \epsilon(D)(I - R_i)$ for $k(c_i)$ odd, and $\equiv 0$
for $k(c_i)$ even, where $I - R_i \in \{0, 2I\}$; and $\langle f, (I -
R^{\pm1})p\rangle \equiv p_1 + p_2$ for $f \in \{\pm e_1, \pm e_2\}$. This
gives the sum of Proposition~\ref{prop:fox} modulo $2$, and the direction
$e_1$ of each turn that of $\sum_r u_r \alpha_{\gamma_r}$. Rotations about
whole points and quarter-turns keep $p_1 + p_2 \bmod 2$, so $\chi_i$ does not
depend on the paths. A whole translation of the origin flips all colours and
adds $\sum_{k(c_i)\text{ odd}} w_i$, even because the rotation of the cycle is
trivial ($\sum_i w_i (4 - k(c_i)) \equiv 0 \pmod 4$); adding $q$ to all $w_i$
adds $q \sum_{k(c_i)\text{ odd}} \chi_i$, the value modulo $2$ of the
relation, a loop with zero translation, hence even. Reversing a turn changes
signs only. For positive lengths the developed parity colours the quad graph
properly, as in Proposition~\ref{prop:parity}.
\end{proof}

The corollary contains Propositions~\ref{prop:checkerboard},
\ref{prop:dislocation}(iii) and~\ref{prop:parity}. As for crystal
defects~\cite{kleman2008disclinations} and in~\cite[\S7]{izmestiev2013torus},
winding numbers do not determine the holonomy:

\begin{proposition}[Winding numbers do not suffice]\label{prop:commutator}
Let $c_a$, $c_b$ be singular vertices of odd valence developed at distinct
points $p_a \neq p_b$. The loop $\lambda_a \lambda_b \lambda_a^{-1}
\lambda_b^{-1}$ winds zero times round every singular vertex, and its holonomy
is the translation by $(I - R_a)(I - R_b)(p_a - p_b)$, of length $2|p_a - p_b|
\neq 0$. No function of the winding numbers gives the holonomy of the loops
with trivial rotation.
\end{proposition}

\begin{proof}
For $w = x_a x_b x_a^{-1} x_b^{-1}$, $\rho(\partial w/\partial x_a) = \rho(1 -
x_a x_b x_a^{-1}) = I - R_b$ and $\rho(\partial w/\partial x_b) = \rho(x_a -
w) = R_a - I$, the rotations commuting; Proposition~\ref{prop:fox} gives the
translation, and for odd valences $I - R_a$, $I - R_b$ are $\sqrt2$ times
rotations. The exponent sums of $w$ are zero.
\end{proof}

\paragraph{On the male body.}
Cut into one chart along a spanning tree of the nodes, the T-mesh gives every
turn a word, and Proposition~\ref{prop:fox} its translation, for all
$1{,}299$ turns. Corollary~\ref{cor:checker} holds on all of them and on the
first $14$ sets of Table~\ref{tab:iis}, fractional multipliers scaled to whole
ones (the set of $15$ turns was not tried); at the left hand the
cycle winds an even number of times round every vertex of odd valence, which
is why $-2F$ is even. Four turns satisfy $w_{18} = w_{22} + w_{101} + w_{65}$
round every singular vertex while their translations do not, in any
quarter-turn frames: a ``winding Burgers formula'' with numbers for weights is
false (measured).

\subsection{Combs}\label{sec:combs}

The twisted windings depend on the chart; combs do not. In the T-meshes of
openclare, traced from the separatrices of the singular vertices only, every
trace of a closed surface starts at a singular vertex (feature lines or seeded
traces would break this; we assume it). For an arc $a$, along $e_a$ from the node $n_0$
to the node $n_1$, the \emph{foot} of $n_j$ is the node itself if it is
singular, and otherwise the singular vertex where the trace through $n_j$
perpendicular to $a$ starts. We assume this vertex unique: a single trace
perpendicular to $a$ meets $n_j$ (not two, from opposite sides), and it is neither a separatrix
connection nor a loop; otherwise a foot must be chosen, and different choices
give different dislocations. The \emph{comb} $\sigma_a$ of $a$ runs from the
foot $c_0$ of $n_0$ along its trace to $n_0$, along $a$ to $n_1$, and along the
other trace to the foot $c_1$ of $n_1$, through T-junctions only
(Figure~\ref{fig:comb}). The \emph{graph of feet} has the singular vertices of
odd valence as vertices and the arcs whose two feet are distinct vertices of
odd valence as edges.

\begin{figure}[t]
\centering
% The comb of an arc (Section "Combs" and Lemma "Combs").
% \input inside a figure; needs \usepackage{tikz} and
% \usetikzlibrary{arrows.meta,decorations.pathreplacing}.
%
% The T-mesh is that of dislocation.tex (b) and separatrix.tex (b): a band cut
% open along the trace tau of the vertex v of valence five, the right side
% glued to the left side two units higher above y = 3 and at the same height
% below; v is drawn at (0,3), (10,3) and (0,5), the vertex c_3 of valence three
% at (1,4), and the empty corner next to c_3 is glued by a quarter-turn about
% c_3 (see dislocation.tex). Every node is at a whole position; check.py
% checks the numbers below.
%
% The arc a is the piece of the separatrix sigma of v from the T-junction
% n_0 = m = (1,5), where the separatrix of c_3 upwards stops on sigma, to the
% T-junction n_1 = q = (10,5), where sigma stops on tau; e_a = (1,0), l_a = 9.
% Foot of n_0: the trace through m perpendicular to a is the separatrix of c_3
% upwards, which starts at c_3: c_0 = c_3. Foot of n_1: the trace through q
% perpendicular to a is tau, which starts at v: c_1 = v. The comb sigma_a runs
% from c_3 up to m (1 unit), along a to q (9 units), and down tau through the
% T-junction n = (10,4) to v (2 units): it passes through T-junctions only.
% Developed along it, p_{c_0} = (1,4) and p_{c_1} = (10,3), and
%   < e_a, p_{c_1} - p_{c_0} > = < (1,0), (9,-1) > = 9 = l_a
% (Lemma "Combs"); the teeth only add the across component -1.
% The feet are distinct, of valences 3 and 5: the comb reads the length of a
% as that of a dislocation (Proposition 3); with two feet of valence 3 it would
% be an exchange (Proposition 7), with a single foot a winding separatrix
% (Proposition 8).
\begin{tikzpicture}[x=0.5cm, y=0.5cm, font=\footnotesize,
    side/.style={draw=black, line width=0.5pt},
    quadgrid/.style={draw=black!18, line width=0.25pt, step={(1,1)}},
    surface/.style={fill=black!3},
    node dot/.style={circle, fill=black, inner sep=0pt, minimum size=2.6pt},
    comb/.style={draw=measured, line width=1.8pt, line join=round},
    arc mark/.style={line width=2.2pt, draw=black},
    dimension/.style={line width=0.5pt, {Bar[width=5pt]}-{Bar[width=5pt]}},
    val3/.pic={\fill[valthree, draw=black, line width=0.3pt]
      (90:3.4pt) -- (210:3.4pt) -- (330:3.4pt) -- cycle;},
    val5/.pic={\fill[valfive, draw=black, line width=0.3pt]
      (90:3.4pt) -- (162:3.4pt) -- (234:3.4pt) -- (306:3.4pt) -- (18:3.4pt) -- cycle;}]
  \definecolor{valthree}{HTML}{E8890C}
  \definecolor{valfive}{HTML}{7C3AAD}
  \definecolor{measured}{HTML}{14934A}

  \begin{scope}
    \fill[surface] (0,1.6) rectangle (10,3) (1,3) rectangle (10,4) (0,4) rectangle (10,8);
    \draw[quadgrid] (0,1.6) grid (10,3);
    \draw[quadgrid] (1,3) grid (10,4);
    \draw[quadgrid] (0,4) grid (10,8);
    % All traces of the T-mesh in the picture.
    \draw[side] (0,1.6) -- (0,3) (10,1.6) -- (10,8) (0,5) -- (0,8);
    \draw[side] (0,3) -- (10,3) (1,4) -- (10,4) (1,4) -- (1,5) (0,5) -- (10,5);
    \draw[side] (1,4) -- (1,3) (1,4) -- (0,4) -- (0,5);
  \end{scope}
  % The comb of a: from its foot c_3 up to m, along a, down tau to its foot v.
  \draw[comb] (1,4) -- (1,5) -- (10,5) -- (10,3);
  \draw[comb, line width=3pt] (1,5) -- (10,5);
  \foreach \p in {(1,3), (0,4), (10,4), (0,6), (1,5), (10,5), (0,7)}
    \node[node dot] at \p {};
  \pic at (1,4) {val3};
  \foreach \p in {(0,3), (10,3), (0,5)} \pic at \p {val5};
  % Names.
  \node[above] at (5.5,5) {$a$};
  \draw[-{Stealth[length=4pt]}, line width=0.5pt] (6.6,5.75) -- (8,5.75)
    node[right, inner sep=1pt] {$e_a$};
  \node[above left=-1pt] at (1,5) {$n_0$};
  \node[above right=-1pt] at (10,5) {$n_1$};
  \node[below right=-1pt] at (1,4) {$c_0 = c_3$};
  \node[right=2pt] at (10,3) {$c_1 = v$};
  \node[left=1pt] at (0,4.95) {$v$};
  \node[left=1pt] at (0,3) {$v$};
  \node[measured, right=1pt] at (10,4.5) {$\sigma_a$};
  \node[right=1pt] at (10,7.4) {$\tau$};
  % The developed feet and the e_a component of their separation.
  \draw[dimension] (1,2.2) -- (10,2.2) node[midway, fill=black!3, inner sep=1.5pt]
    {$\langle e_a,\, p_{c_1} - p_{c_0}\rangle = \ell_a = 9$};
  \draw[black!50, dotted, line width=0.6pt] (1,2.2) -- (1,4) (10,2.2) -- (10,3);
\end{tikzpicture}

\caption{The comb $\sigma_a$ of an arc $a$ (Lemma~\ref{lem:comb}), $\ell_a =
9$: from the foot $c_0$ of one end of $a$, of valence three, up its trace,
along $a$, and down the other perpendicular trace to the foot $c_1$, of
valence five. Every step but $a$ is orthogonal to $a$, so $\langle e_a,
p_{c_1} - p_{c_0} \rangle = 9 = \ell_a$. This T-mesh is that of
Figures~\ref{fig:dislocation}(b) and~\ref{fig:separatrix}(b): the staircase,
the winding separatrix and the comb are three readings of the same shift
$2$.}
\label{fig:comb}
\end{figure}

\begin{lemma}[Combs]\label{lem:comb}
Let $p_{c_0}$, $p_{c_1}$ be the feet of an arc $a$ developed along its comb.
Then, for all balanced lengths,
\[
  \langle e_a,\ p_{c_1} - p_{c_0} \rangle = \ell_a .
\]
If $c_0 \neq c_1$ have odd valence, $2\ell_a$ is, up to sign, the form of the
twisted class $\lambda_{c_0}\lambda_{c_1}^{\pm1} \otimes v$ of the loops round
the two feet reached along the comb, for a whole $v$ with $(I - R^{\pm1})^\top v
= 2e_a$: the pair is a dislocation if the valences differ
(Proposition~\ref{prop:dislocation}) and an exchange if they are equal
(Proposition~\ref{prop:swap}). If $c_0 = c_1$, $\ell_a$ is the shift of a
winding separatrix (Proposition~\ref{prop:winding}).
\end{lemma}

\begin{proof}
Every step of the comb but $a$ runs along a trace perpendicular to $a$, through
no cone point, and adds a vector orthogonal to $e_a$; the step along $a$ adds
$\ell_a e_a$. The loop $\lambda_{c_0} \lambda_{c_1}^{\pm1}$, with the sign that
makes its rotation trivial, has the translation $\pm(I - R^{\pm1})(p_{c_0} -
p_{c_1})$ as in Propositions~\ref{prop:dislocation} and~\ref{prop:swap};
pairing it with $v$ gives $\pm 2\ell_a$, and $v$ is whole because $I -
R^{\pm1}$ maps $\Z^2$ onto the vectors whose coordinates have equal parity,
among them $2e_a$. If $c_0 = c_1$, the second trace is a separatrix of the same
vertex stopping on the first at the distance $\ell_a$.
\end{proof}

A \emph{closed walk of combs} is a sequence of arcs $a_1, \dots, a_k$ whose
feet are $c_{i-1}$ and $c_i$, of odd valence, with $c_k = c_0$; its
\emph{reading} is $\sum_i \ell_{a_i}$. A walk of length one is a winding
separatrix, one of length two between valences three and five the staircase of
Proposition~\ref{prop:staircase}, and one round four vertices of valence three
a cap.

\begin{proposition}[Certificates are sums of combs]\label{thm:combs}
Let $c = \sum_a c_a\ell_a$, $c_a \in \Z_{\ge0}$, be a certificate
(Remark~\ref{rem:definite}) whose arcs have feet of odd valence and whose odd
part, the arcs with $c_a$ odd, is a cycle modulo~$2$ of the graph of feet.
Then $c$ is a sum with whole non-negative weights of comb readings $2\ell_a$,
each the form of a dislocation or an exchange (Lemma~\ref{lem:comb}), and of
readings of closed walks of combs. With real weights, every non-negative form
supported on arcs whose two feet are of odd valence, distinct or equal, is a
non-negative combination of comb readings $2\ell_a$ and of winding separatrices
$\ell_a$, without condition of parity.
\end{proposition}

\begin{proof}
The odd part splits into closed walks along its edges (Veblen), and what
remains is $\sum_a \lfloor c_a/2 \rfloor\,2\ell_a$. With real weights, $c =
\sum_a \tfrac{c_a}{2}\,(2\ell_a)$ on the arcs with two distinct feet, and
$c_a\ell_a$ on those whose feet are one vertex.
\end{proof}

The proof is elementary: the content is in the hypothesis. That the odd part of
a certificate is a cycle of the graph of feet is observed on all $25$
certificates of the body (exact) but not proved; Corollary~\ref{cor:checker}
does not give it, since it controls the parity of $p_{i,1} + p_{i,2}$ while a
comb reads one coordinate (in Figure~\ref{fig:comb} both feet have the same
colour and $\ell_a = 9$).

\begin{remark}[Closed walks as twisted cycles]
We expect the reading of a closed walk of combs, $\sum_i \langle e_{a_i} -
e_{a_{i+1}},\ p_{c_i}\rangle$ by a summation by parts, to be the Burgers form
of a twisted cycle near its combs; but the two combs through $c_i$ develop it
in charts that differ by a rotation, and a proof must develop the whole walk
in one chart. We use closed walks only as forms.
\end{remark}

\begin{proposition}[The Burgers pairing]\label{prop:pairing}
Let the surface be closed of genus zero, with no singular vertex of valence
divisible by four, and let every trace of the T-mesh start at a singular
vertex. Then the map that sends a class of the twisted homology
$H_1(\Sigma^\circ; \R^2_\rho)$ to its form $\sum_i \langle W_i, (I - R_i)p_i
\rangle$ on the balanced lengths is onto the dual of the balanced lengths. Its
kernel contains the class across each separatrix connection. If some singular
vertices have valence divisible by four, no form depends on their developed
positions, and the map need not be onto.
\end{proposition}

\begin{proof}
If every such form vanishes at balanced $\ell$, every loop of $\Sigma^\circ$
with trivial rotation has a holonomy without translation
(Proposition~\ref{prop:fox}). The commutator of $\lambda_a$ and $\lambda_b$
has trivial rotation and translation $(I - R_a)(I - R_b)(p_a - p_b)$
(Proposition~\ref{prop:commutator}); no $R_i$ is the identity, so $I - R_i$ is
invertible ($\sqrt2$ times a rotation, or $2I$), and every singular vertex
develops at one point $q$. As the $\rho_i$ fix $q$ and generate the holonomy on a closed
sphere, every singular vertex develops at $q$ along any path; every arc has
feet, so every comb offset of Lemma~\ref{lem:comb} is zero, and $\ell = 0$.
The image has no common zero but $0$. Across a connection the
separation of its two vertices is zero across the trace, and the twisted class
pairing the across direction with it has the zero form. A vertex of valence
divisible by four has $(I - R_i)p_i = 0$: its position is a relative period,
which the lengths can move without changing any holonomy.
\end{proof}

\paragraph{On the male body.}
Lemma~\ref{lem:comb} holds on all $2{,}069$ arcs, $2{,}046$ of which have two
distinct feet of odd valence: nearly every arc is a local dislocation or
exchange. On this T-mesh (way~1) each node meets a single perpendicular trace,
and one arc, $1676$, ends on the separatrix connection, which joins the
vertices $204$ and $210$; we took as its foot the vertex the connection was
traced from, $204$, and no certificate of way~1 contains the arc. The
twisted homology has dimension $518 = 2(n - 2)$ and the pairing rank $517$,
its kernel the class across the one separatrix connection; we conjecture that
without vertices of valence divisible by four the kernel is spanned by the
classes across the connections. All $25$ certificates of the conflicts
examined have a representative whose odd part is a cycle of the graph of feet
(exact). The wrist is a comb with both feet at $246$, the winding separatrix;
A, C and D are closed walks of two combs (A the staircase $\ell_{611} +
\ell_{624}$), B of four, the right eye of three, E of six, of total index
$-2$, and F of four round a cap; the left hand is five times the dislocation
$\ell_{1805}$ plus eight comb readings on a tree of seven arcs, and the set of
seven of way~7 a tree of seven combs (measured).

\paragraph{Question~\ref{conj:local}, read loosely, is true and empty.}
Allowing every comb reading and every closed walk of combs,
Proposition~\ref{thm:combs} answers it for every certificate whose odd part is
a cycle of the graph of feet, all those of the body; but this follows from
Farkas' lemma and a parity, and says nothing about which configurations
conflict. \emph{Read strictly, it is false}: restricted to staircases of
weight one, monotone walks round a cap and winding separatrices, exact
rational balanced lengths make every such form non-negative and the
certificates of E and F negative, $-2.9967$ and $-1.9973$ (checked in integer
arithmetic). These forms are infinitely many, but each is the length of a walk
on a graph whose steps carry the arc lengths, so a shortest-path search finds
a negative one if any. In floating point the same test also separates the four
turns at the eyes and puts the left hand inside the cone (measured). Whether
every certificate splits into combs on disjoint sets of vertices, each of its
sign, is open: the trees of the left hand and of way~7 share vertices.

\section{Borders and handles}\label{sec:beyond}

Models with borders (the eye sockets of Suzanne, the mouth of Jerry) and
handles (Bob, a torus; part and cage) need the forms below, measured on the
T-meshes of these bench models with an engine version not recorded. A border
node $b$ is a corner of angle $a(b)\pi/2$, of
\emph{index} $2 - a(b)$; for a border component $\beta_j$, $L_j(\ell)$ is the
sum of the lengths of its arcs.

\begin{lemma}[Gauss--Bonnet with borders]\label{lem:gauss-border}
Let $D$ be a region bounded by curves of turns and by border components. Then
\[
  \sum_{c \in D} (4 - k(c)) + \sum_{b \in D \cap \partial\Sigma} (2 - a(b))
  = 4\chi(D) - \Theta,
\]
where $\Theta$ is the turning of the curves of the turns at their corners, in
quarter-turns: zero for the side of a simple curve.
\end{lemma}
\begin{proof}
Gauss--Bonnet for a flat surface with cone points and a piecewise straight
boundary, border nodes being corners of turning $(2 - a)\pi/2$.
\end{proof}

A side of a simple turn holding $b$ border components without corner and $h$
handles has index $4(1 - b - 2h)$, and may hold no singular vertex at all; the
lemma holds on the $1{,}076$ sides of turns of Suzanne and Jerry (measured).
For positive lengths the homomorphism of the next theorem is the cycle parity
of the quad graph, a $\Z/2$-homology invariant of any
quadrangulation~\cite[Lemma~3.1]{kaiser2015colouring}
extended here to turns at all balanced
lengths.

\begin{theorem}[Parity of a turn on any surface]\label{thm:parity-any}
Let $\ell$ be whole and balanced. There is a homomorphism $\Omega_\ell\colon
H_1(\Sigma; \Z/2) \to \Z/2$ such that for every turn
\[
  \alpha_\gamma(\ell) + \shift_\gamma(\ell) \equiv \Omega_\ell([\bar\gamma]) \pmod 2 .
\]
Its value on a border component is the parity of $L_j(\ell)$, the sum of all
$L_j(\ell)$ is even, and for positive lengths $\Omega_\ell(z)$ is the parity of
the length of any cycle of the quad graph in the class $z$; $\Omega_\ell = 0$
exactly when the quad graph is bipartite. On a surface of genus zero with
borders, $\alpha_\gamma + \shift_\gamma$ has the parity of the total length of
the border components on either side of $\bar\gamma$, or of the number of times
$\bar\gamma$ crosses any system of disjoint arcs joining the border components
of odd length in pairs.
\end{theorem}
\begin{proof}
As in Proposition~\ref{prop:parity}, $\gamma \mapsto \pi(\hat h(\gamma))$ is a
homomorphism, zero on the loops round the singular vertices (quarter-turns
about whole points); with values of exponent two it factors through
$H_1(\Sigma; \Z/2)$, those loops spanning the kernel of $H_1(\Sigma^\circ)
\to H_1(\Sigma)$. For a turn $\hat h(\bar\gamma)$ is $(\alpha_\gamma,
\pm\shift_\gamma)$ up to a quarter-turn; a border component comes back
displaced by $\sum_a \pm\ell_a v_a$, of the parity of $L_j(\ell)$, and the
components add up to the boundary, zero in homology. Along a cycle of the quad
graph every edge changes the colour. In genus zero the border components span
$H_1(\Sigma; \Z/2)$ and $\bar\gamma$ is homologous to those on one side; the
last form is Lefschetz duality.
\end{proof}

On Suzanne, whose eye sockets are borders of $23$ and $25$ edges, the turn
round an eye socket is a ring of $23$ quads at the lengths of one way and of
$24$ at another, against Proposition~\ref{prop:parity}; its shift is zero at
every balanced length:

\begin{proposition}[Collars]\label{prop:collar}
Let the curve of a turn $\gamma$ and a border component $\beta_j$ without
corner bound an annulus holding no singular vertex. Then $\shift_\gamma(\ell)
= 0$ and $\alpha_\gamma(\ell) = L_j(\ell)$ for all balanced lengths.
\end{proposition}
\begin{proof}
The holonomy of $\beta_j$ is the translation by $L_j$ along an axis of the
field. Cut along a path from $\bar\gamma$ to $\beta_j$, the annulus is a disc
without cone point, so the translation of $\bar\gamma$ is $L_j u$ with $u \in
\{\pm e_1, \pm e_2\}$, as in Proposition~\ref{prop:burgers}(b). For positive
lengths $\alpha_\gamma > 0$, so $u = \pm e_1$; the identities are linear and
extend to all balanced lengths.
\end{proof}

The theorem predicts the parity of all $660$ turns of Suzanne, Jerry and a
dome, $61$ of them odd, and of every closed loop of the output meshes ($294$
on the male body); the proposition holds on all $40$ collars; on Bob,
$\Omega_\ell$ depends on the lengths, the three ways giving it $(0, 1)$, $(0,
0)$ and $(1, 0)$ on two handle loops (measured). For combinations,
$\sum_r u_r(\alpha_{\gamma_r} + \shift_{\gamma_r}) \equiv \sum_r u_r\,
\Omega_\ell([\bar\gamma_r])$; the colours of Corollary~\ref{cor:checker} may
depend on the paths, but the count on white points still gives the parity when
every curve lies in a planar region without border.

\begin{proposition}[Burgers sum with borders and handles]\label{prop:burgers-any}
Let the loops $\lambda_i$ round the singular vertices, $\beta_j$ round the
border components and $a_k, b_k$ round the handles generate the fundamental
group of $\Sigma^\circ$. For every word $w$ in them,
\[
  \tau(w) = \sum_i \rho\Bigl(\frac{\partial w}{\partial x_i}\Bigr)(I - R_i)\,p_i
  + \sum_{g \in \{\beta_j, a_k, b_k\}} \rho\Bigl(\frac{\partial w}{\partial g}\Bigr)\tau(g).
\]
For a turn, the shift is the Burgers sum of the singular vertices plus a
\emph{period term}, linear in the lengths and not determined by the developed
positions of the singular vertices. A border component without corner has
holonomy the translation by $L_j$ along its developed direction; a handle
inside the region between two turns contributes the commutator $(I -
\rho(b))\tau(a) - (I - \rho(a))\tau(b)$, zero when both its loops have trivial
rotation.
\end{proposition}
\begin{proof}
Both sides satisfy $\phi(uv) = \phi(u) + \rho(u)\phi(v)$ and agree on the
generators. The commutator of $x \mapsto Ax + a$ and $x \mapsto Bx + b$ with
commuting rotations is the translation by $(I - B)a - (I - A)b$.
\end{proof}

With its border term the sum reproduces all $314$ sides holding a border on
which it was tested, and without it none of the $586$ such sides of Suzanne
and Jerry (measured).

\begin{proposition}[Two crossing turns on a torus]\label{prop:torus}
On a closed surface of genus one, let the curves of two turns be simple and
cross once. Cut along both, the surface is a disc $D$ holding every singular
vertex, and its Burgers sum vanishes: $\sum_i Q_i (I - R_i) p_i = 0$, in the
order of first visit, for all balanced lengths. More generally, if $\kappa$ is
a simple loop crossing the curve of a turn once, with rotation $R^j$, then $(I
- R^j)\,T_\gamma = \pm B_D$: for $j \neq 0$ the translation of the turn is set
by the singular vertices, for $j = 0$ it is a free period and the singular
vertices satisfy $B_D = 0$.
\end{proposition}
\begin{proof}
The border of $D$ is the commutator of the two curves, whose holonomies are
translations and commute; Proposition~\ref{prop:burgers}(a) applies to $D$.
With $\kappa$ of rotation $R^j$, the commutator is the translation by $(I -
R^j)T_\gamma$.
\end{proof}

On Bob, of genus one with ten vertices of valence three and ten of valence
five, the Burgers sum vanishes for all $304$ pairs of simple turns crossing
once, and for none of $90$ shuffled orders (measured). With a single 3--5
pair the two cone points would develop onto one point; that no such torus
exists is the theorem of Barnette, Jucovi\v{c} and
Trenkler~\cite{barnette1971toroidal}, also a consequence of Abel's theorem
($c_5 - c_3$ is principal on an elliptic curve only if $c_5 = c_3$), an
obstruction to existence, not to closing. With $j = 0$ the proposition is the
commutator argument of~\cite[\S7]{izmestiev2013torus}. The conflict met on
Bob is of another kind: $3\shift_{33} - 2\shift_{31}$, for two turns winding
three and two times round the same handle, cancels the period and equals two
arc lengths at a vertex of valence three, one held open (measured).

Proposition~\ref{prop:linear}, Lemma~\ref{lem:gauss} for discs and annuli,
Propositions~\ref{prop:burgers} to~\ref{prop:checkerboard} (for annuli),
\ref{prop:staircase}, \ref{prop:swap}, \ref{prop:winding}
and~\ref{prop:commutator}, Remarks~\ref{rem:definite} and~\ref{rem:exist},
Theorem~\ref{thm:forced} and Corollary~\ref{cor:chain} need no closed sphere;
Lemma~\ref{lem:eight}, Proposition~\ref{prop:parity},
Corollary~\ref{cor:parity}, the count of independent rings, the setting of
Proposition~\ref{prop:fox}, Corollary~\ref{cor:checker} and
Question~\ref{conj:local} do; Appendix~\ref{app:orders} treats triangle
meshes.

\section{Turns every quantization within bounds leaves shifted}\label{sec:forced}

Turned around, the obstructions bound from below the number of turns any
quantization in a set $P$ of lengths leaves shifted.

\begin{theorem}[Forced shifts]\label{thm:forced}
Let $H$ be a finite set of turns of a T-mesh and $P$ a non-empty set of lengths
containing every quantization considered: for instance the balanced lengths
meeting given bounds of separation, $\ell \ge b$, possibly with a cap $\frac1n
\sum_a |\ell_a - t_a| \le D$ on the mean deviation from the targets $t$, or with
a given number of quads. Let
\[
  \tau_P(H) = \min\bigl\{\, |X| : X \subseteq H,\ \text{the turns of } H
  \setminus X \text{ can be closed together by lengths in } P \,\bigr\}.
\]
Then for every $\ell \in P$, whole or not, at least $\tau_P(H)$ turns of $H$
have $\shift_\gamma(\ell) \neq 0$. If $H$ contains $k$ pairwise disjoint subsets
none of which can be closed within $P$, then $\tau_P(H) \ge k$.
\end{theorem}

\begin{proof}
At $\ell \in P$ the turns with zero shift are closed by $\ell$ itself, so the
shifted ones form a set $X$ as in the definition. A set whose removal leaves a
closable set meets each of the $k$ disjoint subsets.
\end{proof}

The theorem is elementary; its content is in what makes subsets infeasible.
It uses neither convexity, nor linearity, nor whole lengths. The minimal
non-closable sets are irreducible infeasible subsystems (IIS), $\tau_P(H)$ is
a minimum IIS cover~\cite{parker1996iis,amaldi2003maxfs}, and the theorem is
the packing bound of that covering problem.

\begin{corollary}[Chains of dislocations]\label{cor:chain}
Let $\gamma_0, \gamma_1, \dots, \gamma_m$ be turns and suppose that for $i = 1,
\dots, m$
\[
  \shift_{\gamma_{i-1}}(\ell) - \shift_{\gamma_i}(\ell) = \varepsilon_i\,
  F_i(\ell) \quad \text{for all balanced } \ell,
\]
with $\varepsilon_i \in \{+1, -1\}$ and forms $F_i$ that are positive on the
whole points of $P$; for instance $F_i$ the length of a staircase
(Proposition~\ref{prop:staircase}) between a vertex of valence three and one of
valence five lying between $\gamma_{i-1}$ and $\gamma_i$
(Propositions~\ref{prop:dislocation} and~\ref{prop:burgers}(b)), the
staircase containing an arc held at length at least one. Let $q$ be the number
of indices $i < m$ with $\varepsilon_{i+1} \neq \varepsilon_i$. Then for every
whole $\ell \in P$ at most $q + 1$ of the turns $\gamma_i$ are rings: at least
$m - q$ are helices. In particular $m + 1$ nested parallel turns separated by
$m$ dislocations whose staircases all climb the same way leave at least $m$
helices.
\end{corollary}

\begin{proof}
If $\gamma_i$ and $\gamma_j$, $i < j$, are rings, $\sum_{k=i+1}^{j}
\varepsilon_k F_k(\ell) = 0$ with $F_k(\ell) > 0$, so a change of sign lies
between $i$ and $j$; between two consecutive rings lies one of its own.
\end{proof}

\begin{remark}[What the bound counts]
The theorem counts turns with a non-zero shift; at lengths where a turn is not
a strip of the mesh~\cite[Prop.~2]{gardent2026loops}, or where one spiral
strip carries several turns of $H$, it is not a count of helices.
\end{remark}

The bound is attained, by the standard correspondence between minimum IIS
covers and maximum feasible subsystems~\cite{amaldi2003maxfs,parker1996iis}:

\begin{theorem}[Forced shifts as a transversal number; standard]\label{thm:transversal}
Call $Y \subseteq H$ \emph{closable} in $P$ if some $\ell \in P$ has
$\shift_\gamma(\ell) = 0$ for all $\gamma \in Y$, and a \emph{core} a set that
is not closable while all its proper subsets are. For any $P$:
(i) $H \setminus X$ is closable if and only if $X$ meets every core; hence
$\tau_P(H)$ is the least size of a set meeting every core;
(ii) for every $\ell \in P$ the turns $\ell$ leaves shifted meet every core;
(iii) every least such set is exactly the set of turns left shifted by some
$\ell \in P$;
(iv) $P \subseteq P'$ implies $\tau_P(H) \ge \tau_{P'}(H)$.
\end{theorem}

\begin{proof}
Closability is hereditary: if $X$ misses a core, $H \setminus X$ contains it,
and a non-closable $H \setminus X$ contains a core. The zero-shift set of
$\ell$ is closed by $\ell$. For a least $X$, a length closing $H \setminus X$
shifts a subset of $X$ meeting every core, hence all of $X$.
\end{proof}

So $\tau_P(H)$ lies between the largest number of disjoint cores and the
least hitting set of the cores found so far, exact once its complement is
closable; the \emph{implicit hitting-set} loop of Parker and
Ryan~\cite{parker1996iis} alternates the two, with
cores from the deletion filter of~\cite{chinneck1991locating}.

\subsection{Whole lengths and size}\label{sec:whole}

\begin{proposition}[Whole lengths and size without a cap]\label{prop:whole-size}
If $P$ is given by the balances and bounds $\ell \ge b \ge 0$, a set closable
in real lengths is closable in whole lengths, and closable in real lengths at
any number of quads greater than $\inf_{P_Y} Q$, and at $\inf_{P_Y} Q$ itself
when the infimum is attained, where $P_Y$ is the set of lengths of $P$ closing
$Y$ and $Q(\ell)$ the number of quads.
\end{proposition}

\begin{proof}
Scale a rational point of $P_Y$ by its denominator (Remark~\ref{rem:exist}).
For a size $q$, take $\ell \in P_Y$ with $0 < Q(\ell) \le q$; $c\ell \in P_Y$
for $c \ge 1$, and $Q(c\ell) = c^2 Q(\ell)$ reaches $q$.
\end{proof}

With a cap on the deviation, $P_Y$ is a polytope and $Q(P_Y)$ an interval: two
exact points on either side of a window of sizes certify that $Y$ closes
within it, while the converse needs the extremes of $Q$ over a polytope. The
cap is a budget shared by all turns: on the body the lengths may move by
$980.7$ edges in all, while closing the $109$ strips of way~2 that the bounds
do not force costs $1{,}680.2$, and the cores become covers of a knapsack
(measured).

On a closed surface where no patch borders itself, an arc lies on a side of
each of the two patches it borders, so its column of $B$ has two entries
$\pm1$: $B$ is the incidence matrix of a \emph{bidirected graph} $G$ whose
nodes are the balances and whose edges are the arcs, an arc on a border giving
a half-edge, the setting of~\cite{heistermann2023bimdf}, with the rank of
Zaslavsky~\cite{zaslavsky1982signed}.
A cycle of $G$ is \emph{balanced} if values $\pm1$ on its edges balance each
of its nodes, and \emph{unbalanced} otherwise.

\begin{lemma}[Handshake]\label{lem:handshake}
Let $x \in \Z^n$ and $C$ be a connected component of $G$. Then $\sum_{r \in C}
(Bx)_r \equiv \sum_{a} x_a \pmod 2$, the second sum running over the border
arcs of $C$. On a closed surface, $\sum_r (Bx)_r$ is even.
\end{lemma}

\begin{proof}
It is $\sum_a x_a \sum_{r \in C} B_{ra}$, and each column sums to $0$ or
$\pm2$, or to $\pm1$ for a border arc.
\end{proof}

So on a closed surface a system $Bx = d$ with $\sum_r d_r$ odd has real
solutions but no whole one; this does not affect the balanced lengths:

\begin{proposition}[Free arcs]\label{prop:free}
Let $G$ be connected and contain an unbalanced cycle. Let $T$ be the edges of a
spanning tree of $G$ and one more edge closing an unbalanced cycle, and $F$ the
other edges. Then $B$ has full row rank, and every $x_F \in \Z^F$ extends to
exactly one $x \in \Z^n$ with $Bx = 0$. The whole balanced lengths thus form a
lattice isomorphic to $\Z^F$, of dimension $n - m$ for $m$ balances, with a
basis of vectors $z_f$, $f \in F$, each supported on $T \cup \{f\}$ with
$z_f(f) = 1$ and entries in $\{0, \pm1, \pm2\}$.
\end{proposition}

\begin{proof}
$T$ exists: the sign of a cycle is a homomorphism on the cycle space over
$\Z/2$, which the fundamental cycles of a spanning tree span, so if they were
all balanced every cycle would be. $T$ has $m$ edges. If $B_T y = 0$, a node
met by a single edge of $T$ forces that edge to zero; pruning such edges
leaves the unbalanced cycle, on which each value fixes the next up to sign,
and going round forces all to zero. So $B_T$ is invertible, $B$ has full row
rank, and $x_F$ has exactly one real extension. For $f \in F$, $T \cup \{f\}$
has a one-dimensional space of balanced values. Pruned, it is either three
paths between two nodes, whose cycles are the unbalanced cycle of $T$, one more
unbalanced and one balanced (the number of unbalanced cycles of such a graph
is even), the balanced one containing $f$, since the cycle of $T$ does not,
and carrying values $\pm1$; or two cycles joined by a path, possibly of length
zero, the cycle through $f$ being balanced and carrying $\pm1$, or unbalanced,
and then $\pm1$ on both cycles and $\pm2$ on the path are balanced. Either way
there is a whole balanced $z_f$ supported on $T \cup \{f\}$ with $z_f(f) =
\pm1$; with signs fixed, $x = \sum_{f \in F} x_f z_f$ is the whole extension.
\end{proof}

With a spanning forest and, in each component with an unbalanced cycle, an
edge closing one or a half-edge, the proposition extends beyond the closed
sphere, where full row rank can fail (on part, $145$ of $148$). On the male
body $G$ is connected with an unbalanced cycle and the lattice of balanced
lengths has dimension $2{,}069 - 1{,}552 = 517$: balances alone never obstruct
whole lengths, and integer elimination ends with a single row of greatest
common divisor $2$. The shifts of the held turns leave lattices of dimension
$363$ to $469$, and in ways 4, 5 and 6 elimination needs steps of Euclid's
algorithm: turn rows can destroy the free-arc structure (measured). The one
integer obstruction found on the body is a rounding gap in deviation, not a
parity: rings that close in real numbers at a mean deviation of $0.399$ need
at least $0.536$ in whole numbers~\cite[\S7]{gardent2026loops}.

\subsection{On the male body}

In Table~\ref{tab:forced}, $H$ is the held turns of a way, or the $111$
strips of the engine's mesh for way~2 leaning less than $0.12$, and $P$ is
$P_{\mathrm{eng}}$, then also the engine's cap $D = 0.474$ on the mean
deviation, then also whole lengths and its window of $10{,}379$ to $10{,}801$
quads. With the bounds alone the implicit hitting-set loop is exact in at most
$66$ seconds per family; with the cap and the window, bounds come from
closable subsets, disjoint cores, solver bounds and whole points checked
exactly.

\begin{table}[t]
\centering
\small
\begin{tabular}{lrrrrr}
\toprule
way & turns & bounds & bounds, cap & cap, size, whole & left shifted by the engine \\
\midrule
0 & 48 & 1 & 1 & 1 & 13 \\
1 & 78 & 5 & 5 & 5--6 & 39 \\
2 & 52 & 1 & 1 & 1 & 16 \\
3 & 86 & 4 & 4 & 5--18 & 44 \\
4 & 84 & 1 & 1 & 1--12 & 49 \\
5 & 160 & 6 & 7--17 & 7--86 & 115 \\
6 & 88 & 1 & 1 & 1--17 & 44 \\
7 & 169 & 9 & 9--26 & 9--90 & 123 \\
\midrule
2, strips of the mesh & 111 & 2 & 3--11 & 3--43 & 81 \\
\bottomrule
\end{tabular}
\caption{Male body: turns of each family that every quantization in $P$ leaves
shifted (Theorems~\ref{thm:forced} and~\ref{thm:transversal}; measured), with
$P = P_{\mathrm{eng}}$, then also the engine's cap (real lengths), then also
whole lengths and the engine's window of sizes. A single number is exact; a
range is certified, below by the solver's proven bound or disjoint cores,
above by a closable subset found within the time limit ($120$ to $1{,}800$
seconds). Last column: turns of the family shifted at the engine's whole
lengths. Turns are counted up to sign.}
\label{tab:forced}
\end{table}

With the bounds alone the forced turns are exactly $1$, $5$, $1$, $4$, $1$,
$6$, $1$ and $9$ (Table~\ref{tab:forced}); in ways 0, 2, 4 and 6 the forced
turn is the wrist; in way~7 nine turns are needed although at most
eight of its $21$ cores are disjoint, so the packing bound is not always
attained. With the cap and the window, at whole lengths, the value is still
exactly $1$ in ways 0 and 2: whole lengths closing every held turn but the
wrist exist at $10{,}439$ and $10{,}436$ quads, at mean deviations $0.442$ and
$0.469$. In way~3 integrality
costs at least one turn ($4$ in real numbers, at least $5$ in whole numbers). Of the
$16$ held turns of way~2 the engine leaves shifted, the bounds force $1$, the
cap, the size and integrality none: the heuristic accounts for $15$; in the
other ways it leaves at least $12$ to $37$ turns more than the best lengths
found. The cap is the engine's only approximately (it refused $0.4737$), and optima
that bring other arcs to zero can break the mesh; keeping open every arc the
engine keeps open, the optimum for way~2 is $4$ (measured).

\section{Choosing the turns}\label{sec:choosing}

Call \RC{} the problem: given the balances $B$, bounds $b$, a set $H$ of
linear forms (the shifts) and a number $k$, are there lengths in $P = \{\ell :
B\ell = 0,\ \ell \ge b\}$ that close at least $k$ forms of $H$? The maximum
feasible subsystem problem is NP-hard and hard to approximate in
general~\cite{amaldi1995complexity}, and minimum IIS covers are strongly
NP-hard already for network flows~\cite{joormann2015flowiis}; the hardness
survives the structure of quantization, a case new as far as we found.

\begin{theorem}[Choosing the turns is NP-hard]\label{thm:nphard}
\RC{} is NP-complete, in real or in whole lengths, already when $B$ is the
incidence matrix of a bidirected graph, the bounds are $0$ or $1$, and the
forms have coefficients $0$ or $1$; it is strongly NP-hard, and stays NP-hard
with a cap on the deviation, which may be loose. Moreover $\tau_P(H)$ can
equal the vertex cover number of any graph, with $|H|$ its number of vertices.
So it cannot be approximated within $\sqrt2 - \varepsilon$ unless
P${}={}$NP~\cite{khot2017independent,dinur2018towards,dinur2018nonoptimally,khot2018grassmann},
nor within $2 - \varepsilon$ if the unique games conjecture
holds~\cite{khot2008vertex}, and the largest closable subset of $H$ cannot be
approximated within $|H|^{1-\varepsilon}$ unless
P${}={}$NP~\cite{hastad1999clique,zuckerman2007linear}.
\end{theorem}

\begin{proof}
Membership: closing a given set of $k$ forms is a linear program, and by
Proposition~\ref{prop:whole-size} real and whole closability agree. Hardness,
from \textup{\textsc{Independent Set}}: for a graph $G$, take for each edge $e
= uv$ four arcs $y_{u,e}$, $y_{v,e}$, $d_e$, $d'_e$ with the balances
\[
  y_{u,e} + y_{v,e} - d_e = 0, \qquad y_{u,e} + y_{v,e} - d'_e = 0, \qquad
  d_e - d'_e = 0 ,
\]
in which every arc appears in exactly two rows with coefficients $\pm1$, the
bounds $d_e, d'_e \ge 1$ and $y \ge 0$, and for each vertex $u$ the form
$\shift_u = \sum_{e \ni u} y_{u,e}$. With $y \ge 0$, $\shift_u = 0$ forces
$y_{u,e} = 0$ at $u$, so closing two adjacent vertices gives $d_e = 0 < 1$.
For an independent set $Y$, every edge has at most one end in $Y$: put $y = 0$
on its side and $1$ on the other, or, with no end in $Y$, $1$ on either side
and $0$ on the other; then $d_e = d'_e = 1$, whole lengths in $P$ closing $Y$.
The closable sets are the independent sets of $G$, and $\tau_P(H) = |V| -
\alpha(G)$ is the vertex cover number. The instance is linear in $G$, with
numbers $0$ and $\pm1$; a large cap constrains nothing.
\end{proof}

The reduction was checked by computer on $40$ random graphs (measured); its
third balance puts every arc in two balances, as in a T-mesh. It is not known
to be a T-mesh: its forms have non-negative coefficients and are not built as
the holonomies of straight walks, so the theorem says that quantization's
algebra allows hard instances, not that a surface produces them.

\begin{conjecture}[Realizable instances]\label{conj:realize}
\RC{} stays NP-hard when $B$, $b$ and $H$ are the balances, the bounds of
separation and the turns of a T-mesh of a closed surface of genus zero.
\end{conjecture}

The evidence is that vertex cover stays NP-hard on planar graphs of small
degree~\cite{garey1976simplified}, and that two turns separated by a
dislocation conflict like the two ends of an edge.

Some parts are easy: without turns and window, whole quantization is a
polynomial matching problem~\cite{edmonds1970matching}; with disjoint known
cores $\tau_P(H)$ is their number; with shift forms of fixed rank the closable
sets are the flats of an arrangement, but on the body the held turns of a way
have full rank.

\begin{remark}[The window of sizes; a sketch]
The window on the number of quads makes even closing nothing hard in the weak
sense. Take a disjoint union of tori, each of two patches whose long and short
sides $a_i$ and $c_i$ have the targets $N v_i$ and $1$, so that $Q = \sum_i 2
a_i c_i$, and a cap that lets the
short sides $c_i$ grow by $T_0$ edges in all. For $N$ large, the window $Q =
2N(\sum_i v_i + T_0)$ is met exactly when $\sum_i v_i(c_i - 1) = T_0$ with
whole $c_i \ge 1$: unbounded subset sum, which is
NP-complete~\cite{lueker1975two}, with numbers exponential in the size of the
instance (the gadgets are of genus one and disconnected). Whether the window
alone is strongly NP-hard, with $Q$ polynomial in the size, is open, as is the
complexity of maximizing $Q$ over a polytope of lengths in reals.
\end{remark}

\section{What is machine-checked}\label{sec:lean}

The folder \texttt{lean} beside this note is a Lean~4~\cite{demoura2021lean4}
project on Mathlib~\cite{mathlib2020} (both \texttt{v4.35.0}). Its $235$
theorem declarations, many of them small lemmas, depend only on Lean's
standard axioms; the count is not evidence, and its \texttt{README.md} maps
each statement to its Lean theorems. Lean checks the algebraic and
combinatorial core of each statement, with the topological facts it uses
(that a loop develops into a given motion, that a curve bounds a disc, that
the loops $\lambda_i$ present the fundamental group, that the quad graph of a
sphere is bipartite) as explicit hypotheses. Each statement is formalized on
its own abstraction (a turn as a list of motions, a region as finite sets with
counts, a system as a pair of matrices): the files share no type of T-mesh
and are not composed into one model of it.

\emph{Checked}, statement by statement. For motions with linear parts in any
ring: the Burgers sum, and that the recursion of the gluing maps of a tree of
cuts, a hypothesis, unfolds into the order of first visit
(Proposition~\ref{prop:burgers}(a)); the cut annulus~(b); exchanges; the
conjugation step of Proposition~\ref{prop:winding}. Over the reals: the four
orders and orientations of a dislocation and straight chains
(Proposition~\ref{prop:dislocation}(i), (ii)); staircases; the positivity of
a non-negative form with a held arc. At whole lengths: the checkerboard rule;
Proposition~\ref{prop:parity} for whole balanced lengths of any sign, given a
2-colouring at positive lengths and the linearity of $\alpha_\gamma$ and
$\shift_\gamma$; the second proof of Remark~\ref{rem:borders}. The
Fox--Burgers formula on the free group with the twisted windings and the
insertion of the relation, the case of a simple turn,
Proposition~\ref{prop:commutator} and Corollary~\ref{cor:checker}.
Lemma~\ref{lem:handshake} on a closed surface; Proposition~\ref{prop:free},
stronger than stated (on the edges $T$ the whole image of $B$ is the lattice
of vectors of even sum); Remarks~\ref{rem:definite} and~\ref{rem:exist}, with
Farkas' lemma over any ordered field and whole multipliers;
Theorem~\ref{thm:forced} and Corollary~\ref{cor:chain}. For fields of any
order: the rotation part of Gauss--Bonnet and its counts on triangle and
hexagonal meshes, the lengths of dislocations and the congruence of
Appendix~\ref{app:orders}.

\emph{Checked in a weaker form}: Proposition~\ref{prop:linear}, only the
additivity of the composed translation, not the integer coefficients nor the
reading as a helix; Lemmas~\ref{lem:gauss} and~\ref{lem:eight}, only the
Euler-characteristic count on the grid refined by two and the totals from the
numbers of convex and concave corners, the regions being discs and the
convexity of their corners hypotheses; the completeness of the walks,
Proposition~2 of~\cite{gardent2026loops}, only in a model: on a closed
surface (opposite sides of a patch of equal length, no two unit edges glued
onto the same one) the station walk of the gluing data is a permutation, the
strips are \emph{defined} as the cycles of the successor it gives on the quads
above the stations, and they are in bijection with the cycles of the real walk
from the half-integers $k + \frac12$, in a one-dimensional model of the
gluing; that patch sides are positive, that crossing a patch keeps the offset
and that the model is the quad mesh are assumptions.
\emph{Not checked}: the topology taken as hypotheses, strips with borders, the
measured statements, and those proved on paper only (see AI assistance). The formalization found no error of sign or
convention; it sharpened the hypotheses of Propositions~\ref{prop:parity},
\ref{prop:free} and~\ref{prop:staircase}, of Remark~\ref{rem:definite} and of
the completeness proposition, stated in their sharpened form.

\section{Open problems and limitations}\label{sec:open}

When is a set of turns closable, in terms of the singularity configuration?
Every certificate is a sum of combs, which says nothing about the
singularities, and not of staircases, caps and winding separatrices alone; a
useful answer lies between, for instance combs on disjoint sets of vertices,
each of the sign of the certificate. Computing $\tau_P(H)$ at a given number of quads needs the extremes of an
indefinite quadratic function over a polytope. Is \RC{} NP-hard for a fixed number of singular vertices? And when
the set of turns depends on the lengths, as the strips of the mesh do, what is
the problem at all? When does a quartic differential with all trajectories finite exist for a
given singularity configuration~\cite{lei2020quartic}? Turns are its discrete
version. Finally, a formal model of a complex
of rectangles with integer sides, its turns and its developing map would turn
the topological hypotheses of the Lean statements into theorems.

\paragraph{Limitations and reproducibility.}
The measured statements come from one T-mesh of one body, exported by an
engine without the constraint rounds of the companion paper, and from a few
bench models with an unrecorded engine version; identities were checked in
floating point unless exact arithmetic is stated, and the homotopies of
Propositions~\ref{prop:swap} and~\ref{prop:winding} were not checked. Forced
counts are relative to the engine's bounds of separation and, approximately,
to its cap, and the intervals depend on solver time limits on a shared
machine. The scripts and logs of the exact experiments are not published: the
measured statements can be checked against the branch \texttt{rd-exact}, an
experimental hook of openclare, but not rerun as they stand. The Lean proofs
replay with \texttt{make proofs NOTE=turns-of-a-t-mesh}. This version of the
note was built from commit \texttt{\enginecommit}.

\section*{AI assistance}
This note, its Lean proofs and the software it describes were made by a human
author working with an AI system, which is not an author. The author chose the
problem and its aims (edge loops that close into rings and, as an artist, few
singular vertices and few spirals rather than perfectly square quads),
directed the work, decided which results to pursue, and judged them through
the measurements and drafts presented to him. Claude Opus 5.5 (Anthropic),
used through Claude Code in October 2026 with several sub-agents, wrote the
changes to the engine and the experiment scripts, ran the exports and exact
computations, searched the literature, stated and proved the propositions,
wrote the Lean formalization and drafted the text. Each bibliography entry says how much of the work was read.
Beyond the cores machine-checked in Lean (Section~\ref{sec:lean}), everything
was written with AI help and checked by AI agents rereading it and testing it
on the data, not by a human expert. In particular, Propositions~\ref{thm:combs},
\ref{prop:pairing}, \ref{prop:collar}, \ref{prop:burgers-any}, \ref{prop:torus}
and~\ref{prop:whole-size}, Lemmas~\ref{lem:comb} and~\ref{lem:gauss-border},
Theorems~\ref{thm:parity-any}, \ref{thm:transversal} and~\ref{thm:nphard}, the
case of Lemma~\ref{lem:handshake} with borders, and the statements on
rotations of order six other than the checked ones have only a paper proof
written with AI help. AI review found errors in Propositions~\ref{thm:combs}
and~\ref{prop:pairing} and Theorem~\ref{thm:nphard}, and in
Proposition~\ref{prop:linear} and the completeness proposition
of~\cite{gardent2026loops}, whose Lean checks are weaker than the statements;
all are corrected here. The note has not been peer reviewed; the author, who
publishes it, is responsible for its content.

\bibliographystyle{plainnat}
\bibliography{refs}

\appendix

\section{Rotations of other orders}\label{app:orders}

On triangle and hexagonal meshes the rotations are powers of $\zeta =
e^{i\pi/3}$, a vertex of valence $k$ has index $6 - k$, and a turn along $u$
comes back translated by $T = \alpha u + \shift\, \zeta u$.
The Burgers sum, the order of first visit and the Fox--Burgers formula do not
use the order of the rotations and hold unchanged, as do
Theorem~\ref{thm:forced}, Corollary~\ref{cor:chain} and
Lemma~\ref{lem:handshake}, which use no geometry. A pair of valences five and seven has $B =
\pm\zeta^{\pm1} d$, and the shifts of two turns differ by $\pm d_1$ or
$\pm(d_1 + d_2)$, $d = d_1 u + d_2 \zeta u$. What changes is parity. When every
rotation is a power $z^{j_i}$ of one rotation $z$, every Fox derivative is
congruent to its exponent sum modulo $1 - z$, and $1 - z^j$ to $j(1 - z)$
modulo $(1 - z)^2$, so in any commutative ring
\[
  \tau(w) \equiv (1 - z) \sum_i \epsilon\Bigl(\frac{\partial w}{\partial
  x_i}\Bigr) j_i\, p_i \pmod{(1 - z)^2}.
\]
For quarter-turns $(1 - z)^2 = -2z$, and this is Corollary~\ref{cor:checker}.
For order six $1 - z$ is a unit: no colouring rule on triangle or hexagonal
meshes. When every singular vertex has even valence the rotations have order
three, $(1 - \omega)^2 = -3\omega$ (compare Lemma~11
of~\cite{izmestiev2013torus}), and on a closed sphere $\alpha \equiv \shift
\pmod 3$: every ring has a length divisible by three, as on subdivided
octahedra, while on subdivided icosahedra rings of $5n$ edges take every
residue (measured); a rule modulo three with borders or handles is a
conjecture. A dislocation of order $n$ has $|B| = 2\sin(\pi/n)\,|d|$.

\end{document}
 gives the paper
% without them, for a venue that wants it plain.
\newcommand{\lean}[1]{\texttt{\detokenize{#1}}}
\ifdefined\plain
\newenvironment{evidence}[3][]{}{}
\else
\usepackage[most]{tcolorbox}
\newcommand{\evlevel}[2]{\expandafter\newcommand\csname ev-name-#1\endcsname{#2}}
\evlevel{lean}{Lean, complet}
\evlevel{core}{Lean, cœur algébrique}
\evlevel{data}{Vérifié sur les données du moteur}
\evlevel{measured}{Mesuré}
\evlevel{prose}{Démonstration sur papier seulement}
\evlevel{tested}{Conjecture, vérifiée par ordinateur}
\evlevel{cited}{Classique, cité}
\evlevel{open}{Ouvert, rien n'est affirmé}
\newcommand{\evname}[1]{\csname ev-name-#1\endcsname}
\tcbset{evstyle-lean/.style={}, evstyle-core/.style={}, evstyle-data/.style={},
  evstyle-measured/.style={}, evstyle-tested/.style={}, evstyle-prose/.style={},
  evstyle-cited/.style={},
  evstyle-open/.style={frame hidden, borderline={0.4pt}{0pt}{black!45, dashed}}}
\newtcolorbox{evidence}[3][]{breakable, enhanced, colback=white,
  colframe=black!45, boxrule=0.4pt, arc=0pt, left=4pt, right=4pt, top=2pt,
  bottom=2pt, before skip=6pt, after skip=6pt, evstyle-#2,
  title={\footnotesize\textsc{\evname{#2}}\if\relax\detokenize{#3}\relax\else
    \quad #3\fi},
  coltitle=black, colbacktitle=white, titlerule=0.4pt, toptitle=1pt,
  bottomtitle=1pt, #1}
\fi
\title{Les tours d'un T-mesh~:\\
obstructions à la fermeture des boucles d'arêtes\\[0.5em]
\large Note de recherche --- brouillon, non relue par des pairs}
\author{Marc Gardent%
\thanks{Cette note est diffusée sous la licence Creative Commons Attribution
4.0 International (CC BY 4.0). Le code d'openclare et les preuves Lean sont
sous la licence GPL-3.0-or-later du dépôt. La note décrit les preuves Lean au
commit \texttt{\enginecommit} de
\url{https://github.com/marcgardent/openclare}.}\\
\normalsize Chercheur indépendant\\
\normalsize ORCID~: \href{https://orcid.org/0009-0000-6466-3098}{0009-0000-6466-3098}\\
\normalsize Contact~: par le dépôt,
\url{https://github.com/marcgardent/openclare}}
\date{Octobre 2026}

\begin{document}
\maketitle

\begin{center}
\fbox{\parbox{0.9\textwidth}{\small\textbf{Avertissement.} Ce texte est une
traduction française, faite au mieux (\emph{best effort}) et avec l'aide
d'une IA, de la version anglaise de la note. En cas de différence, c'est la
version anglaise qui fait foi~: \emph{Turns of a T-Mesh: Obstructions to
Closing Edge Loops}, dans le même dépôt,
\nolinkurl{docs/research/turns-of-a-t-mesh/paper.pdf}.}}
\end{center}

\begin{abstract}
Dans la quantification d'un T-mesh, une boucle d'arêtes se ferme en anneau
lorsqu'un \emph{tour}, marche en ligne droite à travers les patchs qui revient
à son premier côté, a un \emph{décalage} nul, le décalage étant une forme
linéaire à coefficients entiers des longueurs des arcs\,; sinon, elle
s'enroule en hélice. Nous étudions pourquoi les tours ne peuvent pas tous être
fermés. En lisant le décalage comme l'holonomie de translation de la métrique
plate, nous montrons que, pour un tour simple, il est la composante transverse
du vecteur de Burgers des sommets singuliers situés d'un côté\,; qu'une paire
3--5 ou un échange de deux sommets, reliés par un escalier d'arcs, force l'un
de deux tours à être décalé, et qu'une séparatrice qui fait une fois le tour
d'un tour et s'arrête sur une autre séparatrice de son sommet force ce seul tour
à être décalé\,; et que, sur une sphère fermée, le décalage et la longueur d'un
tour ont la même parité. Le calcul de Fox étend la somme de Burgers à toute
combinaison de tours, et des \emph{peignes}, dislocations locales portées par
un seul arc, écrivent tous les certificats de conflit trouvés sur un corps
humain. Les conflits minorent le nombre de tours que toute quantification
respectant des bornes données laisse décalés~: sur ce corps, dans les bornes
de séparation d'un moteur de remaillage, exactement un à neuf des tours qu'il
retient, là où il en laisse de 13 à 123 décalés. Choisir les tours à fermer
est NP-difficile pour des équilibres bidirigés avec des bornes $0$--$1$, la
forme algébrique de la quantification\,; savoir si de vrais T-meshes donnent
des instances difficiles est une question ouverte. Le cœur algébrique des démonstrations est vérifié en Lean\,; la topologie de la
surface, que les marches et les régions soient celles d'un T-mesh, y est une
hypothèse.
\end{abstract}

\section{Introduction}

Le remaillage automatique en quads produit souvent des boucles d'arêtes qui
s'enroulent en hélice autour d'une partie de la forme au lieu de se refermer
en anneaux~\cite{bommes2011structure}. Qu'une boucle se ferme ou non se décide
lorsqu'un quantificateur attribue des longueurs entières aux arcs d'un T-mesh
(maillage en T) tracé dans une paramétrisation sans
couture~\cite{campen2015qgp,lyon2021constrained,heistermann2023bimdf}~: une
boucle suit un \emph{tour}, marche en ligne droite à travers les patchs qui
revient à son premier côté, et se ferme lorsque le \emph{décalage} du tour,
une forme linéaire à coefficients entiers des longueurs des arcs, s'annule. Un
article compagnon~\cite{gardent2026loops} ajoute ces équations à un
quantificateur\,; sur des modèles réels, elles ne peuvent pas toutes être
satisfaites en même temps que les bornes qui gardent les sommets singuliers
séparés. Le présent article demande pourquoi, et ce qui s'ensuit.

Les sections~\ref{sec:setting} à~\ref{sec:beyond} vont de l'holonomie d'un
tour aux conflits entre tours, à toute combinaison de tours, puis aux bords et
aux anses\,; les sections~\ref{sec:forced} et~\ref{sec:choosing} minorent le
nombre de tours que toute quantification laisse décalés et montrent que les
choisir est NP-difficile\,; les sections~\ref{sec:lean} et~\ref{sec:open}
disent ce qui est vérifié par machine et ce qui reste ouvert. Les ingrédients
sont classiques (holonomie des métriques plates à singularités coniques,
vecteurs de Burgers, parité des quadrangulations, calcul de Fox, couvertures
de sous-systèmes infaisables)\,; leurs énoncés pour les tours, le minorant du
nombre d'hélices et la difficulté pour les instances de la quantification sont
les nôtres, pour autant que nous ayons pu le vérifier. \ifdefined\plain
Un énoncé est dit
\emph{prouvé} lorsqu'il est accompagné d'une démonstration, \emph{mesuré}
lorsqu'il provient d'un calcul en arithmétique exacte ou de la sortie d'un
solveur sur les systèmes qu'exporte le remailleur openclare, et
\emph{conjecturé} sinon.
\else
Chaque énoncé est encadré, et son étiquette dit ce qui le vérifie~:
\textsc{Lean, complet}, démontré en Lean sans autre hypothèse que celles de
l'énoncé\,; \textsc{Lean, cœur algébrique}, l'algèbre démontrée en Lean et la
topologie de la surface, que les marches et les régions sont celles d'un
T-mesh, supposée comme le dit l'étiquette\,; \textsc{vérifié sur les données
du moteur}, exactement, sur le T-mesh qu'exporte le remailleur openclare,
contre les définitions Lean\,; \textsc{mesuré}, par un script ou un solveur
sur cet export\,; \textsc{démonstration sur papier seulement}, celle qui
suit, qu'aucune machine n'a lue\,; \textsc{conjecture, vérifiée par
ordinateur}, vérifiée exactement sur des T-mesh aléatoires qui ne viennent
pas du moteur, l'étiquette disant sur combien d'instances et de quel genre
(les tests de propriétés du paquet de données,
\texttt{experiments/property})~: un test réfute ou corrobore et ne démontre
rien, et une étiquette d'un autre niveau dit \emph{testé} là où les mêmes
tests ont aussi tourné\,; \textsc{classique, cité}\,; et
\textsc{ouvert, rien n'est affirmé}, une question ou une conjecture. Les noms
Lean sont ceux du commit de la première page, construit sans
\texttt{sorry}.
\fi

Ce travail est né d'un projet de programmation, openclare, un remailleur en
quads libre, et il est devenu une recherche quand l'ingénierie seule ne
suffisait plus à dire quels résultats pouvaient être garantis et quels calculs
pouvaient être épargnés. Les résultats répondent ici du côté de ce qui ne peut
pas être fait~: une réparation qui garde la longueur d'un tour de décalage
impair ne peut pas le fermer sur une sphère fermée\,; un tour dont le décalage
est une forme de signe fixe avec un arc maintenu ouvert reste une hélice, et
il est inutile de chercher à le fermer\,; les tours que toute quantification
laisse décalés mesurent ce qu'une heuristique peut encore gagner\,; et la
difficulté de leur choix explique qu'un moteur s'en tienne à une heuristique.

\section{Travaux antérieurs}\label{sec:related}

Les valences d'un maillage de quads fermé de genre $g$ vérifient $\sum_v (4 -
k(v)) = 8(1-g)$, ce qui dit combien de sommets singuliers il faut, non où les
placer. Un maillage de quads fermé induit une différentielle quartique dont
les zéros et les pôles sont ses sommets singuliers\,; Lei et
al.~\cite{lei2020quartic} citent comme problème ouvert la garantie que toutes
ses trajectoires soient finies, et une trajectoire qui n'est pas finie est une
boucle d'arêtes qui ne se ferme pas. Une différentielle quadratique dont
toutes les trajectoires non critiques se ferment est de
Jenkins--Strebel~\cite{strebel1984quadratic}, une réunion de cylindres
plats\,; une hélice est un cylindre recollé avec un décalage d'un nombre
entier de rangées. Chen et al.~\cite{chen2019metric} caractérisent les
métriques plates des maillages de quads, en supposant que les lignes du champ
sont des boucles finies. Izmestiev et al.~\cite{izmestiev2013torus} démontrent
des résultats de non-existence pour les maillages de quads, parmi lesquels le
tore avec un sommet de valence trois et un de valence
cinq, qu'ils rappellent comme connu, au moyen de l'holonomie de l'application
développante, dont la translation, lorsque la rotation est triviale, est le
vecteur de Burgers de la cristallographie\,; une paire de disclinaisons
opposées est une dislocation~\cite{kleman2008disclinations}. Bommes et
al.~\cite{bommes2011structure} retirent les hélices de maillages finis\,; les
paramétrisations en bandes des tubes~\cite{kalberer2011stripe} arrondissent à
un nombre entier le saut autour d'un tube, notre décalage. Nous n'avons trouvé
aucun minorant du nombre d'hélices pour une configuration de singularités
donnée\,; le nôtre est une couverture minimale de sous-systèmes infaisables
irréductibles~\cite{parker1996iis,amaldi2003maxfs}.

\section{Tours, décalages et holonomie}\label{sec:setting}

Un T-mesh est une partition de la surface en rectangles, les \emph{patchs},
dont les côtés sont faits d'\emph{arcs}\,; les nœuds où une trace s'arrête sur
une autre sont des jonctions en T, et les sommets singuliers du champ sont des
nœuds. Chaque patch $P$ a des coordonnées locales dont l'origine est son
premier coin et dont les axes portent ses côtés. Une quantification attribue à
chaque arc $a$ une longueur $\ell_a \in \Z_{\ge 0}$ telle que les côtés
opposés de chaque patch aient la même longueur, ce sont les
\emph{équilibres} $B\ell = 0$\,; le maillage de quads est alors la grille
entière de chaque patch.

\begin{definition}[Tour]
Un \emph{tour} $\gamma$ est une marche en ligne droite à travers les patchs,
partant d'une abscisse de son premier côté qui ne rencontre aucun nœud~: elle
entre dans un patch par un côté, en sort par le côté opposé à la même distance
du coin correspondant, entre dans le patch voisin à travers l'arc qu'elle
atteint, et ainsi de suite, jusqu'à revenir dans son premier patch par son
premier côté. La rotation des passages composés est alors l'identité,
puisque la marche rentre dans son premier côté vers l'intérieur.
\end{definition}

Passer du patch $P$ au patch $Q$ à travers un arc compose une rotation d'un
quart de tour du repère local avec une translation dont les composantes sont
des sommes signées de longueurs d'arcs le long des bords de $P$ et de $Q$.

\begin{evidence}{core}{\lean{TMesh.Turn.shift_linear}, \lean{TMesh.Turn.mem_corridor_iff}, \lean{TMesh.Turn.comes_back}, \lean{TMesh.Turn.winds}\,; supposé~: un recollement orienté\,; seule l'abscisse transverse à la marche est suivie\,; le corridor sans demi-entier n'est pas formalisé}
\begin{proposition}[Le décalage d'un tour est linéaire]\label{prop:linear}
Pour un tour $\gamma$, il existe des coefficients entiers $c_a$, indépendants
des longueurs, tels qu'après un tour complet un point du premier côté s'est
déplacé transversalement à la direction de marche de
\[
  \shift_\gamma(\ell) = \sum_a c_a\, \ell_a .
\]
Pour des longueurs entières $\ell$, soit $o$ une abscisse demi-entière du
premier côté dont la marche aux longueurs $\ell$ traverse les mêmes patchs que
$\gamma$ (nous appelons l'ensemble de ces abscisses le \emph{couloir} de
$\gamma$ en $\ell$). Alors la ligne duale du maillage de quads en $o$ revient
en $o + \shift_\gamma(\ell)$~: elle se ferme après une révolution si et
seulement si $\shift_\gamma(\ell) = 0$. Si de plus $o + j\,\shift_\gamma(\ell)$
appartient au couloir pour $j = 1, \dots, t - 1$, elle fait $t$ fois le tour
en hélice de pas $|\shift_\gamma(\ell)|$ au sens
de~\cite{bommes2011structure}. Si le couloir ne contient aucun demi-entier,
aucune bande du maillage ne suit $\gamma$.
\end{proposition}
\end{evidence}

\begin{proof}[Esquisse de démonstration]
Chaque passage est une application $x \mapsto R_k x + t$, où $R_k$ est une
rotation de $k$ quarts de tour et $t$ un vecteur dont les composantes sont des
sommes de $\pm\ell_a$. La composée le long de $\gamma$ a une rotation
triviale, c'est donc une translation $T_\gamma(\ell)$, linéaire en $\ell$ à
coefficients entiers\,; sa composante transverse à la direction de marche est
$\shift_\gamma(\ell)$. Les singularités sont des nœuds, donc le maillage de
quads est une grille régulière le long de la marche, et une ligne duale à la
distance $o$ du premier coin qui traverse les mêmes patchs revient en $o +
\shift_\gamma(\ell)$. Le couloir, découpé par deux inégalités par passage,
peut contenir $o$ mais non $o + \shift_\gamma(\ell)$\,; la révolution
suivante suit alors d'autres patchs.
\end{proof}

Dans le repère du premier côté, avec $e_1$ la direction de marche et $e_2$
celle des abscisses croissantes, $T_\gamma(\ell) = -\alpha_\gamma(\ell)\, e_1 +
\shift_\gamma(\ell)\, e_2$~: il ramène la fin de la marche sur le premier
côté. $\alpha_\gamma(\ell)$ est la \emph{longueur} du tour, qui, aux longueurs entières, est le nombre de
quads de la bande par révolution. Aux longueurs réelles de la paramétrisation,
le décalage est le \emph{décalage réel}, et sa part de la longueur
l'\emph{inclinaison}. Aux longueurs entières équilibrées dont les côtés de
patch sont strictement positifs, toute bande du maillage de quads est une
telle marche~\cite[prop.~2]{gardent2026loops}, et tout décalage est
entier\,; fermer un anneau, c'est lui demander d'être nul.

\paragraph{Métrique plate et holonomie.}
Sauf mention contraire, la surface $\Sigma$ est fermée, et \emph{genre zéro}
signifie une sphère fermée. Des longueurs équilibrées $\ell > 0$ font de
chaque patch un rectangle euclidien, et de $\Sigma$ une surface plate à points
coniques aux sommets singuliers~: un sommet de valence $k$ a un angle
$k\pi/2$ et une courbure $(4-k)\pi/2$, et les jonctions en T ne sont pas des
points coniques. Soit $\Sigma^\circ$ la surface $\Sigma$ privée de ses
sommets singuliers, $x_0$ un point intérieur à un patch, et $R$ le quart de
tour positif. Développer patch après patch un lacet $\gamma$ de $\Sigma^\circ$
basé en $x_0$ donne son \emph{holonomie}, un déplacement $\hat
h(\gamma)\colon x \mapsto R^j x + t$ des coordonnées du premier patch\,; elle
ne dépend que de la classe d'homotopie de $\gamma$, et $\hat h(\alpha\beta) =
\hat h(\alpha) \circ \hat h(\beta)$~\cite{izmestiev2013torus}. Un petit lacet
autour d'un sommet singulier $c$ de valence $k$, parcouru avec $c$ à sa
droite et atteint le long d'un chemin qui développe $c$ en $p$, a pour
holonomie $x \mapsto R^{4-k}(x - p) + p$. Lorsque la rotation de $\hat
h(\gamma)$ est l'identité, sa translation est le \emph{vecteur de Burgers} de
$\gamma$. Pour un tour $\gamma$, soit $\bar\gamma$ la courbe fermée qui suit
la marche puis, de retour dans son premier patch, revient à son départ le long
d'un segment parallèle au premier côté\,; $\hat h(\bar\gamma)$ est une
translation. Elle prolonge la carte le long du lacet, tandis que les passages
de la proposition~\ref{prop:linear} développent le dernier patch dans le
premier~: $\hat h(\bar\gamma)$ est l'inverse de $T_\gamma(\ell)$, la
translation de vecteur $\alpha_\gamma(\ell)\, e_1 - \shift_\gamma(\ell)\, e_2$,
et $\shift_\gamma(\ell) = -\langle e_2, \hat h(\bar\gamma) \rangle$. Tous les
décalages de cet article sont ceux de la proposition~\ref{prop:linear}, la
convention des décalages qu'exporte le moteur et des figures. Les décalages, les positions développées et les vecteurs de Burgers
sont des formes linéaires à coefficients entiers en $\ell$ tant que les
chemins traversent les mêmes patchs, si bien qu'une identité entre eux qui
vaut sur un ouvert de longueurs équilibrées strictement positives vaut pour
toutes les longueurs équilibrées\,; pour des longueurs entières, les nœuds se
développent sur $\Z^2$. Une forme $f$ est égale à $g$ \emph{modulo les
équilibres} si $f - g$ est une combinaison des lignes de $B$.

\paragraph{Données.}
Les énoncés mesurés portent sur le T-mesh du corps masculin de la
bibliothèque Essentials de Blender remaillé par openclare à $10\,590$ quads,
exporté avec tous ses tours (branche \texttt{rd-exact}, moteur antérieur au
commit \texttt{c8bf2be}) et étudié dans l'article
compagnon~\cite[\S7]{gardent2026loops}~: $261$ sommets singuliers ($133$ de
valence trois, $127$ de valence cinq, un de valence deux) parmi $1\,295$
nœuds, $2\,069$ arcs, $776$ patchs, $1\,552$ équilibres. Les tours sont
trouvés selon huit \emph{variantes}, la variante $m + 2j + 4i$ marchant depuis
les fractions $\{0,5\}$ ($m = 0$) ou $\{0,2 ; 0,5 ; 0,8\}$ ($m = 1$) de chaque
côté et \emph{retenant} (cherchant à fermer) un tour si son décalage réel
dépasse $0,05$, si son inclinaison est inférieure à $0,05$ ($j = 0$) ou à
$0,12$ ($j = 1$) et si sa longueur est inférieure à $2$ ($i = 0$) ou à $6$ ($i
= 1$) côtés du T-mesh\,; le moteur garde la variante~2. Les variantes
trouvent $1\,299$ tours et en retiennent $765$ (en sommant sur les
variantes). La recherche entière maintient de $28$ à $36$ arcs à une longueur
d'au moins un, les \emph{bornes de séparation}\,; $P_{\mathrm{eng}}$,
l'ensemble des longueurs équilibrées qui les respectent, n'est ni un
sous-ensemble ni un sur-ensemble des quantifications valides, qui demandent
seulement qu'une chaîne entre deux sommets singuliers garde une longueur
strictement positive. Les identités sont vérifiées en virgule flottante sauf
mention de l'exactitude\,; les programmes ont été résolus avec HiGHS
1.8.0~\cite{huangfu2018highs} via SciPy 1.16.3~\cite{virtanen2020scipy}, sur
un Intel Core i7-7700K (4~cœurs, 2017, 62\,Go) partagé avec d'autres tâches.

\section{Tours et singularités}\label{sec:holonomy}

\begin{evidence}{core}{\lean{gauss_bonnet_turn}, \lean{gauss_bonnet}\,; testé~: 3\,833 régions de strips simples sur 50 instances aléatoires, 35 sphères et 15 tores\,; supposé~: la courbe borde un disque (ou une couronne) de quads}
\begin{lemma}[Gauss--Bonnet pour les tours]\label{lem:gauss}
Soit $\gamma$ un tour dont la courbe $\bar\gamma$ est simple. Si $\bar\gamma$
borde un disque $D \subset \Sigma$, les sommets singuliers intérieurs à $D$
vérifient $\sum_{c \in D} (4 - k(c)) = 4$. Si les courbes de deux tours sont
disjointes et bordent une couronne $A$, alors $\sum_{c \in A} (4 - k(c)) =
0$.
\end{lemma}
\end{evidence}

\begin{proof}
Pour une région d'une surface plate à points coniques bordée de segments
droits, la formule de Gauss--Bonnet s'écrit $\sum_c \kappa_c + \sum_i
\theta_i = 2\pi\chi$, où $\kappa_c = (4 - k(c))\pi/2$ sont les courbures
intérieures et $\theta_i$ les angles de rotation du
bord~\cite{bommes2013survey}. La courbe $\bar\gamma$ ne tourne qu'aux
extrémités de son segment de fermeture, de $+\pi/2$ et de $-\pi/2$, si bien
que les angles de rotation ont une somme nulle sur chaque courbe, vue de l'un
ou l'autre côté\,; $\chi(D) = 1$ et $\chi(A) = 0$.
\end{proof}

Sur une sphère fermée, un tour simple sépare donc les sommets singuliers en
deux groupes d'indice total quatre, chacun d'au moins deux sommets~: un bout
de doigt au-delà d'un anneau contient l'équivalent de quatre sommets de
valence trois.

\begin{evidence}{core}{\lean{one_crossing_on_walk}, \lean{one_crossing_on_segment}\,; testé~: 42 points doubles sur les sphères de 50 instances aléatoires\,; supposé~: les régions coupées au point double sont des disques, et lesquels de leurs coins sont convexes}
\begin{lemma}[Un seul point double]\label{lem:eight}
Supposons la surface fermée de genre zéro, et que la courbe $\bar\gamma$ d'un
tour se croise elle-même exactement une fois, transversalement, sans sommet
singulier sur elle. Elle découpe la surface en trois disques~: deux
\emph{lobes}, chacun bordé par l'un des deux lacets de $\bar\gamma$, et une
troisième région bordée par les deux. Si le croisement est sur la marche, les
totaux de $4 - k$ sur les sommets singuliers des lobes et de la troisième
région sont $3$, $3$ et $2$\,; s'il est sur le segment de fermeture, ils sont
$2$, $2$ et $4$.
\end{lemma}
\end{evidence}

\begin{proof}
La marche se croise elle-même à angle droit, si bien que chaque lacet quitte
le croisement $X$ par l'un des passages et y revient par l'autre, ses deux
rayons étant à un quart de tour l'un de l'autre. Chaque lacet est simple et
borde un disque qui contient le quadrant compris entre ses rayons, un lobe qui
a en $X$ un coin d'angle $\pi/2$\,; les lobes occupent des quadrants opposés
et la troisième région les deux autres, et les trois faces de ce graphe à un
sommet sur la sphère sont des disques. Par Gauss--Bonnet, un disque dont le
bord tourne de $\theta_i$ contient un indice total $4 - \sum_i
\theta_i/(\pi/2)$, ce qui donne $3$, $3$ et $2$ avant les coins du segment de
fermeture, de rotations $+\pi/2$ et $-\pi/2$ successivement. Si le croisement
est sur la marche, ces deux coins sont sur un même lacet, chaque région voit
les deux, et ils se compensent. Sinon, supposons que la marche parte de $A$
selon $e_1$, revienne en $Z$ et croise le segment de fermeture entre eux\,;
elle le croise selon $-e_1$, puisque selon $e_1$ le tour se serait terminé là.
Le lacet qui passe par $A$ tourne alors dans le même sens en $A$ qu'en $X$, de
même que le lacet qui passe par $Z$, en $Z$ et en $X$ (on le vérifie sur les
quatre directions)~: chaque lobe voit un second coin convexe et contient $2$,
et la troisième région le reste du total $8$.
\end{proof}

\begin{evidence}{core}{\lean{Motion.burgers_eq_sum}, \lean{Motion.composeAll_rotations}, \lean{Forest.glue_eq_preorder}, \lean{Motion.annulus_holonomy}\,; testé~: (c) sur 898 couronnes sans point conique de 50 instances aléatoires\,; supposé~: l'holonomie de la courbe est la composée des rotations autour des points coniques}
\begin{proposition}[Le décalage est fixé par les singularités]\label{prop:burgers}
Soient $\gamma$, $\gamma_1$, $\gamma_2$ des tours à courbes simples.
\begin{enumerate}
\item[(a)] Supposons que $\bar\gamma$ borde un disque $D$ contenant les
  sommets singuliers $c_1, \dots, c_m$, et soit $x_0$ un point de
  $\bar\gamma$, par exemple le départ de la marche. Soit $T \subset D$ un
  arbre dont les sommets comprennent $c_1, \dots, c_m$ et qui ne rencontre
  $\bar\gamma$ qu'en $x_0$. Découpé le long de $T$, le disque $D$ devient un
  disque $\Omega$ sans point conique, qui se développe dans le plan, avec le
  repère $(e_1, e_2)$ en $x_0$. Numérotons les $c_i$ dans l'ordre où un
  parcours de $T$ partant de $x_0$ et gardant $T$ à sa droite les atteint pour
  la première fois, et soit $p_i$ la position développée de $c_i$ à cette
  première visite. Alors $\shift_\gamma(\ell) = \langle e_2, B_D \rangle$ si
  $D$ est à gauche de $\bar\gamma$, et $\shift_\gamma(\ell) = -\langle e_2,
  B_D \rangle$ s'il est à sa droite, avec
  \[
    B_D = \sum_{i=1}^m Q_i\,(I - R_i)\,p_i ,
    \qquad R_i = R^{4 - k(c_i)},
    \qquad Q_i = R_m R_{m-1} \cdots R_{i+1},
  \]
  la translation de $\rho_m \circ \dots \circ \rho_1$, où $\rho_i(x) =
  R_i(x - p_i) + p_i$. Avec $D$ à sa gauche, $B_D = T_\gamma(\ell) =
  -\alpha_\gamma(\ell)\, e_1 + \shift_\gamma(\ell)\, e_2$\,; avec $D$ à sa
  droite, $B_D$ est la translation inverse.
\item[(b)] Supposons que $\bar\gamma_1$ et $\bar\gamma_2$ soient disjointes et
  bordent une couronne $A$, dont les deux tours font le tour dans le même
  sens, et qu'un chemin $\kappa$ intérieur à $A$, de $\bar\gamma_1$ à
  $\bar\gamma_2$, transporte la direction de marche de $\gamma_2$ sur celle de
  $\gamma_1$ (les tours sont \emph{parallèles})\,; mesurons les deux
  décalages selon la direction $e_2$ transportée par $\kappa$. Alors
  \(
    \shift_{\gamma_1}(\ell) - \shift_{\gamma_2}(\ell) = -\langle e_2, B_A \rangle,
  \)
  où $B_A$ est le vecteur de Burgers du bord du disque obtenu en découpant
  $A$ le long de $\kappa$, donné par la même somme sur les sommets singuliers
  intérieurs à $A$, développés dans le disque découpé.
\item[(c)] En particulier, des tours parallèles qui bordent une couronne sans
  sommet singulier ont le même décalage pour tout $\ell$~: une hélice est une
  propriété de tout un cylindre de tours parallèles, non de l'un d'eux.
\end{enumerate}
\end{proposition}
\end{evidence}

\begin{figure}[t]
\centering
% The shift of a turn is set by the singular vertices on one side of it
% French copy of holonomy.tex: only the labels differ.
% (Proposition 2, "The shift is set by the singularities").
% \input inside a figure; needs \usepackage{tikz} and
% \usetikzlibrary{arrows.meta,decorations.pathreplacing}.
%
% A cap D of a T-mesh (a fingertip, say), cut along a tree T into a disc
% Omega that develops into the plane; every node is at a whole position and the
% faint grid is the quad mesh at these lengths. Omega is the band [0,11]x[.,4]
% (it goes on below), the block [3,11]x[4,5] and the top, the L-shaped
% [0,3]x[4,6] + [0,2]x[6,7]. T runs from x_0 up the left side to c_1 and then
% round the border of the top: c_1, c_2, ..., c_6, which are the six singular
% vertices of D. The cut sides are glued in pairs: the right side x = 11 to the
% left side one unit lower, (11,y) ~ (0,y-1); and the top of the block, from
% c_6 = (3,5) rightwards, onto the border of the top, turning round it:
% [3,4] onto [(3,5),(3,6)], [4,5] onto [(3,6),(2,6)], [5,6] onto [(2,6),(2,7)],
% [6,8] onto [(2,7),(0,7)], [8,11] onto [(0,7),(0,4)]. So c_1 = (11,5) = (0,4),
% c_2 = (8,5) = (0,7), c_3 = (6,5) = (2,7), c_4 = (5,5) = (2,6), c_5 = (4,5) =
% (3,6), c_6 = (3,5). Angles: 90 + 90 + 90 at c_1; 90 + 180 at c_2, c_3, c_5;
% 180 + 90 at c_6; 270 + 180 at c_4. Valences 3, 3, 3, 5, 3, 3; total index
% 5 - 1 = 4 (Lemma 1, a disc).
%
% Traces besides the cut: c_1 to the right (stops at (2,4)), c_4 to the left
% (stops on the cut at (0,6) = (9,5)) and downwards (x = 2), c_6 to the left
% (stops at (2,5)) and downwards (x = 3), and the separatrices downwards from
% c_5, c_4, c_3, c_2 at x = 4, 5, 6, 8.
%
% The turn gamma from x_0 = (0, 3.5) crosses x = 2, 3, 4, 5, 6, 8, 11 and comes
% back at (0, 2.5): length 11, s = -1 (one row lower; offsets measured
% upwards). The closing segment of its curve runs up the first side from 2.5
% to 3.5; continued along the walk, the chart brings x_0 back at x_0' =
% (11, 4.5), where the walk round T (keeping T on its right) starts.
% First visits on that walk, with R the quarter-turn x -> (-y, x), R_i = R for
% valence 3 and R^{-1} for valence 5, Q_i = R_6 ... R_{i+1}:
%   p_1 = (11,5) v3  (I-R)p = (16,-6)  Q_1 = R^3  -> ( -6,-16)
%   p_2 = ( 8,5) v3  (I-R)p = (13,-3)  Q_2 = R^2  -> (-13,  3)
%   p_3 = ( 6,5) v3  (I-R)p = (11,-1)  Q_3 = R    -> (  1, 11)
%   p_4 = ( 5,5) v5  (I-R^-1)p = (0,10) Q_4 = R^2 -> (  0,-10)
%   p_5 = ( 4,5) v3  (I-R)p = ( 9, 1)  Q_5 = R    -> ( -1,  9)
%   p_6 = ( 3,5) v3  (I-R)p = ( 8, 2)  Q_6 = I    -> (  8,  2)
%   B_D = sum = (-11, -1) = x_0 - x_0'.
% Its along component is the length 11, its across component the shift -1:
% |s_gamma| = |<e_2, B_D>| (Proposition 2(a)). The other order of visit gives
% (1,1), which is not a translation of x_0: the order matters.
% gamma' from (0, 1.5) comes back at (0, 0.5): s = -1, length 11. The curves of
% gamma and gamma' bound an annulus without singular vertex (shaded), and the
% shifts agree (Proposition 2(c)).
% Parity on a sphere (Proposition 5): 11 = -1 mod 2.
% check.py builds this surface and checks these numbers by walking the strips
% of the quad mesh: the valences above, every other vertex of valence 4,
% s = -1 and length 11 from o = 1.5, 2.5, 3.5, and the sum B_D.
\begin{tikzpicture}[x=0.55cm, y=0.55cm, font=\footnotesize,
    side/.style={draw=black, line width=0.5pt},
    quadgrid/.style={draw=black!18, line width=0.25pt, step={(1,1)}},
    surface/.style={fill=black!3},
    node dot/.style={circle, fill=black, inner sep=0pt, minimum size=2.6pt},
    turn/.style={line width=1pt, -{Stealth[length=5pt]}},
    closing/.style={line width=0.7pt, densely dotted},
    cut/.style={draw=black, line width=1.1pt, line join=round},
    brace/.style={decorate, decoration={brace, amplitude=3pt}, line width=0.5pt},
    burgers/.style={draw=measured, line width=1.1pt, -{Stealth[length=6pt]}},
    val3/.pic={\fill[valthree, draw=black, line width=0.3pt]
      (90:3.4pt) -- (210:3.4pt) -- (330:3.4pt) -- cycle;},
    val5/.pic={\fill[valfive, draw=black, line width=0.3pt]
      (90:3.4pt) -- (162:3.4pt) -- (234:3.4pt) -- (306:3.4pt) -- (18:3.4pt) -- cycle;}]
  \definecolor{loopspiral}{HTML}{D55E00}
  \definecolor{annulus}{HTML}{CFE0F3}
  \definecolor{valthree}{HTML}{E8890C}
  \definecolor{valfive}{HTML}{7C3AAD}
  \definecolor{measured}{HTML}{14934A}

  % The disc Omega: band, block, top.
  \fill[surface] (0,-0.6) rectangle (11,4) (3,4) rectangle (11,5)
    (0,4) rectangle (3,6) (0,6) rectangle (2,7);
  % The annulus between the curves of gamma' and gamma.
  \fill[annulus] (0,0.5) -- (0.15,0.5) -- (0.15,1.5) -- (11,1.5) -- (11,3.5)
    -- (0.15,3.5) -- (0.15,2.5) -- (0,2.5) -- cycle;
  \draw[quadgrid] (0,-0.6) grid (11,4);
  \draw[quadgrid] (3,4) grid (11,5);
  \draw[quadgrid] (0,4) grid (3,6);
  \draw[quadgrid] (0,6) grid (2,7);
  % Traces. The cut T: the left side up to c_1 and the border of the top; its
  % other side: the right side and the top of the block.
  \draw[side] (0,-0.6) -- (0,3.5) (11,-0.6) -- (11,4.5);
  \draw[cut] (0,3.5) -- (0,7) -- (2,7) -- (2,6) -- (3,6) -- (3,5) -- (11,5) -- (11,4.5);
  \draw[side] (0,4) -- (2,4) (2,6) -- (2,-0.6) (3,5) -- (2,5) (3,5) -- (3,-0.6) (2,6) -- (0,6);
  \foreach \x in {4,5,6,8} \draw[side] (\x,5) -- (\x,-0.6);
  \foreach \p in {(2,4), (2,5), (0,6), (9,5)} \node[node dot] at \p {};
  % The turns, their closing segments and where they come back.
  \foreach \o/\name in {3.5/{\gamma}, 1.5/{\gamma'}} {
    \draw[turn, loopspiral] (0,\o) -- (11.1,\o);
    \draw[closing, loopspiral] (0.15,\o-1) -- (0.15,\o);
    \draw[loopspiral, line width=0.5pt, fill=white] (0,\o-1) circle (1.8pt);
    \draw[loopspiral, line width=0.5pt, fill=white] (0,\o) circle (1.8pt);
    \node[above, loopspiral] at (9.6,\o) {$\name$};
    \node[right, loopspiral] at (11.2,\o) {$s = -1$};
  }
  \draw[loopspiral, line width=0.5pt, fill=white] (11,4.5) circle (1.8pt);
  \node[left=1pt] at (0,3.5) {$x_0$};
  \node[right=1pt] at (11,4.5) {$x_0'$};
  % The singular vertices: first visits p_i on the top of the block, and
  % the other copies on the border of the top.
  \pic at (3,5) {val3};
  \foreach \p in {(0,4), (11,5), (0,7), (8,5), (2,7), (6,5), (3,6), (4,5)} \pic at \p {val3};
  \foreach \p in {(2,6), (5,5)} \pic at \p {val5};
  \begin{scope}[font=\scriptsize, inner sep=1pt]
    \foreach \i/\x in {1/11, 2/8, 3/6, 4/5, 5/4} \node[above=2pt] at (\x,5) {$p_{\i}$};
    \node[above right=1pt] at (3,5) {$p_6$};
    \node[left=2pt] at (0,4) {$c_1$};
    \node[left=2pt] at (0,7) {$c_2$};
    \node[above=2pt] at (2,7) {$c_3$};
    \node[above left=0pt] at (2,6) {$c_4$};
    \node[right=2pt] at (3,6) {$c_5$};
  \end{scope}
  \draw[brace] (-1.5,3.5) -- (-1.5,7) node[midway, left=3pt] {$D$};
  \node[above] at (1,7) {$T$};
  % The Burgers vector of D, from x_0' to x_0, and its two components.
  \draw[burgers] (11,-1.6) -- (0,-2.6);
  \draw[measured, dashed, line width=0.5pt] (11,-1.6) -- (0,-1.6) -- (0,-2.6);
  \node[measured, above] at (5.5,-1.6) {en long~: $\alpha_\gamma = 11$};
  \node[measured, left] at (0,-2.1) {en travers~: $s_\gamma = -1$};
  \node[measured, below] at (6.5,-2.3) {$B_D = \sum_i Q_i (I - R_i)\, p_i = (-11, -1)$};
\end{tikzpicture}

\caption{Le décalage lu sur les singularités
(proposition~\ref{prop:burgers}). Une calotte $D$ à six sommets singuliers
(valences $3, 3, 3, 5, 3, 3$, indice total $4$), découpée le long d'un arbre
$T$ et développée\,; chaque $c_i$ est dessiné à sa première visite $p_i$ lors
du parcours de $T$. Le tour $\gamma$, de longueur $11$, revient une rangée
plus bas, et $D$ est à sa gauche~: $B_D = \sum_i Q_i (I - R_i) p_i = (-11, -1)
= T_\gamma(\ell)$, moins la longueur le long de la marche et le décalage $-1$
en travers. Le tour $\gamma'$
borde avec $\gamma$ une couronne sans sommet singulier (en grisé) et a le même
décalage (partie~(c))\,; $11 \equiv -1 \pmod 2$
(proposition~\ref{prop:parity}). Dans toutes les figures, les abscisses sont
mesurées vers le haut.}
\label{fig:holonomy}
\end{figure}

L'ordre compte~: sur la figure~\ref{fig:holonomy}, parcourir $T$ dans l'autre
sens donne $(1, 1)$ au lieu de $B_D = (-11, -1)$.

\begin{proof}
(a) Orientons chaque arête $f$ de $T$ en l'éloignant de $x_0$\,; ses deux
côtés sur le bord de $\Omega$ sont recollés par un déplacement $G_f$, du côté
que le parcours de $T$ suit en premier vers l'autre. Supposons que $f$
aboutisse en $v$, d'où partent $f_1, \dots, f_r$ dans l'ordre du parcours. Une
petite courbe qui tourne dans le sens des aiguilles d'une montre autour de
$v$, du premier côté de $f$ au second, traverse $f_1, \dots, f_r$ et parcourt
tout l'angle en $v$, donc $G_f = G_{f_r} \circ \dots \circ G_{f_1} \circ
\rho_v$, avec $\rho_v(x) = R^{4-k(v)}(x - p_v) + p_v$, $p_v$ la position à la
première visite ($\rho_v = I$ si $v$ n'est pas singulier). En $x_0$,
$\bar\gamma$ étant parcourue avec $D$ à sa droite à la place de la petite
courbe, $\hat h(\bar\gamma) = G_{f_r} \circ \dots \circ G_{f_1}$. Dérouler la
récurrence compose les $\rho_{c_i}$ dans l'ordre de première visite\,; la
rotation est $R^{\sum_i(4 - k(c_i))} = I$ d'après le lemme~\ref{lem:gauss},
le développement donne $B_D = \hat h(\bar\gamma) = T_\gamma(\ell)^{-1}$, et
avec $D$ à sa gauche $\bar\gamma$ a l'holonomie inverse, $B_D =
T_\gamma(\ell)$. (b) Avec $x_0$ sur $\bar\gamma_1$ au départ de $\kappa$,
le lacet $\lambda = \bar\gamma_1\, \kappa\, \bar\gamma_2^{-1} \kappa^{-1}$
borde le disque obtenu en découpant $A$ le long de $\kappa$. Développer le
long de $\kappa$ lit $\hat h(\bar\gamma_2)$ dans les coordonnées en $x_0$, où
la marche de $\gamma_2$ suit $e_1$\,; ainsi $\hat h(\lambda)$ est la
translation de vecteur $\hat h(\bar\gamma_1) - \hat h(\bar\gamma_2)$, dont
la composante selon $e_2$ est $\shift_{\gamma_2}(\ell) - \shift_{\gamma_1}(\ell)$,
et la partie~(a) s'applique au disque découpé, dont les sommets singuliers ont un
indice total nul. (c) La somme est vide.
\end{proof}

Dans les T-meshes d'openclare, dont les traces partent des sommets
singuliers, une couronne sans sommet singulier ne contient aucun nœud, à moins
qu'une trace qui fait le tour n'y entre par un segment de fermeture, si bien
que deux tours comme en (c) sont en général dans un même couloir.

La configuration la plus simple entre deux tours parallèles est une paire de
valences trois et cinq, qui permet au nombre d'arêtes autour d'un membre de
changer. Une paire de disclinaisons opposées de rotation $\Omega$ est une
dislocation dont le vecteur de Burgers vaut $2\sin(\Omega/2)$ fois leur
écartement~\cite[éq.~(5)]{kleman2008disclinations}. Nous la lisons sur les
décalages de deux tours (figures~\ref{fig:holonomy}
et~\ref{fig:dislocation}).

\begin{evidence}{core}{\lean{dislocation_burgers}, \lean{dislocation}, \lean{dislocation_closes_iff}\,; testé~: (i) sur 5 couronnes contenant une paire 3--5 de 50 instances aléatoires\,; supposé~: comme dans la proposition ci-dessus}
\begin{proposition}[Une dislocation entre deux tours]\label{prop:dislocation}
Dans le cadre de la proposition~\ref{prop:burgers}(b), supposons que la
couronne $A$ contienne exactement deux sommets singuliers, $c_3$ de valence
trois et $c_5$ de valence cinq, et soit $d = p_5 - p_3 = d_1 e_1 + d_2 e_2$
leur écartement, développé dans le disque découpé. Alors~:
\begin{enumerate}
\item[(i)] $\shift_{\gamma_1}(\ell) - \shift_{\gamma_2}(\ell) = \varepsilon_1
  d_1 + \varepsilon_2 d_2$ avec des signes $\varepsilon_1, \varepsilon_2 \in
  \{+1, -1\}$ qui dépendent de la configuration mais pas de $\ell$, et $|B_A|
  = \sqrt2\,|d|$. Les deux tours ne peuvent se fermer ensemble que si $c_5$
  se trouve, vu de $c_3$, sur une diagonale particulière du champ.
\item[(ii)] Si $c_3$ et $c_5$ sont sur un même segment de droite du champ
  intérieur à $A$, de longueur $\delta$ (dans un T-mesh~: ils sont reliés par
  une chaîne d'arcs le long d'une même trace, et $\delta$ est la somme de
  leurs longueurs), et si un chemin $\kappa$ comme dans la
  proposition~\ref{prop:burgers}(b) ne traverse pas ce segment, alors
  $|\shift_{\gamma_1}(\ell) - \shift_{\gamma_2}(\ell)| = \delta$. Les deux
  tours ne peuvent être tous deux des anneaux que si $\delta = 0$,
  c'est-à-dire si les deux singularités fusionnent\,; tant que les bornes de
  séparation maintiennent $\delta \ge 1$, l'un d'eux au moins est une hélice.
\item[(iii)] Pour des longueurs entières, $\shift_{\gamma_1}(\ell) -
  \shift_{\gamma_2}(\ell) \equiv d_1 + d_2 \pmod 2$~: une paire dont les deux
  sommets sont à un nombre impair d'arêtes l'un de l'autre, comptées selon les
  deux directions du champ, ne se trouve jamais entre deux anneaux.
\end{enumerate}
\end{proposition}
\end{evidence}

\begin{proof}
Avec $R_3 = R$ et $R_5 = R^{-1}$, les deux ordres donnent $B_A = R^{-1}(I -
R)p_3 + (I - R^{-1})p_5 = (I - R^{-1})d$ ou $B_A = R(I - R^{-1})p_5 + (I -
R)p_3 = (I - R)(-d)$. Avec $R(a, b) = (-b, a)$, on a $(I - R^{-1})(d_1, d_2) =
(d_1 - d_2,\ d_1 + d_2)$ et $(I - R)(-d_1, -d_2) = (-d_1 - d_2,\ d_1 -
d_2)$\,; inverser l'orientation du bord inverse les rotations et donne les
deux autres combinaisons de signes (vérifié symboliquement). Dans les quatre
cas, les deux composantes sont $\pm d_1 \pm d_2$ et $|B_A|^2 = 2|d|^2$, d'où
(i). Pour (ii), prenons $\kappa$ ne traversant pas le segment qui joint $c_3$
à $c_5$ (le traverser ferait tourner le repère d'un quart de tour autour de
$c_3$) et un arbre contenant ce segment, et développons $c_5$ depuis $c_3$ le
long de celui-ci~: $d$ vaut $\delta$ fois l'un des vecteurs $\pm e_1$, $\pm
e_2$. Pour (iii), $d \in \Z^2$ et $\pm d_1 \pm d_2 \equiv d_1 + d_2$.
\end{proof}

\begin{figure}[t]
\centering
% A 3-5 dislocation between two parallel turns (Propositions 3 and 6).
% French copy of dislocation.tex: only the labels differ.
% \input inside a figure; needs \usepackage{tikz} and
% \usetikzlibrary{arrows.meta,decorations.pathreplacing}.
%
% Each panel is a band of a T-mesh on a cylinder, cut open along a trace of
% the vertex c_5 of valence five: the right side is glued to the left side, the
% lower part (below c_5) at the same height, the upper part two units higher.
% Every node is at a whole position and the faint grid is the quad mesh at
% these lengths. c_5 is drawn three times (its angle 5 x 90 degrees is split
% between the three copies: 90 at the bottom left, 180 at the right, 180 at
% the left), c_3 once (270 degrees). The empty corner next to c_3 is not part of
% the surface: its two thick sides are glued to each other by a quarter-turn
% about c_3. The turn gamma_1 (below the pair) and gamma_2 (above it) bound an
% annulus that holds exactly c_3 and c_5 (Proposition 2(b)).
%
% (a) Straight chain (Proposition 3(ii)). c_3 = (2,3), c_5 = (0,3) = (10,3) =
% (2,5); the gluings are (10,y) ~ (0,y) for y < 3, (10,y) ~ (2,y+2) for y > 3,
% and the quarter-turn about c_3 taking [c_3, (0,3)] onto [c_3, (2,5)]. c_3 and
% c_5 are joined along the trace y = 3 by the chain of arcs b = [c_3, (1,3)]
% and c = [(1,3), c_5]; a trace from below (its vertex is out of the picture)
% stops at (1,3). delta = l_b + l_c = 1 + 1 = 2.
%   gamma_1 at o = 1.5: crosses x = 1, 2, 10; comes back at 1.5: s = 0,
%   length 10 (a ring).
%   gamma_2 at o = 5.5 on the side x = 2: crosses x = 10 at height 5.5, which
%   is height 7.5 on the side x = 2: s = 2, length 10 - 2 = 8 (a helix).
%   |s_2 - s_1| = 2 = delta (Proposition 3(ii)); the lengths differ by 2 too:
%   the Burgers vector of the pair is (2, 2), of norm sqrt(2) |d| with d = (2, 0)
%   from c_3 to c_5 (Proposition 3(i)); s_2 - s_1 = 2 = d_1 + d_2 mod 2 (iii).
%   The separatrix y = 5 of c_5 goes once round and stops on the trace x = 10
%   two units above c_5 (Proposition 8 gives s_2 = 2 again).
%   Both turns can be rings only if delta = 0: c_3 and c_5 merge.
%   Parity on a sphere (Proposition 5): 10 = 0 and 8 = 2 mod 2.
%
% (b) Staircase (Proposition 6). c_3 = (1,4), c_5 = (0,3) = (10,3) = (0,5);
% gluings (10,y) ~ (0,y) for y < 3, (10,y) ~ (0,y+2) for y > 3, and the
% quarter-turn about c_3 taking [c_3, (1,3)] onto [c_3, (0,4)] and
% [(1,3), (0,3)] onto [(0,4), (0,5)]. c_3 and c_5 are on no common trace; the
% staircase from c_3 to c_5 is b = [c_3, (1,3)], down, then c = [(1,3), c_5],
% left: d = p_5 - p_3 = (-1, -1), and with epsilon = (-1, -1) each step has
% epsilon . v = 1 (b: v = (0,-1); c: v = (-1,0)).
%   gamma_1 at o = 1.5: s = 0, length 10. gamma_2 at o = 5.5: crosses x = 10
%   at height 5.5 = height 7.5 on the side x = 0: s = 2, length 10.
%   s_2 - s_1 = 2 = l_b + l_c = epsilon . d (Proposition 6); here the lengths
%   are equal, the Burgers vector of the pair is (0, 2). Both are rings only
%   if l_b = l_c = 0. Parity: 10 = 0 = 2 mod 2.
% check.py builds both bands and checks these numbers by walking the strips
% of the quad mesh: valences 5 and 3 at the copies listed above, every other
% vertex of valence 4, and the turns above.
\begin{tikzpicture}[x=0.5cm, y=0.5cm, font=\footnotesize,
    side/.style={draw=black, line width=0.5pt},
    quadgrid/.style={draw=black!18, line width=0.25pt, step={(1,1)}},
    surface/.style={fill=black!3},
    node dot/.style={circle, fill=black, inner sep=0pt, minimum size=2.6pt},
    arc mark/.style={line width=1.8pt, draw=black},
    turn/.style={line width=1pt, -{Stealth[length=5pt]}},
    brace/.style={decorate, decoration={brace, amplitude=3pt}, line width=0.5pt},
    val3/.pic={\fill[valthree, draw=black, line width=0.3pt]
      (90:3.4pt) -- (210:3.4pt) -- (330:3.4pt) -- cycle;},
    val5/.pic={\fill[valfive, draw=black, line width=0.3pt]
      (90:3.4pt) -- (162:3.4pt) -- (234:3.4pt) -- (306:3.4pt) -- (18:3.4pt) -- cycle;}]
  \definecolor{loopspiral}{HTML}{D55E00}
  \definecolor{loopclosed}{HTML}{0B3C8C}
  \definecolor{valthree}{HTML}{E8890C}
  \definecolor{valfive}{HTML}{7C3AAD}

  % ---------------------------------------------------------------- (a)
  \begin{scope}
    % The surface: the lower band [0,10]x[0,3] and the upper part [2,10]x[3,8].
    \fill[surface] (0,-0.6) rectangle (10,3) (2,3) rectangle (10,8.6);
    \draw[quadgrid] (0,-0.6) grid (10,3);
    \draw[quadgrid] (2,3) grid (10,8.6);
    % Traces: the cut (a trace of c_5, both copies), the trace from c_3 round
    % to c_5, c_3 downwards, a trace from below, the winding separatrix.
    \draw[side] (0,-0.6) -- (0,3) (10,-0.6) -- (10,8.6) (2,5) -- (2,8.6);
    \draw[side] (2,3) -- (10,3) (2,-0.6) -- (2,3) (1,-0.6) -- (1,3) (2,5) -- (10,5);
    % The chain from c_3 to c_5, and its glued copy.
    \draw[arc mark] (0,3) -- (2,3) (2,3) -- (2,5);
    \node[below] at (1.5,3) {$b$};
    \node[below] at (0.5,3) {$c$};
    \foreach \p in {(1,3), (2,4), (10,5), (2,7)} \node[node dot] at \p {};
    % The turns.
    \draw[turn, loopclosed] (0,1.5) -- (10.1,1.5);
    \draw[loopclosed, line width=0.5pt, fill=white] (0,1.5) circle (1.8pt);
    \node[above, loopclosed] at (6,1.5) {$\gamma_1$};
    \draw[turn, loopspiral] (2,5.5) -- (10.1,5.5);
    \draw[loopspiral, line width=0.5pt, fill=white] (2,5.5) circle (1.8pt);
    \draw[loopspiral, dashed, line width=0.8pt, -{Stealth[length=4pt]}] (2,7.5) -- (4,7.5);
    \draw[loopspiral, line width=0.5pt, fill=white] (2,7.5) circle (1.8pt);
    \node[above, loopspiral] at (6,5.5) {$\gamma_2$};
    \draw[brace] (1.7,5.5) -- (1.7,7.5) node[midway, left=3pt] {$2$};
    % The singular vertices.
    \pic at (2,3) {val3};
    \foreach \p in {(0,3), (10,3), (2,5)} \pic at \p {val5};
    \node[above right=-1pt] at (2,3) {$c_3$};
    \node[left=1pt] at (0,3) {$c_5$};
    % Right: shifts and lengths.
    \node[right, align=left, loopspiral] at (10.3,5.5) {$s = 2$\\ 8 quads};
    \node[right, align=left, loopclosed] at (10.3,1.5) {$s = 0$\\ 10 quads};
    \node[anchor=north, align=center] at (6,-1.1)
      {$s_{\gamma_2} - s_{\gamma_1} = \ell_b + \ell_c = \delta = 2$};
    \node[anchor=north] at (6,-2.4) {\small (a) chaîne droite};
  \end{scope}

  % ---------------------------------------------------------------- (b)
  \begin{scope}[shift={(16.8,0)}]
    % The surface: [0,10]x[0,3], [1,10]x[3,4], [0,10]x[4,8].
    \fill[surface] (0,-0.6) rectangle (10,3) (1,3) rectangle (10,4) (0,4) rectangle (10,8.6);
    \draw[quadgrid] (0,-0.6) grid (10,3);
    \draw[quadgrid] (1,3) grid (10,4);
    \draw[quadgrid] (0,4) grid (10,8.6);
    % Traces: the cut (a trace of c_5), the loop of c_5 at y = 3, the three
    % separatrices of c_3, the winding separatrix of c_5.
    \draw[side] (0,-0.6) -- (0,3) (10,-0.6) -- (10,8.6) (0,5) -- (0,8.6);
    \draw[side] (0,3) -- (10,3) (1,4) -- (10,4) (1,4) -- (1,5) (0,5) -- (10,5);
    % The staircase from c_3 to c_5, and its glued copy.
    \draw[arc mark] (1,4) -- (1,3) -- (0,3) (1,4) -- (0,4) -- (0,5);
    \node[right] at (1,3.5) {$b$};
    \node[below] at (0.5,3) {$c$};
    \foreach \p in {(1,3), (0,4), (10,4), (0,6), (1,5), (10,5), (0,7)}
      \node[node dot] at \p {};
    % The turns.
    \draw[turn, loopclosed] (0,1.5) -- (10.1,1.5);
    \draw[loopclosed, line width=0.5pt, fill=white] (0,1.5) circle (1.8pt);
    \node[above, loopclosed] at (6,1.5) {$\gamma_1$};
    \draw[turn, loopspiral] (0,5.5) -- (10.1,5.5);
    \draw[loopspiral, line width=0.5pt, fill=white] (0,5.5) circle (1.8pt);
    \draw[loopspiral, dashed, line width=0.8pt, -{Stealth[length=4pt]}] (0,7.5) -- (2,7.5);
    \draw[loopspiral, line width=0.5pt, fill=white] (0,7.5) circle (1.8pt);
    \node[above, loopspiral] at (6,5.5) {$\gamma_2$};
    \draw[brace] (-0.3,5.5) -- (-0.3,7.5) node[midway, left=3pt] {$2$};
    % The singular vertices.
    \pic at (1,4) {val3};
    \foreach \p in {(0,3), (10,3), (0,5)} \pic at \p {val5};
    \node[above right=-1pt] at (1,4) {$c_3$};
    \node[left=1pt] at (0,3) {$c_5$};
    % Right: shifts and lengths.
    \node[right, align=left, loopspiral] at (10.3,5.5) {$s = 2$\\ 10 quads};
    \node[right, align=left, loopclosed] at (10.3,1.5) {$s = 0$\\ 10 quads};
    \node[anchor=north, align=center] at (6,-1.1)
      {$s_{\gamma_2} - s_{\gamma_1} = \ell_b + \ell_c = 2$};
    \node[anchor=north] at (6,-2.4) {\small (b) escalier};
  \end{scope}
\end{tikzpicture}

\caption{Une dislocation 3--5 entre un anneau $\gamma_1$ et une hélice
$\gamma_2$ d'une bande, ouverte le long d'une trace du sommet de valence cinq,
qui est dessiné trois fois. (a)~La paire est reliée le long d'une même trace
par deux arcs, $\delta = \ell_b + \ell_c = 2$~: $\gamma_2$ revient $2$ rangées
plus haut, avec $8$ quads par révolution, et le vecteur de Burgers de la paire
est $(2, 2)$ (proposition~\ref{prop:dislocation}). (b)~La paire est reliée
par un escalier de deux arcs, $d = \varepsilon = (-1, -1)$~: $\gamma_2$
revient $2$ rangées plus haut, avec $10$ quads
(proposition~\ref{prop:staircase}). Dans les deux cas, les décalages
diffèrent de $\ell_b + \ell_c = 2$, et les deux tours ne sont des anneaux que
si les deux sommets fusionnent.}
\label{fig:dislocation}
\end{figure}

\begin{evidence}{core}{\lean{checkerboard}, \lean{oddColours_translate}\,; testé~: 1\,960 couronnes sur les sphères de 50 instances aléatoires\,; supposé~: la somme de Burgers de la proposition ci-dessus}
\begin{proposition}[Règle du damier]\label{prop:checkerboard}
Dans le cadre de la proposition~\ref{prop:burgers}(b), pour des longueurs
entières, colorions en noir ou en blanc les points $p \in \Z^2$ du disque
découpé selon la parité de $p_1 + p_2$ (l'une ou l'autre couleur peut être
appelée blanche). Alors $\shift_{\gamma_1}(\ell) - \shift_{\gamma_2}(\ell)$ a
la parité du nombre de sommets singuliers de valence impaire intérieurs à $A$
qui se développent sur des points blancs. En particulier, deux anneaux ne
bordent une couronne que si ce nombre est pair.
\end{proposition}
\end{evidence}

\begin{proof}
Dans la somme de la proposition~\ref{prop:burgers}, un sommet de valence paire
a $R_i = \pm I$ et $(I - R_i)p_i \in \{0, 2p_i\}$\,; pour une valence
impaire, les deux composantes de $(I - R^{\pm1})(a, b) = (a \pm b,\ b \mp a)$
ont la parité de $a + b$. Les quarts de tour conservent la propriété que les
deux composantes ont une parité donnée, et les sommes ajoutent les parités, ce
qui donne le compte. Un quart de tour conserve les couleurs\,; une translation
entière avec $t_1 + t_2$ impair les échange, ce qui conserve la parité du
compte parce que les sommets de valence impaire intérieurs à $A$ sont en
nombre pair, leurs indices impairs ayant avec les indices pairs une somme
nulle (lemme~\ref{lem:gauss}). Développer un sommet le long d'un autre chemin
intérieur au disque découpé le fait tourner autour d'un autre sommet
singulier, ce qui conserve sa couleur.
\end{proof}

La règle vaut aussi sur $13\,000$ configurations aléatoires vérifiées avec
SymPy.

\paragraph{Parité sur une sphère.}
Chaque tour d'une sphère fermée porte sa propre condition de parité, qui ne
demande ni courbe simple ni holonomie.

\begin{evidence}{lean}{\lean{TMesh.parity_of_turn_sphere}, à partir de \lean{TMesh.bipartite_of_sphere} (le graphe des quads d'un T-mesh connexe à $P - A + N = 2$ est biparti, par \lean{QuadSurface.bipartite_of_euler} et \lean{TMesh.eulerQuad_eq_nodeCount}) et de \lean{TMesh.parity_of_turn_of_colouring}\,; la sphère est $P - A + N = 2$ et connexe, et la marche suit le tour à toutes longueurs entières équilibrées strictement positives\,; testé~: \lean{TMesh.IsColouring} décidé sur les 35 sphères de 50 instances aléatoires, 4\,388 tours}
\begin{proposition}[Parité d'un tour sur une sphère]\label{prop:parity}
Supposons la surface fermée de genre zéro et que le T-mesh admette des
longueurs strictement positives qui équilibrent chaque patch (ses longueurs
réelles en sont). Pour tout tour $\gamma$ et toutes longueurs entières $\ell$
équilibrant chaque patch, de signe quelconque,
\[
  \alpha_\gamma(\ell) \equiv \shift_\gamma(\ell) \pmod 2 .
\]
\end{proposition}
\end{evidence}

\begin{proof}
Supposons d'abord $\ell$ entier, équilibré et strictement positif~: les
patchs forment un maillage de quads de la sphère, dont le graphe est biparti,
puisque $H_1(S^2; \Z/2) = 0$ fait de tout cycle une somme de faces à quatre
arêtes. Dans une 2-coloration, la couleur du point $(i, j)$ d'un patch $P$ est
$i + j + \kappa_P \bmod 2$. Pour $M(x) = R^k x + t$, $t \in \Z^2$, la somme
des coordonnées de $M(x)$ a la parité de celle de $x$ plus $\pi(M) = t_1 +
t_2$. Soient $P_0, \dots, P_N = P_0$ les patchs traversés par $\gamma$ et
$D_k$ le déplacement qui développe $P_k$ dans $P_0$, avec $D_N =
T_\gamma(\ell)$. Un nœud de l'arc par lequel on passe de $P_k$ à $P_{k+1}$, de
coordonnées $y$, $y'$ avec $D_k(y) = D_{k+1}(y')$, a la couleur $(D_k y)_1 +
(D_k y)_2 + \pi(D_k) + \kappa_{P_k}$, et de même depuis $P_{k+1}$\,; donc
$\kappa_{P_k} + \pi(D_k)$ est constant, et en $k = N$, $\pi(T_\gamma(\ell))
\equiv 0$. Dans le repère du premier côté, celui de $P_0$ à un quart de tour
et des signes près, $T_\gamma(\ell) = (\alpha_\gamma(\ell),
\pm\shift_\gamma(\ell))$. Pour $\ell$ entier équilibré quelconque~: les
longueurs équilibrées forment un sous-espace rationnel qui a des points réels
strictement positifs, il existe donc des longueurs entières équilibrées
strictement positives $p$, et la forme linéaire $\alpha_\gamma +
\shift_\gamma$ prend en $\ell = (\ell + Np) - Np$, pour $N$ grand, la
différence de deux valeurs paires.
\end{proof}

Lean la vérifie sur le T-mesh combinatoire sans développement~: le sommet
voisin de la station est suivi le long de la marche, et la coloration $i + j +
\kappa_P$ du graphe des quads vient de $V - E + F = 2$ par algèbre linéaire
sur $\Z/2$ sur les chaînes du maillage de quads, la sphère étant $P - A + N =
2$ pour $P$ patchs, $A$ arcs et $N$ nœuds, connexe. La démonstration
n'utilise que le caractère biparti du graphe des quads\,;
elle tombe en défaut en genre supérieur (une grille $3 \times 3$ sur un tore a
des anneaux de trois quads) et avec des bords
(théorème~\ref{thm:parity-any}). Sur une sphère fermée, elle donne aussitôt~:

\begin{evidence}{core}{\lean{ring_even}, \lean{odd_pitch_odd_length}, \lean{repair_keeps_length_cannot_close}, \lean{odd_change_of_pitch_changes_length}\,; supposé~: comme la proposition ci-dessus}
\begin{corollary}\label{cor:parity}
Sur une surface fermée de genre zéro~:
\begin{enumerate}
\item[(a)] tout anneau de quads a un nombre pair de quads\,;
\item[(b)] une hélice de pas impair fait le tour en un nombre impair de quads
  par révolution\,;
\item[(c)] fermer une hélice de pas impair change la longueur de sa marche
  d'un nombre impair d'arêtes~: une réparation qui garde la longueur d'un tour
  ne peut pas le fermer si son pas est impair, et changer le pas d'une
  quantité impaire change toujours la longueur.
\end{enumerate}
\end{corollary}
\end{evidence}

\begin{proof}
(a) Un anneau est la bande d'une marche de décalage $\shift =
0$~\cite[prop.~2]{gardent2026loops}, donc $\alpha$ est pair. (b) et (c) sont
la proposition lue avant et après le changement.
\end{proof}

La partie (a) est le caractère biparti classique des quadrangulations de la
sphère~\cite[lemme~3.1]{kaiser2015colouring}. Pour un tour qui borde un
disque, l'argument de la règle du damier donne aux deux composantes de $B_D =
(\alpha_\gamma, \shift_\gamma)$ la parité du nombre de sommets de valence
impaire situés sur des points blancs, ce qui est la proposition pour ce tour.

\begin{evidence}{core}{\lean{colour_hol_trans}\,; supposé~: le tour est un mot en les lacets autour des points coniques}
\begin{remark}[Bords\,; une seconde démonstration]\label{rem:borders}
La proposition et la partie (a) du corollaire demandent que la surface soit
fermée~: un cylindre ouvert de trois quads de tour est un anneau de trois
quads. Une seconde démonstration, issue de la formalisation, ne demande ni le
graphe biparti ni des longueurs strictement positives~: tout lacet de
$\Sigma^\circ$ est un produit de lacets autour des sommets singuliers
(section~\ref{sec:fox}), des quarts de tour autour de points entiers qui
conservent la parité de $p_1 + p_2$, si bien que la translation de tout lacet
a une somme de coordonnées paire.
\end{remark}
\end{evidence}

\paragraph{Combien d'anneaux indépendants.}
Des anneaux qui ne se croisent pas sont des courbes fermées simples disjointes
de $\Sigma^\circ$\,; sur une sphère fermée à $n \ge 3$ sommets singuliers, au
plus $n - 3$ d'entre elles sont deux à deux non homotopes et non
périphériques (la taille d'une décomposition en pantalons), un anneau autour
d'un seul sommet singulier a un décalage nul (sa valence est alors un
multiple de quatre), et des anneaux homotopes ont le
même décalage, si bien que leurs conditions de fermeture sont de rang au plus
$n - 3$ ($3g - 3 + n + b$, plus au plus $b$ classes de colliers, en genre $g$
avec $b$ bords).

\paragraph{Certificats.}
Fermer un ensemble de tours dans des bornes est une question de faisabilité
linéaire.

\begin{evidence}{lean}{\lean{TMesh.not_closable_iff}, \lean{TMesh.closable_iff}}
\begin{remark}[Tours de signe fixe\,; certificats de Farkas]\label{rem:definite}
Si le décalage d'un tour est égal, modulo les équilibres, à une forme $\sum_a
c_a \ell_a$ dont tous les $c_a \ge 0$, alors pour tout $\ell \ge 0$ équilibré
le décalage $\shift_\gamma(\ell)$ est toujours positif ou nul, et il n'est nul
que si tout arc de coefficient $c_a > 0$ a une longueur nulle. Si l'un de ces
arcs est maintenu ouvert par une borne de séparation, le tour est une hélice
pour toute quantification qui respecte les bornes\,; il ne devrait jamais être
retenu, et un programme linéaire par tour trouve de tels tours. Plus
généralement, d'après le lemme de Farkas (avec $x = b + z$, $z \ge 0$), le
système $B\ell = 0$, $\shift_\gamma(\ell) = 0$ pour $\gamma \in H$, $\ell \ge
b$ n'a pas de solution exactement lorsqu'une combinaison $\sum_{\gamma \in H}
u_\gamma \shift_\gamma$ est égale, modulo les équilibres, à une forme $\sum_a
c_a \ell_a$ avec $c \ge 0$ et $\sum_a c_a b_a > 0$\,; et un arc $a$ a une
longueur nulle dans toute solution avec $\ell \ge 0$ exactement lorsqu'une
telle combinaison a $c_a > 0$. Les données étant entières, bornes comprises,
les multiplicateurs des décalages et des équilibres peuvent être pris entiers.
\end{remark}
\end{evidence}

\begin{evidence}{lean}{\lean{whole_solution_of_real}, \lean{clear_denominators}}
\begin{remark}[Des solutions entières existent à une certaine taille]\label{rem:exist}
Les contraintes d'équilibre et de fermeture sont homogènes, et les bornes de
séparation sont des minorations $\ell \ge b$ avec $b \ge 0$. Si elles ont une
solution réelle, elles en ont une rationnelle, sommet d'un polyèdre rationnel,
et la multiplier par le dénominateur commun $N$ donne des longueurs entières
$N\ell \ge Nb \ge b$ qui les satisfont toutes. Fermer en nombres entiers n'est
donc jamais impossible en soi\,; ce qui peut échouer, c'est fermer \emph{à une
taille donnée} (un nombre de quads donné) ou avec un écart aux cibles donné.
\end{remark}
\end{evidence}

\paragraph{Sur le corps masculin.}
Les énoncés de cette section valent partout où ils s'appliquent (mesuré). Les
applications de passage donnent le décalage exporté des $1\,299$ tours. Sur
les $1\,436$ côtés des $718$ tours simples, d'indice $4$ et de $4$ à $257$
sommets singuliers, l'holonomie est égale à la somme de la
proposition~\ref{prop:burgers}(a) dans ses deux composantes, alors que l'ordre
de dernière visite échoue sur $150$ des $158$ disques de la variante~2\,; la
règle du damier y vaut, ainsi que sur les $30\,148$ couronnes comprises entre
des tours simples disjoints, d'indice $0$\,; les $294$ paires sans sommet
singulier entre elles ont des décalages égaux\,; sur les dix couronnes qui
contiennent une paire 3--5, la proposition~\ref{prop:dislocation}(i) vaut, les
deux sommets n'étant jamais sur une même trace. La
proposition~\ref{prop:parity} vaut sur tous les tours et sur les $1\,838$
révolutions des bandes des huit maillages. Des $101$ tours qui se croisent une
fois, $97$ se découpent en $(3, 3, 2)$ et $4$, qui croisent leur segment de
fermeture, en $(2, 2, 4)$\,; $480$ se croisent jusqu'à $107$ fois, et seuls
$356$ des $765$ tours retenus sont simples~: la section~\ref{sec:fox} traite
toute courbe.

\section{Conflits entre tours}\label{sec:conflicts}

Les tours retenus du corps se ferment ensemble si l'on demande seulement aux
longueurs d'être positives ou nulles, et non d'appartenir à
$P_{\mathrm{eng}}$. Le tableau~\ref{tab:iis} liste des ensembles infaisables
\emph{irréductibles}, qui ne peuvent pas être fermés dans les bornes alors
que chacun de leurs sous-ensembles stricts le peut\,; les multiplicateurs de
Farkas ont été vérifiés en arithmétique exacte, et pour chaque tour omis, des
longueurs entières fermant les autres ont été trouvées par HiGHS et vérifiées
en arithmétique entière exacte (mesuré).

\begin{table}[t]
\begin{evidence}[unbreakable]{data}{les 25 certificats de ses 15 lignes, valides avec les multiplicateurs imprimés, sur l\'export du corps masculin d\'où viennent ces nombres (\texttt{make lean-check})}
\centering
\small
\setlength{\tabcolsep}{5pt}
\begin{tabular}{llllll}
\toprule
variantes & taille & tours & multiplicateurs & arcs bornés & vérif.\ exacte \\
\midrule
0, 1, 2, 4 & 1 & 42 (variante 0) & $-1$ & 1937 & variante 0 \\
1, 3, 5, 7 & 2 & 54, 61 (variante 1) & $1, -1$ & 1587 & --- \\
1, 3 & 2 & 45, 49 (variante 1) & $1, -1$ & 1069 & --- \\
1, 3, 7 & 2 & 26, 37 (variante 1) & $-1, 1$ & 619 & --- \\
5, 7 & 2 & 58, 94 (variante 5) & $1, -1$ & 858 & --- \\
5 & 2 & 49, 71 & $-1, 1$ & 619, 625 & --- \\
6 & 2 & 57, 58 & $1, 1$ & 1241 & --- \\
7 & 2 & 8, 15 & $1, 1$ & 191 & --- \\
1 & 3 & 36, 37, 39 & $1, 1, -1$ & 619 & oui \\
1 & 4 & 40, 56, 59, 61 & $1, 1, -1, -1$ & 1587 & oui \\
1 & 4 & 66, 71, 75, 76 & $-\tfrac76, -\tfrac56, -\tfrac13, -\tfrac12$ & 1967, 1980, 1998 & oui \\
3 & 4 & 73, 78, 82, 83 & $-\tfrac74, -\tfrac54, -\tfrac12, -\tfrac34$ & 1980, 1998 & oui \\
5 & 7 & 59, 91, 96, 98, 99, 157, 158 & $2, 4, 1, -4, -1, 1, -1$ & 1036 & oui \\
7 & 7 & 112--116, 145, 147 & $8, \tfrac12, 1, \tfrac12, \tfrac12, -\tfrac12, \tfrac12$ & 1157 & oui \\
5 & 15 & 0, 2--9, 11, 14, 16--18, 64 & entiers, jusqu'à 23 & 17 & 14 sur 15 \\
\bottomrule
\end{tabular}
\caption{Ensembles infaisables irréductibles de tours retenus sur le corps
masculin, bornes de séparation gardées (mesuré). Les numéros des tours, les
multiplicateurs et les arcs bornés sont ceux de la première variante
indiquée\,; «~arcs bornés~»~: les arcs maintenus à une longueur d'au moins un
dans le support de la forme à coefficients positifs ou nuls. Une autre
variante maintient d'autres arcs à une longueur d'au moins un, et son
certificat peut en borner d'autres. Les lignes sans vérification exacte viennent du programme
linéaire en virgule flottante\,; leurs tours sont chacun fermables seuls, si
bien qu'une paire est irréductible. Dans la dernière ligne, les longueurs
entières fermant l'ensemble privé du tour 3 n'ont pas été trouvées dans la
limite de temps.}
\label{tab:iis}
\end{evidence}
\end{table}

Dans quatre variantes, le tour autour du poignet gauche ne peut même pas être
fermé seul~: son décalage vaut $-(\ell_{1921} + \ell_{1937})$ modulo les
équilibres, et l'arc $1937$, long de $0,08$ arête, est maintenu à une
longueur d'au moins un (exact pour la variante 0\,; nous n'avons pas vérifié
si toute quantification valide respecte cette borne). On pourrait espérer que
tout autre conflit contienne deux tours parallèles dont la couronne
l'interdit. Le plus petit contre-exemple compte trois tours près de l'œil
droit, dont deux quelconques se ferment ensemble~: $\shift_{36} + \shift_{37}
- \shift_{39} = \ell_{604} + \ell_{619} + \ell_{620} + \ell_{637}$ modulo les
équilibres, avec $\ell_{619} \ge 1$. Tout ensemble trouvé est déjà infaisable
en nombres réels, si bien qu'aucun n'est une obstruction de parité. Les
conflits entre deux tours se trouvent dans des amas de singularités aux yeux,
aux mains et aux pieds, la somme ou la différence des décalages étant une
somme de deux à huit longueurs d'arcs, dont une bornée\,; pourtant, dans aucun
des $13$ que nous avons examinés (les paires du tableau~\ref{tab:iis} sauf
$49$, $71$ de la variante~5), les tours ne sont séparés par une paire 3--5 sur
une même trace, comme dans la proposition~\ref{prop:dislocation}(ii). Ce qui
les sépare est un escalier.

\begin{evidence}{core}{\lean{staircase}, \lean{staircase_rings}\,; supposé~: le vecteur entre les sommets est la somme le long de l'escalier}
\begin{proposition}[Escaliers]\label{prop:staircase}
Soient $\gamma_1$, $\gamma_2$ des tours et $c$, $c'$ des sommets singuliers
tels que
\[
  \shift_{\gamma_1}(\ell) - \shift_{\gamma_2}(\ell) = \varepsilon_1 d_1 +
  \varepsilon_2 d_2
\]
pour tout $\ell$ équilibré, avec $\varepsilon_1, \varepsilon_2 \in \{+1,
-1\}$ et $d = p_{c'} - p_c = d_1 e_1 + d_2 e_2$ l'écartement de $c'$ par
rapport à $c$, développé le long d'un chemin $\pi$ de $c$ à $c'$ (les
propositions~\ref{prop:dislocation}(i) et~\ref{prop:swap} donnent de telles
identités). Soit $\sigma$ un chemin de $c$ à $c'$ le long d'arcs $a_1, \dots,
a_r$, qui ne passe par aucun autre sommet singulier et qui est homotope à
$\pi$, extrémités fixées, dans $\Sigma$ privée des autres sommets singuliers.
Développé le long de $\sigma$, chaque arc $a_i$ suit une direction $v_i \in
\{\pm e_1, \pm e_2\}$. Si $\varepsilon_1 v_{i,1} + \varepsilon_2 v_{i,2} = 1$
pour tout $i$, de sorte que $\sigma$ est un \emph{escalier}, qui utilise deux
directions à un quart de tour l'une de l'autre et monte à travers la diagonale
sur laquelle les deux tours pourraient se fermer, alors
\[
  \shift_{\gamma_1}(\ell) - \shift_{\gamma_2}(\ell) = \sum_{i=1}^r \ell_{a_i}
\]
pour tout $\ell$ équilibré. Les deux tours ne sont tous deux des anneaux que
si tout arc de $\sigma$ a une longueur nulle. Si cette somme est égale, modulo
les équilibres, à une forme à coefficients positifs ou nuls contenant un arc
maintenu à une longueur d'au moins un, l'un des deux tours est une hélice pour
toute quantification qui respecte les bornes.
\end{proposition}
\end{evidence}

\begin{proof}
Développer le long de $\sigma$ ajoute $\ell_{a_i} v_i$ arc après arc, les
nœuds franchis étant des jonctions en T. Comme $\sigma$ et $\pi$ sont
homotopes à extrémités fixées hors des autres points coniques, $d = \sum_i
\ell_{a_i} v_i$, donc $\varepsilon_1 d_1 + \varepsilon_2 d_2 = \sum_i
\ell_{a_i}$, d'abord pour des longueurs équilibrées strictement positives,
puis pour toutes par linéarité.
\end{proof}

La démonstration s'applique à tout chemin sur chacune des étapes duquel la
forme de la différence des décalages vaut un\,; la
proposition~\ref{prop:dislocation}(ii) est le cas rectiligne. Les échanges
donnent les autres identités.

\begin{evidence}{core}{\lean{Motion.exchange}, \lean{exchange_quarter}\,; supposé~: l'homotopie entre les lacets}
\begin{proposition}[Échanger deux sommets]\label{prop:swap}
Soient $\gamma_1$, $\gamma_2$ des tours, et $\kappa$ un chemin du départ de
$\bar\gamma_1$ à celui de $\bar\gamma_2$ qui transporte la direction de marche
de $\gamma_2$ sur celle de $\gamma_1$, comme dans la
proposition~\ref{prop:burgers}(b)\,; soit $\lambda = \bar\gamma_1\, \kappa\,
\bar\gamma_2^{-1} \kappa^{-1}$. Supposons que $\lambda$ soit homotope dans
$\Sigma^\circ$ à $\lambda_c \lambda_{c'}^{-1}$, où $\lambda_c$ et
$\lambda_{c'}$ sont de petits lacets parcourus avec $c$, $c'$ à leur droite
autour de deux sommets singuliers de même valence $k$, atteints le long de
chemins qui les développent en $p_c$ et $p_{c'}$. Alors
\[
  \shift_{\gamma_1}(\ell) - \shift_{\gamma_2}(\ell) = \bigl\langle e_2,\
  (I - R^{4-k})(p_{c'} - p_c) \bigr\rangle .
\]
Pour la valence trois ou cinq, cela vaut $\pm d_1 \pm d_2$ avec $d = p_c -
p_{c'}$, comme pour une dislocation, et la proposition~\ref{prop:staircase}
s'applique.
\end{proposition}
\end{evidence}

\begin{proof}
Comme dans la proposition~\ref{prop:burgers}(b), $\hat h(\lambda)$ est la
translation de vecteur $\hat h(\bar\gamma_1) - \hat h(\bar\gamma_2)$, dont la
composante selon $e_2$ est $\shift_{\gamma_2}(\ell) - \shift_{\gamma_1}(\ell)$.
Avec $j = 4 - k$, $\hat h(\lambda_c) \circ \hat h(\lambda_{c'})^{-1}(x) =
R^j\bigl(R^{-j}(x - p_{c'}) + p_{c'} - p_c\bigr) + p_c = x + (I - R^j)(p_c -
p_{c'})$\,; pour $j = \pm1$, les deux composantes sont $\pm d_1 \pm d_2$.
\end{proof}

Les $13$ conflits se répartissent en six configurations
(tableau~\ref{tab:pairs}, mesuré). En A, la couronne entre deux tours simples
contient une paire 3--5, et l'escalier qui la relie, $\ell_{611} +
\ell_{624}$, long de $0,990$, est la différence des décalages réels\,; modulo
les équilibres, il est égal au certificat $\ell_{611} + \ell_{616} +
\ell_{619} + \ell_{635}$, qui maintient l'arc $619$ (exact)\,; aux longueurs entières du moteur, les deux
sommets de la paire sont à une arête l'un de l'autre sur une même ligne de la
grille, la proposition~\ref{prop:dislocation}(ii) avec $\delta = 1$ (son
hypothèse sur $\kappa$ n'étant pas vérifiée). En B, C et D, les tours se
croisent et leur combinaison s'enroule $+1$ et $-1$ fois autour de deux
sommets de valence trois, un échange en homologie (l'homotopie de la
proposition~\ref{prop:swap} n'a pas été vérifiée, l'identité de l'escalier
l'a été). En E et F, la combinaison fait une fois le tour de sommets d'indice
total quatre, comme un tour autour d'une calotte, et aucun escalier, ni
aucune somme de deux escaliers, ne s'est trouvé égal à elle.

\begin{table}[t]
\begin{evidence}[unbreakable]{measured}{sur l'export du corps masculin}
\centering
\small
\setlength{\tabcolsep}{4pt}
\newcolumntype{L}[1]{>{\raggedright\arraybackslash}p{#1}}
\begin{tabular}{@{}lL{4.1cm}L{3.9cm}L{4.4cm}L{2.0cm}@{}}
\toprule
& conflit & entre les tours & escalier & mécanisme \\
\midrule
A & $\shift_{37} - \shift_{26}$, variantes 1, 3, 7\,; œil droit\,;
  disjoints, simples
  & couronne~: 78 (val.\ 3), 76 (val.\ 5)
  & 2 arcs, $0,990$\,; arc maintenu 619 par les équilibres
  & dislocation \\
B & $\shift_{45} - \shift_{49}$, variantes 1, 3\,; main droite\,; simples,
  sécants
  & lentilles~: 102, 103 (val.\ 3)
  & 4 arcs, $1,731$\,; arc maintenu 1069 par les équilibres
  & échange \\
C & $\shift_{58} - \shift_{94}$, variantes 5, 7\,; main droite
  & 109, 118 (val.\ 3)
  & 3 arcs, $2,203$\,; arc maintenu 858 par les équilibres
  & échange \\
D & $\shift_{57} + \shift_{58}$, variante 6\,; pied gauche\,; en forme de
  huit
  & 155, 156 (val.\ 3)
  & 2 arcs, $1,452$, dont un maintenu
  & échange \\
E & $\shift_{54} - \shift_{61}$, variantes 1, 3, 5, 7\,; œil gauche
  & 11 sommets d'indice 3, et 207 (val.\ 3)
  & aucun\,; une chaîne fermée de 6 peignes
  & peignes \\
F & $\shift_{8} + \shift_{15}$, variante 7\,; pied droit
  & 14, 23, 24, et 15 (val.\ 3)
  & aucun\,; une chaîne fermée de 4 peignes autour d'une calotte
  & peignes \\
\bottomrule
\end{tabular}
\caption{Les conflits entre deux tours sur le corps masculin (mesuré). Les
numéros des tours sont ceux de la première variante indiquée, les signes ceux
des décalages exportés\,; «~entre les tours~»~: les sommets singuliers,
numérotés comme nœuds, autour desquels s'enroule la combinaison des deux
courbes\,; «~escalier~»~: un escalier situé entre eux dont la longueur est
égale à la combinaison modulo les équilibres, et l'endroit où se trouve l'arc
maintenu du certificat. A~: propositions~\ref{prop:dislocation}(i)
et~\ref{prop:staircase}\,; B--D~: propositions~\ref{prop:swap}
et~\ref{prop:staircase}\,; E, F~: section~\ref{sec:combs}.}
\label{tab:pairs}
\end{evidence}
\end{table}

\paragraph{Une séparatrice qui s'enroule.}
Le tour qui ne peut pas être fermé seul est d'une autre nature
(figure~\ref{fig:separatrix}).

\begin{evidence}{prose}{une étape en Lean, la conjugaison~: \lean{Motion.conj_translation}\,; sur 25 séparatrices qui s'enroulent de tores aléatoires, aucun strip voisin n'a eu le décalage $\delta$, l'homotopie n'étant pas vérifiée}
\begin{proposition}[Une séparatrice qui s'enroule]\label{prop:winding}
Soient $\sigma$ et $\tau$ des séparatrices d'un sommet singulier $v$ telles
que $\sigma$, suivie depuis $v$, s'arrête sur $\tau$ en un point $q$ situé à
la distance $\delta(\ell)$ de $v$ le long de $\tau$, une somme de longueurs
d'arcs. Soit $\lambda$ le lacet qui suit $\sigma$ de $v$ à $q$ puis revient le
long de $\tau$ jusqu'à $v$, écarté de $v$ vers l'intérieur du secteur qu'il
balaie. Si $\lambda$ est librement homotope dans $\Sigma^\circ$ à la courbe
$\bar\gamma$ d'un tour $\gamma$, la direction de marche de $\gamma$ allant sur
celle de $\sigma$, alors $\shift_\gamma(\ell) = \pm\delta(\ell)$ pour tout
$\ell$ équilibré, le signe dépendant du côté de $\sigma$ où se trouve $\tau$
($+$ sur la figure~\ref{fig:separatrix}(b), $-$ au poignet gauche du corps
masculin). Le tour n'est un anneau que si $\delta = 0$, c'est-à-dire si
$\sigma$ revient sur son propre sommet\,; si $\delta$ est égal, modulo les
équilibres, à une forme à coefficients positifs ou nuls contenant un arc
maintenu à une longueur d'au moins un, $\gamma$ est une hélice pour toute
quantification qui respecte les bornes.
\end{proposition}
\end{evidence}

\begin{proof}
Développons $\lambda$ depuis un point de $\sigma$ voisin de $v$, avec $v$ à
l'origine et $\sigma$ selon $e_1$. La séparatrice s'arrête sur $\tau$ à angle
droit, si bien que le retour est un segment de longueur $\delta$ selon $\pm
e_2$. Homotope à la courbe d'un tour, $\lambda$ a pour holonomie une
translation, qui envoie $v$ développé au départ sur $v$ développé à
l'arrivée, à la distance $\delta$ de $q$ transversalement à $\sigma$. Une
homotopie libre conjugue l'holonomie par un déplacement dont la rotation
transporte la direction de marche de $\gamma$ sur celle de $\sigma$, ici
l'identité.
\end{proof}

\begin{figure}[t]
\centering
% A separatrix that winds round a cylinder, and the turn it keeps shifted
% French copy of separatrix.tex: only the labels differ.
% (Proposition 8, the mechanism of the left wrist of the male body).
% \input inside a figure; needs \usepackage{tikz} and
% \usetikzlibrary{arrows.meta,decorations.pathreplacing}.
%
% (b) is the upper part of the band of dislocation.tex (b), the same T-mesh
% (check.py checks its numbers):
% cut open along the trace tau of the vertex v of valence five, the right side
% glued to the left side two units higher, (10,y) ~ (0,y+2) for y > 3; v is at
% (10,3) and (0,5) (and at (0,3), out of the picture with the lower band),
% the vertex c_3 of valence three at (1,4), its corner glued by a quarter-turn
% (see dislocation.tex). Traces: tau = the side x = 10 from v upwards (and its
% copy x = 0 from v upwards); sigma = the separatrix of v that leaves it to the
% right along y = 5; the separatrix of c_3 along y = 4, which stops on tau at
% n = (10,4); the separatrix of c_3 upwards, which stops on sigma at (1,5).
%   sigma goes once round, 10 units, and stops on tau at q = (10,5), at the
%   distance delta = l_e + l_f = 1 + 1 = 2 from v along tau, with
%   e = [v, n] and f = [n, q].
%   The turn gamma at o = 5.5, just above sigma, crosses tau at (10,5.5) =
%   (0,7.5): it comes back 2 higher, s_gamma = 2 = delta, length 10 = the
%   length of sigma (the along component of the holonomy of the loop lambda:
%   sigma from v to q, then tau back to v). gamma is a ring only if delta = 0,
%   that is if sigma comes back onto v.
%   By the balances of the two patches between y = 3, 4 and 5 (right of
%   x = 1), l_e equals the arc of c_3 downwards and l_f its arc upwards: if a
%   bound of separation holds either at length at least one, gamma is a helix
%   under every quantization meeting the bounds.
%   Parity (Proposition 5): 10 = 2 = 0 mod 2.
% (a) is a sketch of the same band on the cylinder, not to scale: sigma and
% gamma rise by delta = 2 rows in one round (the twist of the band above v).
\begin{tikzpicture}[x=0.46cm, y=0.46cm, font=\footnotesize,
    side/.style={draw=black, line width=0.5pt},
    quadgrid/.style={draw=black!18, line width=0.25pt, step={(1,1)}},
    surface/.style={fill=black!3},
    node dot/.style={circle, fill=black, inner sep=0pt, minimum size=2.6pt},
    separatrix/.style={draw=valfive, line width=1.1pt},
    gap/.style={draw=measured, line width=2pt},
    turn/.style={line width=1pt, -{Stealth[length=5pt]}},
    brace/.style={decorate, decoration={brace, amplitude=3pt}, line width=0.5pt},
    val3/.pic={\fill[valthree, draw=black, line width=0.3pt]
      (90:3.4pt) -- (210:3.4pt) -- (330:3.4pt) -- cycle;},
    val5/.pic={\fill[valfive, draw=black, line width=0.3pt]
      (90:3.4pt) -- (162:3.4pt) -- (234:3.4pt) -- (306:3.4pt) -- (18:3.4pt) -- cycle;}]
  \definecolor{loopspiral}{HTML}{D55E00}
  \definecolor{valthree}{HTML}{E8890C}
  \definecolor{valfive}{HTML}{7C3AAD}
  \definecolor{measured}{HTML}{14934A}

  % ---------------------------------------------------------------- (a)
  % The cylinder; a point (x, y) of the band above v sits at the angle
  % 90 - 36 x degrees (v and tau in front) and the height y + 2x/10 - 5.
  \begin{scope}[shift={(-7.5,1.4)}]
    \pgfmathsetmacro{\R}{2.4}
    \pgfmathsetmacro{\k}{0.22}
    \pgfmathsetmacro{\H}{6.6}
    % Back halves of sigma and gamma (x from 2.5 to 7.5), through the cylinder.
    \draw[valfive!70, dashed, line width=0.7pt]
      plot[domain=2.5:7.5, samples=40, variable=\x]
        ({\R*cos(90-36*\x)}, {1.6 + 0.2*\x - \k*\R*sin(90-36*\x)});
    \draw[loopspiral!70, dashed, line width=0.7pt]
      plot[domain=2.5:7.5, samples=40, variable=\x]
        ({\R*cos(90-36*\x)}, {2.1 + 0.2*\x - \k*\R*sin(90-36*\x)});
    % The cylinder, translucent.
    \fill[black!25, fill opacity=0.18] (-\R,0) -- (-\R,\H) arc[start angle=180, end angle=360,
      x radius=\R, y radius=\k*\R] -- (\R,0) arc[start angle=0, end angle=-180,
      x radius=\R, y radius=\k*\R] -- cycle;
    \draw[line width=0.5pt] (-\R,0) -- (-\R,\H) (\R,0) -- (\R,\H);
    \draw[line width=0.5pt] (0,\H) ellipse[x radius=\R, y radius=\k*\R];
    \draw[line width=0.5pt] (\R,0) arc[start angle=0, end angle=-180,
      x radius=\R, y radius=\k*\R];
    \draw[line width=0.35pt, black!40, dashed] (\R,0) arc[start angle=0, end angle=180,
      x radius=\R, y radius=\k*\R];
    % tau, in front, from v upwards.
    \draw[separatrix] (0,1.6-\k*\R) -- (0,\H-\k*\R);
    % Front halves of sigma (from v) and gamma.
    \draw[separatrix]
      plot[domain=0:2.5, samples=20, variable=\x]
        ({\R*cos(90-36*\x)}, {1.6 + 0.2*\x - \k*\R*sin(90-36*\x)});
    \draw[separatrix]
      plot[domain=7.5:10, samples=20, variable=\x]
        ({\R*cos(90-36*\x)}, {1.6 + 0.2*\x - \k*\R*sin(90-36*\x)});
    \draw[loopspiral, line width=1pt]
      plot[domain=0:2.5, samples=20, variable=\x]
        ({\R*cos(90-36*\x)}, {2.1 + 0.2*\x - \k*\R*sin(90-36*\x)});
    \draw[loopspiral, line width=1pt, -{Stealth[length=5pt]}]
      plot[domain=7.5:9.9, samples=20, variable=\x]
        ({\R*cos(90-36*\x)}, {2.1 + 0.2*\x - \k*\R*sin(90-36*\x)});
    % The other separatrices of v, started, and c_3 next to it.
    \foreach \a in {162, 234, 306}
      \draw[valfive!55, line width=0.8pt] (0,1.6-\k*\R) -- ++(\a:0.55);
    \pic at ({\R*cos(54)}, {0.8 - \k*\R*sin(54)}) {val3};
    % The gap delta on tau, from v to q.
    \draw[gap] (0,1.6-\k*\R) -- (0,3.6-\k*\R);
    \pic at (0,1.6-\k*\R) {val5};
    \node[node dot] at (0,3.6-\k*\R) {};
    \node[left=3pt] at (0,1.6-\k*\R) {$v$};
    \node[right=1pt] at (0,3.6-\k*\R) {$q$};
    \node[valfive, left] at (0,4.8-\k*\R) {$\tau$};
    \node[valfive, right] at (\R,1.95) {$\sigma$};
    \node[loopspiral, right] at (\R,2.75) {$\gamma$};
    \node[anchor=north] at (0,-2.0) {\small (a) le cylindre};
  \end{scope}

  % ---------------------------------------------------------------- (b)
  \begin{scope}
    \fill[surface] (0,1.6) rectangle (10,3) (1,3) rectangle (10,4) (0,4) rectangle (10,8.6);
    \draw[quadgrid] (0,1.6) grid (10,3);
    \draw[quadgrid] (1,3) grid (10,4);
    \draw[quadgrid] (0,4) grid (10,8.6);
    % Traces: the loop of v at y = 3, the separatrices of c_3, and the cut
    % edges of the corner of c_3; then tau and sigma.
    \draw[side] (0,1.6) -- (0,3) (10,1.6) -- (10,3);
    \draw[side] (0,3) -- (10,3) (1,4) -- (10,4) (1,4) -- (1,5);
    \draw[side] (1,4) -- (1,3) (1,4) -- (0,4) -- (0,5);
    \draw[separatrix] (10,3) -- (10,8.6) (0,5) -- (0,8.6);
    \draw[separatrix] (0,5) -- (10,5);
    \draw[gap] (10,3) -- (10,5);
    \foreach \p in {(1,3), (0,4), (10,4), (0,6), (1,5), (10,5), (0,7)}
      \node[node dot] at \p {};
    % The turn gamma, and where it comes back.
    \draw[turn, loopspiral] (0,5.5) -- (10.1,5.5);
    \draw[loopspiral, line width=0.5pt, fill=white] (0,5.5) circle (1.8pt);
    \draw[loopspiral, dashed, line width=0.8pt, -{Stealth[length=4pt]}] (0,7.5) -- (2,7.5);
    \draw[loopspiral, line width=0.5pt, fill=white] (0,7.5) circle (1.8pt);
    \node[above, loopspiral] at (5,5.5) {$\gamma$};
    \draw[brace] (-0.3,5.5) -- (-0.3,7.5) node[midway, left=3pt] {$s_\gamma = 2$};
    % Singular vertices and names.
    \pic at (1,4) {val3};
    \foreach \p in {(0,3), (10,3), (0,5)} \pic at \p {val5};
    \node[above right=-1pt] at (1,4) {$c_3$};
    \node[left=1pt] at (0,4.9) {$v$};
    \node[right=1pt] at (10,3) {$v$};
    \node[right=1pt] at (10,4) {$n$};
    \node[right=1pt] at (10,5) {$q$};
    \node[valfive, below] at (5,5) {$\sigma$};
    \node[valfive, right] at (10,8) {$\tau$};
    \node[measured, left] at (10,3.5) {$e$};
    \node[measured, left] at (10,4.5) {$f$};
    \draw[brace, decoration={mirror}] (11.3,3) -- (11.3,5)
      node[midway, right=3pt] {$\delta = \ell_e + \ell_f = 2$};
  \end{scope}
  \node[anchor=north] at (5,0.9) {$s_\gamma = \delta = 2$~: un anneau seulement si $\sigma$ revient sur $v$};
  \node[anchor=north] at (5,-0.6) {\small (b) ouvert le long de $\tau$};
\end{tikzpicture}

\caption{Une séparatrice qui s'enroule (proposition~\ref{prop:winding}).
(a)~Sur un cylindre comme un poignet, la séparatrice $\sigma$ du sommet $v$ de
valence cinq fait une fois le tour et s'arrête sur une autre séparatrice
$\tau$ de $v$. (b)~Ouvert le long de $\tau$~: $\sigma$, de longueur $10$,
s'arrête en $q$, à $\delta = \ell_e + \ell_f = 2$ de $v$, et le tour $\gamma$
voisin de $\sigma$ revient $2 = \delta$ rangées plus haut\,; il n'est un
anneau que si $\sigma$ revient sur $v$. Au poignet gauche du corps masculin,
$\delta = \ell_{1920} + \ell_{1921}$, maintenu strictement positif par une
borne de séparation.}
\label{fig:separatrix}
\end{figure}

Au poignet gauche, le tour ($42$ de la variante 0, de longueur $6,99$, de
décalage réel $-0,305$) borde une calotte de quatre sommets de valence trois,
dont la somme de Burgers donne son décalage comme une forme de signe non
fixé. Le signe vient de l'extérieur~: la séparatrice $964$ du sommet $246$,
de valence cinq, parcourt $6,994$ arêtes en faisant une fois le tour du
poignet et s'arrête sur une autre de ses séparatrices à $\delta = \ell_{1920}
+ \ell_{1921} = 0,305$ de $246$\,; le décalage vaut $-\delta$, et
$\ell_{1920}$ est égal à la longueur de l'arc maintenu $1937$, tous deux
modulo les équilibres~: d'où le certificat et $\shift \le -1$ à toute
longueur entière admissible. Nous avons vérifié l'identité, non l'homotopie
(mesuré).

\paragraph{Conflits plus grands.}
Les trois tours de l'œil droit s'enroulent autour d'une paire 3--5, et leur
combinaison est la longueur d'un escalier de trois arcs qui passe par l'arc
maintenu $619$\,; les quatre des yeux s'enroulent autour d'une paire 3--5 et
de deux sommets de valence trois, une dislocation et un échange en homologie.
Les multiplicateurs fractionnaires du tableau~\ref{tab:iis} viennent de ce
qu'on demande aux arcs bornés de peser un au total\,; à la main gauche, où les
tours se croisent de $10$ à $46$ fois, les multiplicateurs primitifs donnent
\[
  -(7\shift_{66} + 5\shift_{71} + 2\shift_{75} + 3\shift_{76}) = 2F =
  10\,\ell_{1805} + 2\ell_{1817} + 4\ell_{1853} + 2(\ell_{1870} + \ell_{1879}
  + \ell_{1969} + \ell_{1998} + \ell_{2029})
\]
modulo les équilibres. Le coefficient $5$ est une dislocation~: une
séparatrice du sommet $231$ (valence trois) s'arrête sur une séparatrice de
$228$ (valence cinq), et l'arc $1805$, long de $0,063$ arête, relie cette
jonction à $228$, si bien que $\ell_{1805}$ est l'écart transverse de la paire
par rapport à la ligne du champ qui les relie presque, la
proposition~\ref{prop:dislocation}(ii) avec $\delta = \ell_{1805}$ (hypothèse
sur $\kappa$ non vérifiée)\,; leurs enroulements tordus
(section~\ref{sec:fox}) contribuent exactement $10\,\ell_{1805}$ (mesuré).

\paragraph{Dislocations locales.}
Les courbes d'un ensemble irréductible ne bordent pas nécessairement une même
région plane autour d'un amas de singularités~: ce n'est le cas ni des
configurations B à F ni des ensembles de quatre et de sept tours. Ce que les
conflits ont en commun est local.

\begin{evidence}{open}{}
\begin{question}[Les conflits sont-ils des sommes de dislocations locales\,?]\label{conj:local}
Sur une surface fermée de genre zéro, soit $H$ un ensemble irréductible de
tours qui ne peuvent pas être fermés ensemble dans les bornes de séparation.
Existe-t-il une combinaison $\sum_{\gamma \in H} u_\gamma \shift_\gamma$ à
multiplicateurs entiers qui certifie le conflit (remarque~\ref{rem:definite})
et qui soit égale, modulo les équilibres, à une somme à poids entiers positifs
ou nuls de
\begin{enumerate}
\item[(i)] longueurs d'escaliers (proposition~\ref{prop:staircase}) entre un
  sommet de valence trois et un de valence cinq, ou entre deux sommets de même
  valence échangés comme dans la proposition~\ref{prop:swap}, et
\item[(ii)] décalages de calottes~: des formes de signe fixe portées par des
  courbes qui font une fois le tour de sommets singuliers d'indice total
  quatre, comme le tour du poignet et sa séparatrice
  (proposition~\ref{prop:winding}),
\end{enumerate}
un terme au moins contenant un arc maintenu ouvert\,?
\end{question}
\end{evidence}

La réponse est oui, avec un seul terme et aux homotopies non vérifiées près,
pour le poignet, les configurations A à D et l'œil droit\,; la
section~\ref{sec:combs} tranche la question sur le corps~: non telle qu'elle
est posée, oui, mais sans contenu, lue au sens large.

\section{Toute combinaison de tours}\label{sec:fox}

Les nombres d'enroulement sont la classe d'homologie d'un cycle pondéré
$\sum_\gamma u_\gamma \bar\gamma$, qui est commutative, alors que l'holonomie
ne l'est pas\,; le calcul différentiel libre de Fox~\cite{fox1953free} garde
l'ordre. Soit $\Sigma$ fermée de genre zéro, de sommets singuliers $c_1,
\dots, c_n$ (la proposition~\ref{prop:burgers-any} ajoute les bords et les
anses). Fixons $x_0$ intérieur à un patch, avec son repère, et des chemins
disjoints $\pi_i$ de $x_0$ à $c_i$\,; soit $\lambda_i$ le lacet qui suit
$\pi_i$, fait le tour de $c_i$ en le gardant à sa droite, et revient. Le
groupe fondamental de $\Sigma^\circ$ est engendré par les $\lambda_i$ avec une
seule relation, leur produit dans l'ordre où les chemins quittent $x_0$. Soit
$p_i$ la position développée de $c_i$ au bout de $\pi_i$ et $R_i = R^{4 -
k(c_i)}$, de sorte que $\hat h(\lambda_i) = \rho_i$, $\rho_i(x) = R_i(x - p_i)
+ p_i$. Écrivons $\hat h(\gamma)(x) = \rho(\gamma)x + \tau(\gamma)$\,; alors
$\tau(\alpha\beta) = \tau(\alpha) + \rho(\alpha)\,\tau(\beta)$. Soit $F$ le
groupe libre sur $x_1, \dots, x_n$, envoyé sur le groupe fondamental par $x_i
\mapsto \lambda_i$, muni des dérivées de Fox $\partial x_j/\partial x_i =
\delta_{ij}$, $\partial(uv)/\partial x_i = \partial u/\partial x_i +
u\,\partial v/\partial x_i$. Étendue linéairement à l'algèbre du groupe,
$\rho$ prend les valeurs $a_0 I + a_1 R + a_2 R^2 + a_3 R^3$, et
l'augmentation $\epsilon$ envoie un tel élément sur $a_0 + a_1 + a_2 + a_3$.
La première identité ci-dessous est la propriété qui définit les dérivées de
Fox pour un homomorphisme croisé~\cite{fox1953free}, et les $p_i$, positions
développées des sommets singuliers, sont définis à un déplacement commun
près\,; ce qui est nouveau est la
lecture sur des combinaisons de tours.

\begin{evidence}{core}{\lean{fox_burgers}, \lean{fox_fundamental}, \lean{fox_burgers_combination}\,; supposé~: la présentation du groupe fondamental par les lacets autour des sommets singuliers}
\begin{proposition}[Formule de Fox--Burgers]\label{prop:fox}
Pour tout mot $w$ de $F$,
\[
  \tau(w) = \sum_{i=1}^n \rho\Bigl(\frac{\partial w}{\partial x_i}\Bigr)
  (I - R_i)\,p_i .
\]
Soient $\gamma_1, \dots, \gamma_m$ des tours, $\kappa_r$ un chemin de $x_0$
au départ de $\bar\gamma_r$, $f_r$ la direction $e_2$ de $\gamma_r$ ramenée en
$x_0$ le long de $\kappa_r$ (l'un des vecteurs $\pm e_1$, $\pm e_2$), et $w_r$
un mot représentant $\kappa_r \bar\gamma_r \kappa_r^{-1}$. Pour des
multiplicateurs $u_r$, définissons l'\emph{enroulement tordu} du cycle pondéré
autour de $c_i$,
\[
  W_i = -\sum_{r=1}^m u_r\, \rho\Bigl(\frac{\partial w_r}{\partial x_i}\Bigr)^{\!\top} f_r
  \in \Z^2 \text{ (pour des } u_r \text{ entiers)} .
\]
Alors, pour toutes longueurs équilibrées,
\[
  \sum_{r=1}^m u_r\, \shift_{\gamma_r}(\ell) = \sum_{i=1}^n \bigl\langle W_i,\
  (I - R_i)\,p_i \bigr\rangle ,
  \qquad
  \sum_{i=1}^n (I - R_i)^\top W_i = 0 ,
\]
si bien que la somme ne change pas quand tous les $p_i$ subissent une même
translation~: c'est une combinaison d'écartements $p_i - p_j$. Le nombre
d'enroulement ordinaire de $\sum_r u_r \bar\gamma_r$ autour de $c_i$ est, à
une constante commune à tous les $i$ près, $w_i = \sum_r u_r\,
\epsilon(\partial w_r/\partial x_i)$. Les mots $w_r$ sont définis à la
relation près, et les $W_i$ en conséquence~: insérer un conjugué $g\varrho
g^{-1}$ de la relation $\varrho$ dans $w_r$, après un préfixe $a$ de $w_r$,
retranche $u_r \rho(a g\, \partial \varrho/\partial x_i)^\top f_r$ de chaque $W_i$
(avec $a$ vide, $u_r \rho(g\, \partial \varrho/\partial x_i)^\top f_r$), et
laisse la somme inchangée.
\end{proposition}
\end{evidence}

\begin{proof}
L'application $w \mapsto \tau(w)$ et le membre de droite vérifient tous deux
$\phi(uv) = \phi(u) + \rho(u)\phi(v)$, et ils coïncident sur les lettres et
sur le mot vide, donc partout, par récurrence sur la longueur (pour
$x_i^{-1}$, utiliser $0 = \phi(x_i x_i^{-1})$). Pour un tour, $\hat
h(\kappa_r \bar\gamma_r \kappa_r^{-1})$ est la translation de vecteur $M_r
T_r$, où $T_r$ est la translation de $\bar\gamma_r$ dans son propre repère et
$M_r$ la rotation de $\hat h(\kappa_r)$\,; donc $\shift_{\gamma_r} = -\langle
M_r e_2, M_r T_r \rangle = -\langle f_r, \tau(w_r) \rangle$, et $\langle f,
\rho(D) b\rangle = \langle \rho(D)^\top f, b \rangle$ donne la formule. La
formule fondamentale $\sum_i (\partial w/\partial x_i)(x_i - 1) = w - 1$
donne, par $\rho$, $\sum_i \rho(\partial w_r/\partial x_i)(R_i - I) =
\rho(w_r) - I = 0$, la rotation d'un tour étant triviale\,; transposée et
sommée avec les $u_r f_r$, c'est la seconde identité. Enfin, $\epsilon(\partial
w/\partial x_i)$ est la somme des exposants de $x_i$ dans $w$, la classe
d'homologie de $w$ sur la base des $\lambda_i$, définie à la relation près,
dans laquelle chaque $\lambda_i$ apparaît une fois.
\end{proof}

Pour un tour simple, $\rho(\partial w/\partial x_i) = Q_i$ et la formule est
la proposition~\ref{prop:burgers}(a). En général, un cycle peut faire
plusieurs fois le tour de $c_i$ après des rotations différentes\,; modulo
$2$, ces différences disparaissent.

\begin{evidence}{core}{\lean{checkerboard_word}, \lean{checkerboard_combination}\,; supposé~: comme dans la formule ci-dessus}
\begin{corollary}[Règle du damier pour toute combinaison]\label{cor:checker}
Sur une surface fermée de genre zéro, soient $\ell$ entier et équilibré et les
multiplicateurs $u_r$ entiers. Colorions chaque sommet singulier $c_i$ selon
la parité $\chi_i$ de $p_{i,1} + p_{i,2}$\,; elle ne dépend pas des chemins
$\pi_i$, et pour des longueurs strictement positives c'est la couleur de $c_i$
dans la bipartition du maillage de quads. Alors
\[
  \sum_r u_r\, \shift_{\gamma_r}(\ell) \equiv \sum_r u_r\, \alpha_{\gamma_r}(\ell)
  \equiv \sum_{i:\ k(c_i) \text{ impair}} w_i\, \chi_i \pmod 2 ,
\]
où les $w_i$ sont les nombres d'enroulement de $\sum_r u_r \bar\gamma_r$,
quelle que soit la constante commune et quelle que soit l'orientation de
chaque $\bar\gamma_r$. En particulier, une combinaison dont le cycle fait un
nombre pair de fois le tour de chaque sommet de valence impaire prend des
valeurs paires à toute longueur entière.
\end{corollary}
\end{evidence}

\begin{proof}
Puisque $(R - I)^2 = -2R$, on a $R^k(I - R^{\pm1}) \equiv I - R^{\pm1}$
modulo $2$, donc $\rho(D)(I - R_i) \equiv \epsilon(D)(I - R_i)$ pour $k(c_i)$
impair, et $\equiv 0$ pour $k(c_i)$ pair, où $I - R_i \in \{0, 2I\}$\,; et
$\langle f, (I - R^{\pm1})p\rangle \equiv p_1 + p_2$ pour $f \in \{\pm e_1,
\pm e_2\}$. Cela donne la somme de la proposition~\ref{prop:fox} modulo $2$,
et la direction $e_1$ de chaque tour celle de $\sum_r u_r
\alpha_{\gamma_r}$. Les rotations autour de points entiers et les quarts de
tour conservent $p_1 + p_2 \bmod 2$, donc $\chi_i$ ne dépend pas des chemins.
Une translation entière de l'origine inverse toutes les couleurs et ajoute
$\sum_{k(c_i)\text{ impair}} w_i$, pair parce que la rotation du cycle est
triviale ($\sum_i w_i (4 - k(c_i)) \equiv 0 \pmod 4$)\,; ajouter $q$ à tous les
$w_i$ ajoute $q \sum_{k(c_i)\text{ impair}} \chi_i$, la valeur modulo $2$ de
la relation, un lacet de translation nulle, donc pair. Inverser un tour ne
change que des signes. Pour des longueurs strictement positives, la parité
développée colore proprement le graphe des quads, comme dans la
proposition~\ref{prop:parity}.
\end{proof}

Le corollaire contient les propositions~\ref{prop:checkerboard},
\ref{prop:dislocation}(iii) et~\ref{prop:parity}. Comme pour les défauts des
cristaux~\cite{kleman2008disclinations} et
dans~\cite[\S7]{izmestiev2013torus}, les nombres d'enroulement ne déterminent
pas l'holonomie~:

\begin{evidence}{lean}{\lean{commutator_holonomy}, \lean{winding_numbers_do_not_suffice}}
\begin{proposition}[Les nombres d'enroulement ne suffisent pas]\label{prop:commutator}
Soient $c_a$, $c_b$ des sommets singuliers de valence impaire développés en
des points distincts $p_a \neq p_b$. Le lacet $\lambda_a \lambda_b
\lambda_a^{-1} \lambda_b^{-1}$ fait zéro fois le tour de chaque sommet
singulier, et son holonomie est la translation de vecteur $(I - R_a)(I -
R_b)(p_a - p_b)$, de longueur $2|p_a - p_b| \neq 0$. Aucune fonction des
nombres d'enroulement ne donne l'holonomie des lacets à rotation triviale.
\end{proposition}
\end{evidence}

\begin{proof}
Pour $w = x_a x_b x_a^{-1} x_b^{-1}$, on a $\rho(\partial w/\partial x_a) =
\rho(1 - x_a x_b x_a^{-1}) = I - R_b$ et $\rho(\partial w/\partial x_b) =
\rho(x_a - w) = R_a - I$, les rotations commutant\,; la
proposition~\ref{prop:fox} donne la translation, et pour des valences
impaires, $I - R_a$ et $I - R_b$ sont $\sqrt2$ fois des rotations. Les sommes
d'exposants de $w$ sont nulles.
\end{proof}

\paragraph{Sur le corps masculin.}
Découpé en une seule carte le long d'un arbre couvrant des nœuds, le T-mesh
donne à chaque tour un mot, et la proposition~\ref{prop:fox} sa translation,
pour les $1\,299$ tours. Le corollaire~\ref{cor:checker} vaut sur tous et sur
les $14$ premiers ensembles du tableau~\ref{tab:iis}, multiplicateurs
fractionnaires ramenés à des entiers (l'ensemble de $15$ tours n'a pas été
essayé)\,; à la
main gauche, le cycle fait un nombre pair de fois le tour de chaque sommet de
valence impaire, ce qui explique que $-2F$ soit pair. Quatre tours vérifient
$w_{18} = w_{22} + w_{101} + w_{65}$ autour de chaque sommet singulier alors
que leurs translations ne le vérifient pas, quels que soient les repères, à
des quarts de tour près~: une «~formule de Burgers des enroulements~» qui
aurait des nombres pour poids est fausse (mesuré).

\subsection{Peignes}\label{sec:combs}

Les enroulements tordus dépendent de la carte\,; les peignes, non. Dans les
T-meshes d'openclare, tracés à partir des seules séparatrices des sommets
singuliers, toute trace d'une surface fermée part d'un sommet singulier (des
lignes caractéristiques ou des traces amorcées casseraient cette propriété\,;
nous la supposons). Pour un arc $a$, orienté selon $e_a$ du nœud $n_0$ au nœud
$n_1$, le \emph{pied} de $n_j$ est le nœud lui-même s'il est singulier, et
sinon le sommet singulier d'où part la trace qui passe par $n_j$
perpendiculairement à $a$. Nous supposons ce sommet unique~: une seule trace
perpendiculaire à $a$ rencontre $n_j$ (et non deux, venant de côtés opposés),
et ce n'est ni une connexion de séparatrices ni une trace refermée sur
elle-même\,; sinon, un pied doit être choisi, et des choix différents donnent
des dislocations différentes. Le \emph{peigne} $\sigma_a$ de $a$ va du pied
$c_0$ de $n_0$ le long de sa trace jusqu'à $n_0$, le long de $a$ jusqu'à
$n_1$, et le long de l'autre trace jusqu'au pied $c_1$ de $n_1$, en ne passant
que par des jonctions en T (figure~\ref{fig:comb}). Le \emph{graphe des pieds}
a pour sommets les sommets singuliers de valence impaire et pour arêtes les
arcs dont les deux pieds sont des sommets distincts de valence impaire.

\begin{figure}[t]
\centering
% The comb of an arc (Section "Combs" and Lemma "Combs").
% French copy of comb.tex: it has no words, only the comma of the pairing
% differs (written for icomma, a comma followed by a space).
% \input inside a figure; needs \usepackage{tikz} and
% \usetikzlibrary{arrows.meta,decorations.pathreplacing}.
%
% The T-mesh is that of dislocation.tex (b) and separatrix.tex (b): a band cut
% open along the trace tau of the vertex v of valence five, the right side
% glued to the left side two units higher above y = 3 and at the same height
% below; v is drawn at (0,3), (10,3) and (0,5), the vertex c_3 of valence three
% at (1,4), and the empty corner next to c_3 is glued by a quarter-turn about
% c_3 (see dislocation.tex). Every node is at a whole position; check.py
% checks the numbers below.
%
% The arc a is the piece of the separatrix sigma of v from the T-junction
% n_0 = m = (1,5), where the separatrix of c_3 upwards stops on sigma, to the
% T-junction n_1 = q = (10,5), where sigma stops on tau; e_a = (1,0), l_a = 9.
% Foot of n_0: the trace through m perpendicular to a is the separatrix of c_3
% upwards, which starts at c_3: c_0 = c_3. Foot of n_1: the trace through q
% perpendicular to a is tau, which starts at v: c_1 = v. The comb sigma_a runs
% from c_3 up to m (1 unit), along a to q (9 units), and down tau through the
% T-junction n = (10,4) to v (2 units): it passes through T-junctions only.
% Developed along it, p_{c_0} = (1,4) and p_{c_1} = (10,3), and
%   < e_a, p_{c_1} - p_{c_0} > = < (1,0), (9,-1) > = 9 = l_a
% (Lemma "Combs"); the teeth only add the across component -1.
% The feet are distinct, of valences 3 and 5: the comb reads the length of a
% as that of a dislocation (Proposition 3); with two feet of valence 3 it would
% be an exchange (Proposition 7), with a single foot a winding separatrix
% (Proposition 8).
\begin{tikzpicture}[x=0.5cm, y=0.5cm, font=\footnotesize,
    side/.style={draw=black, line width=0.5pt},
    quadgrid/.style={draw=black!18, line width=0.25pt, step={(1,1)}},
    surface/.style={fill=black!3},
    node dot/.style={circle, fill=black, inner sep=0pt, minimum size=2.6pt},
    comb/.style={draw=measured, line width=1.8pt, line join=round},
    arc mark/.style={line width=2.2pt, draw=black},
    dimension/.style={line width=0.5pt, {Bar[width=5pt]}-{Bar[width=5pt]}},
    val3/.pic={\fill[valthree, draw=black, line width=0.3pt]
      (90:3.4pt) -- (210:3.4pt) -- (330:3.4pt) -- cycle;},
    val5/.pic={\fill[valfive, draw=black, line width=0.3pt]
      (90:3.4pt) -- (162:3.4pt) -- (234:3.4pt) -- (306:3.4pt) -- (18:3.4pt) -- cycle;}]
  \definecolor{valthree}{HTML}{E8890C}
  \definecolor{valfive}{HTML}{7C3AAD}
  \definecolor{measured}{HTML}{14934A}

  \begin{scope}
    \fill[surface] (0,1.6) rectangle (10,3) (1,3) rectangle (10,4) (0,4) rectangle (10,8);
    \draw[quadgrid] (0,1.6) grid (10,3);
    \draw[quadgrid] (1,3) grid (10,4);
    \draw[quadgrid] (0,4) grid (10,8);
    % All traces of the T-mesh in the picture.
    \draw[side] (0,1.6) -- (0,3) (10,1.6) -- (10,8) (0,5) -- (0,8);
    \draw[side] (0,3) -- (10,3) (1,4) -- (10,4) (1,4) -- (1,5) (0,5) -- (10,5);
    \draw[side] (1,4) -- (1,3) (1,4) -- (0,4) -- (0,5);
  \end{scope}
  % The comb of a: from its foot c_3 up to m, along a, down tau to its foot v.
  \draw[comb] (1,4) -- (1,5) -- (10,5) -- (10,3);
  \draw[comb, line width=3pt] (1,5) -- (10,5);
  \foreach \p in {(1,3), (0,4), (10,4), (0,6), (1,5), (10,5), (0,7)}
    \node[node dot] at \p {};
  \pic at (1,4) {val3};
  \foreach \p in {(0,3), (10,3), (0,5)} \pic at \p {val5};
  % Names.
  \node[above] at (5.5,5) {$a$};
  \draw[-{Stealth[length=4pt]}, line width=0.5pt] (6.6,5.75) -- (8,5.75)
    node[right, inner sep=1pt] {$e_a$};
  \node[above left=-1pt] at (1,5) {$n_0$};
  \node[above right=-1pt] at (10,5) {$n_1$};
  \node[below right=-1pt] at (1,4) {$c_0 = c_3$};
  \node[right=2pt] at (10,3) {$c_1 = v$};
  \node[left=1pt] at (0,4.95) {$v$};
  \node[left=1pt] at (0,3) {$v$};
  \node[measured, right=1pt] at (10,4.5) {$\sigma_a$};
  \node[right=1pt] at (10,7.4) {$\tau$};
  % The developed feet and the e_a component of their separation.
  \draw[dimension] (1,2.2) -- (10,2.2) node[midway, fill=black!3, inner sep=1.5pt]
    {$\langle e_a, p_{c_1} - p_{c_0}\rangle = \ell_a = 9$};
  \draw[black!50, dotted, line width=0.6pt] (1,2.2) -- (1,4) (10,2.2) -- (10,3);
\end{tikzpicture}

\caption{Le peigne $\sigma_a$ d'un arc $a$ (lemme~\ref{lem:comb}), $\ell_a =
9$~: depuis le pied $c_0$ d'une extrémité de $a$, de valence trois, en
remontant sa trace, le long de $a$, puis en descendant l'autre trace
perpendiculaire jusqu'au pied $c_1$, de valence cinq. Toutes les étapes sauf
$a$ sont orthogonales à $a$, donc $\langle e_a, p_{c_1} - p_{c_0} \rangle = 9
= \ell_a$. Ce T-mesh est celui des figures~\ref{fig:dislocation}(b)
et~\ref{fig:separatrix}(b)~: l'escalier, la séparatrice qui s'enroule et le
peigne sont trois lectures du même décalage $2$.}
\label{fig:comb}
\end{figure}

\begin{evidence}{core}{\lean{comb_reading}\,; les pieds existent sur 850 des 3\,009 arcs de 50 instances aléatoires\,; supposé~: les pieds existent et les traces perpendiculaires se développent en droites}
\begin{lemma}[Peignes]\label{lem:comb}
Soient $p_{c_0}$, $p_{c_1}$ les pieds d'un arc $a$, développés le long de son
peigne. Alors, pour toutes longueurs équilibrées,
\[
  \langle e_a,\ p_{c_1} - p_{c_0} \rangle = \ell_a .
\]
Si $c_0 \neq c_1$ sont de valence impaire, $2\ell_a$ est, au signe près, la
forme de la classe tordue $\lambda_{c_0}\lambda_{c_1}^{\pm1} \otimes v$ des
lacets autour des deux pieds atteints le long du peigne, pour un $v$ entier
tel que $(I - R^{\pm1})^\top v = 2e_a$~: la paire est une dislocation si les
valences diffèrent (proposition~\ref{prop:dislocation}) et un échange si elles
sont égales (proposition~\ref{prop:swap}). Si $c_0 = c_1$, $\ell_a$ est le
décalage d'une séparatrice qui s'enroule (proposition~\ref{prop:winding}).
\end{lemma}
\end{evidence}

\begin{proof}
Chaque étape du peigne sauf $a$ suit une trace perpendiculaire à $a$, sans
passer par un point conique, et ajoute un vecteur orthogonal à $e_a$\,;
l'étape le long de $a$ ajoute $\ell_a e_a$. Le lacet $\lambda_{c_0}
\lambda_{c_1}^{\pm1}$, avec le signe qui rend sa rotation triviale, a pour
translation $\pm(I - R^{\pm1})(p_{c_0} - p_{c_1})$ comme dans les
propositions~\ref{prop:dislocation} et~\ref{prop:swap}\,; l'accoupler à $v$
donne $\pm 2\ell_a$, et $v$ est entier parce que $I - R^{\pm1}$ envoie $\Z^2$
sur les vecteurs dont les coordonnées ont la même parité, parmi lesquels
$2e_a$. Si $c_0 = c_1$, la seconde trace est une séparatrice du même sommet
qui s'arrête sur la première à la distance $\ell_a$.
\end{proof}

Une \emph{marche fermée de peignes} est une suite d'arcs $a_1, \dots, a_k$
dont les pieds sont $c_{i-1}$ et $c_i$, de valence impaire, avec $c_k =
c_0$\,; sa \emph{lecture} est $\sum_i \ell_{a_i}$. Une marche de longueur un
est une séparatrice qui s'enroule, une marche de longueur deux entre des
valences trois et cinq est l'escalier de la proposition~\ref{prop:staircase},
et une marche autour de quatre sommets de valence trois est une calotte.

\begin{evidence}{prose}{sur 50 instances aléatoires, 13 des 17 certificats exacts dont les arcs ont des pieds de valence impaire échouent à l'hypothèse sur la partie impaire, 6 sur des sphères}
\begin{proposition}[Les certificats sont des sommes de peignes]\label{thm:combs}
Soit $c = \sum_a c_a\ell_a$, $c_a \in \Z_{\ge0}$, un certificat
(remarque~\ref{rem:definite}) dont les arcs ont des pieds de valence impaire
et dont la partie impaire, les arcs de coefficient $c_a$ impair, est un cycle
modulo~$2$ du graphe des pieds. Alors $c$ est une somme à poids entiers
positifs ou nuls de lectures de peignes $2\ell_a$, chacune la forme d'une
dislocation ou d'un échange (lemme~\ref{lem:comb}), et de lectures de marches
fermées de peignes. Avec des poids réels, toute forme à coefficients positifs
ou nuls supportée par des arcs dont les deux pieds sont de valence impaire,
distincts ou égaux, est une combinaison à coefficients positifs ou nuls de
lectures de peignes $2\ell_a$ et de séparatrices qui s'enroulent $\ell_a$,
sans condition de parité.
\end{proposition}
\end{evidence}

\begin{proof}
La partie impaire se décompose en marches fermées le long de ses arêtes
(Veblen), et ce qui reste est $\sum_a \lfloor c_a/2 \rfloor\,2\ell_a$. Avec
des poids réels, $c = \sum_a \tfrac{c_a}{2}\,(2\ell_a)$ sur les arcs à deux
pieds distincts, et $c_a\ell_a$ sur ceux dont les pieds sont un même sommet.
\end{proof}

La démonstration est élémentaire\,: le contenu est dans l'hypothèse. Que la
partie impaire d'un certificat soit un cycle du graphe des pieds est constaté
sur les $25$ certificats du corps (exact), et n'est pas général\,: sur des
T-mesh aléatoires qui ne viennent pas du moteur (les tests de propriétés du
paquet de données, \texttt{experiments/property}), sous les bornes $\ell \ge
1$ sur chaque arc, $13$ des $17$ certificats exacts dont les arcs ont des
pieds de valence impaire ont une partie impaire qui n'est pas un cycle, $6$
d'entre eux sur des sphères, et pour aucun des $13$ un représentant modulo les
balances ne peut l'être, même modulo $2$\,; le cas typique est un seul arc
$\ell_a$ entre deux pieds distincts, le décalage d'un tour étant le décalage
d'une paire 3--5. Le corollaire~\ref{cor:checker} ne donne pas la parité non
plus, puisqu'il contrôle celle de $p_{i,1} + p_{i,2}$ alors qu'un peigne lit
une seule coordonnée (sur la figure~\ref{fig:comb}, les deux pieds ont la
même couleur et $\ell_a = 9$).

\begin{evidence}{open}{une attente}
\begin{remark}[Les marches fermées comme cycles tordus]
Nous nous attendons à ce que la lecture d'une marche fermée de peignes,
$\sum_i \langle e_{a_i} - e_{a_{i+1}},\ p_{c_i}\rangle$ par une sommation par
parties, soit la forme de Burgers d'un cycle tordu voisin de ses peignes\,;
mais les deux peignes qui passent par $c_i$ le développent dans des cartes qui
diffèrent d'une rotation, et une démonstration doit développer toute la
marche dans une même carte. Nous n'utilisons les marches fermées que comme
formes.
\end{remark}
\end{evidence}

\begin{evidence}{tested}{les 35 sphères sans valence $4k$ parmi 50 instances aléatoires~: rang égal à la dimension sur chacune\,; sur les tores le terme de période le fait échouer (11 sur 15)}
\begin{proposition}[L'accouplement de Burgers]\label{prop:pairing}
Supposons la surface fermée de genre zéro, sans sommet singulier de valence
divisible par quatre, et toute trace du T-mesh issue d'un sommet singulier.
Alors l'application qui envoie une classe de l'homologie tordue
$H_1(\Sigma^\circ; \R^2_\rho)$ sur sa forme $\sum_i \langle W_i, (I - R_i)p_i
\rangle$ sur les longueurs équilibrées est surjective sur le dual des
longueurs équilibrées. Son noyau contient la classe transverse à chaque
connexion de séparatrices. Si certains sommets singuliers ont une valence
divisible par quatre, aucune forme ne dépend de leurs positions développées,
et l'application n'est pas nécessairement surjective.
\end{proposition}
\end{evidence}

\begin{proof}
Si toutes ces formes s'annulent en un $\ell$ équilibré, tout lacet de
$\Sigma^\circ$ à rotation triviale a une holonomie sans translation
(proposition~\ref{prop:fox}). Le commutateur de $\lambda_a$ et $\lambda_b$ a
une rotation triviale et la translation $(I - R_a)(I - R_b)(p_a - p_b)$
(proposition~\ref{prop:commutator})\,; aucun $R_i$ n'est l'identité, donc $I
- R_i$ est inversible ($\sqrt2$ fois une rotation, ou $2I$), et tous les
sommets singuliers se développent en un même point $q$. Comme les $\rho_i$
fixent $q$ et engendrent l'holonomie sur une sphère fermée, tout sommet
singulier se développe en $q$ le long de tout chemin\,; tout arc a des pieds,
donc tous les écarts de peigne du lemme~\ref{lem:comb} sont nuls, et $\ell =
0$. L'image n'a pas d'autre zéro commun que $0$. À travers une connexion,
l'écartement de ses deux sommets est nul transversalement à la trace, et la
classe tordue qui accouple la direction transverse avec elle a la forme
nulle. Un sommet de valence divisible par quatre a $(I - R_i)p_i = 0$~: sa
position est une période relative, que les longueurs peuvent déplacer sans
changer aucune holonomie.
\end{proof}

\paragraph{Sur le corps masculin.}
Le lemme~\ref{lem:comb} vaut sur les $2\,069$ arcs, dont $2\,046$ ont deux
pieds distincts de valence impaire~: presque tout arc est une dislocation ou
un échange local. Sur ce T-mesh (variante~1), chaque nœud ne rencontre qu'une
trace perpendiculaire, et un seul arc, $1676$, aboutit sur la connexion de
séparatrices, qui joint les sommets $204$ et $210$\,; nous avons pris pour
pied le sommet d'où la connexion a été tracée, $204$, et aucun certificat de la
variante~1 ne contient cet arc. L'homologie tordue est de dimension $518 = 2(n - 2)$ et
l'accouplement de rang $517$, de noyau la classe transverse à l'unique
connexion de séparatrices\,; nous conjecturons que, sans sommet de valence
divisible par quatre, le noyau est engendré par les classes transverses aux
connexions. Les $25$ certificats des conflits examinés ont tous un
représentant dont la partie impaire est un cycle du graphe des pieds (exact).
Le poignet est un peigne dont les deux pieds sont en $246$, la séparatrice qui
s'enroule\,; A, C et D sont des marches fermées de deux peignes (A,
l'escalier $\ell_{611} + \ell_{624}$), B de quatre, l'œil droit de trois, E de
six, d'indice total $-2$, et F de quatre autour d'une calotte\,; la main
gauche vaut cinq fois la dislocation $\ell_{1805}$ plus huit lectures de
peignes sur un arbre de sept arcs, et l'ensemble de sept tours de la
variante~7 est un arbre de sept peignes (mesuré).

\paragraph{Lue au sens large, la question~\ref{conj:local} est vraie et vide.}
En admettant toute lecture de peigne et toute marche fermée de peignes, la
proposition~\ref{thm:combs} y répond pour tout certificat dont la partie
impaire est un cycle du graphe des pieds, ce qui est le cas de tous ceux du
corps\,; mais cela découle du lemme de Farkas et d'une parité, et ne dit rien
des configurations qui entrent en conflit. \emph{Lue au sens strict, elle est
fausse}~: si on la restreint aux escaliers de poids un, aux marches monotones
autour d'une calotte et aux séparatrices qui s'enroulent, des longueurs
équilibrées rationnelles exactes rendent toutes ces formes positives ou nulles
et les certificats de E et F négatifs, $-2,9967$ et $-1,9973$ (vérifié en
arithmétique entière). Ces formes sont en nombre infini, mais chacune est la
longueur d'une marche sur un graphe dont les arêtes portent les longueurs des
arcs, si bien qu'une recherche de plus court chemin en trouve une négative
s'il en existe. En virgule flottante, le même test sépare aussi les quatre
tours des yeux et place la main gauche à l'intérieur du cône (mesuré). On ne
sait pas si tout certificat se décompose en peignes portés par des ensembles
disjoints de sommets, chacun de son signe~: les arbres de la main gauche et de
la variante~7 partagent des sommets.

\section{Bords et anses}\label{sec:beyond}

Les modèles à bords (les orbites des yeux de Suzanne, la bouche de Jerry) et à
anses (Bob, un tore\,; part et cage) demandent les formes ci-dessous, mesurées
sur les T-meshes de ces modèles du banc avec une version du moteur non
enregistrée. Un nœud de bord $b$ est un coin d'angle $a(b)\pi/2$,
d'\emph{indice} $2 - a(b)$\,; pour une composante de bord $\beta_j$,
$L_j(\ell)$ est la somme des longueurs de ses arcs.

\begin{evidence}{tested}{à $\Theta = 0$ et sans bord, 3\,833 régions de 50 instances aléatoires\,; le terme de bord n'est pas testé}
\begin{lemma}[Gauss--Bonnet avec bords]\label{lem:gauss-border}
Soit $D$ une région bordée par des courbes de tours et par des composantes de
bord. Alors
\[
  \sum_{c \in D} (4 - k(c)) + \sum_{b \in D \cap \partial\Sigma} (2 - a(b))
  = 4\chi(D) - \Theta,
\]
où $\Theta$ est la rotation des courbes des tours en leurs coins, en quarts de
tour~: nulle pour le côté d'une courbe simple.
\end{lemma}
\end{evidence}
\begin{proof}
C'est la formule de Gauss--Bonnet pour une surface plate à points coniques et
à bord rectiligne par morceaux, les nœuds de bord étant des coins de rotation
$(2 - a)\pi/2$.
\end{proof}

Un côté d'un tour simple qui contient $b$ composantes de bord sans coin et $h$
anses a pour indice $4(1 - b - 2h)$, et peut ne contenir aucun sommet
singulier\,; le lemme vaut sur les $1\,076$ côtés de tours de Suzanne et de
Jerry (mesuré). Pour des longueurs strictement positives, l'homomorphisme du
théorème suivant est la parité des cycles du graphe des quads, un invariant
d'homologie à coefficients dans $\Z/2$ de toute
quadrangulation~\cite[lemme~3.1]{kaiser2015colouring}, étendu ici aux tours à
toutes longueurs équilibrées.

\begin{evidence}{core}{\lean{parity_any_surface}, \lean{colour_comp_trans}, \lean{colour_rotAbout_trans}, \lean{colour_border}, \lean{colour_turn}\,; testé~: 2\,027 strips de graphes quad bipartis et 1\,960 couronnes de 50 instances aléatoires\,; supposé~: que $\pi \circ \hat h$ passe au quotient $H_1(\Sigma\,; \Z/2)$, et que sa valeur sur une classe est la parité de la longueur d'un cycle du graphe quad dans cette classe}
\begin{theorem}[Parité d'un tour sur toute surface]\label{thm:parity-any}
Soit $\ell$ entier et équilibré. Il existe un homomorphisme
$\Omega_\ell\colon H_1(\Sigma; \Z/2) \to \Z/2$ tel que, pour tout tour,
\[
  \alpha_\gamma(\ell) + \shift_\gamma(\ell) \equiv \Omega_\ell([\bar\gamma]) \pmod 2 .
\]
Sa valeur sur une composante de bord est la parité de $L_j(\ell)$, la somme de
tous les $L_j(\ell)$ est paire, et pour des longueurs strictement positives
$\Omega_\ell(z)$ est la parité de la longueur de tout cycle du graphe des
quads dans la classe $z$\,; $\Omega_\ell = 0$ exactement lorsque le graphe des
quads est biparti. Sur une surface de genre zéro à bords, $\alpha_\gamma +
\shift_\gamma$ a la parité de la longueur totale des composantes de bord
situées de l'un ou l'autre côté de $\bar\gamma$, ou du nombre de fois que
$\bar\gamma$ traverse un système quelconque d'arcs disjoints qui relient deux
à deux les composantes de bord de longueur impaire.
\end{theorem}
\end{evidence}
\begin{proof}
Comme dans la proposition~\ref{prop:parity}, $\gamma \mapsto \pi(\hat
h(\gamma))$ est un homomorphisme, nul sur les lacets autour des sommets
singuliers (quarts de tour autour de points entiers)\,; à valeurs dans un
groupe d'exposant deux, il se factorise par $H_1(\Sigma; \Z/2)$, ces lacets
engendrant le noyau de $H_1(\Sigma^\circ) \to H_1(\Sigma)$. Pour un tour,
$\hat h(\bar\gamma)$ vaut $(\alpha_\gamma, \pm\shift_\gamma)$ à un quart de
tour près\,; une composante de bord revient déplacée de $\sum_a \pm\ell_a
v_a$, de la parité de $L_j(\ell)$, et les composantes ont pour somme le bord,
nul en homologie. Le long d'un cycle du graphe des quads, chaque arête change
la couleur. En genre zéro, les composantes de bord engendrent $H_1(\Sigma;
\Z/2)$ et $\bar\gamma$ est homologue à celles d'un de ses côtés\,; la
dernière forme est la dualité de Lefschetz.
\end{proof}

Sur Suzanne, dont les orbites sont des bords de $23$ et $25$ arêtes, le tour
autour d'une orbite est un anneau de $23$ quads aux longueurs d'une variante
et de $24$ dans une autre, contrairement à la proposition~\ref{prop:parity}\,;
son décalage est nul à toute longueur équilibrée~:

\begin{evidence}{core}{\lean{collar}\,; supposé~: l'holonomie du tour est conjuguée à celle de la composante de bord, et $\alpha_\gamma > 0$}
\begin{proposition}[Colliers]\label{prop:collar}
Supposons que la courbe d'un tour $\gamma$ et une composante de bord
$\beta_j$ sans coin bordent une couronne ne contenant aucun sommet singulier.
Alors $\shift_\gamma(\ell) = 0$ et $\alpha_\gamma(\ell) = L_j(\ell)$ pour
toutes longueurs équilibrées.
\end{proposition}
\end{evidence}
\begin{proof}
L'holonomie de $\beta_j$ est la translation de longueur $L_j$ selon un axe du
champ. Découpée le long d'un chemin de $\bar\gamma$ à $\beta_j$, la couronne
est un disque sans point conique, si bien que la translation de $\bar\gamma$
est $L_j u$ avec $u \in \{\pm e_1, \pm e_2\}$, comme dans la
proposition~\ref{prop:burgers}(b). Pour des longueurs strictement positives,
$\alpha_\gamma > 0$, donc $u = \pm e_1$\,; les identités sont linéaires et
s'étendent à toutes les longueurs équilibrées.
\end{proof}

Le théorème prédit la parité des $660$ tours de Suzanne, de Jerry et d'un
dôme, dont $61$ impairs, et de toute boucle fermée des maillages produits
($294$ sur le corps masculin)\,; la proposition vaut sur les $40$ colliers\,;
sur Bob, $\Omega_\ell$ dépend des longueurs, les trois variantes lui donnant
les valeurs $(0, 1)$, $(0, 0)$ et $(1, 0)$ sur deux lacets d'anse (mesuré).
Pour les combinaisons, $\sum_r u_r(\alpha_{\gamma_r} + \shift_{\gamma_r})
\equiv \sum_r u_r\, \Omega_\ell([\bar\gamma_r])$\,; les couleurs du
corollaire~\ref{cor:checker} peuvent dépendre des chemins, mais le compte sur
les points blancs donne encore la parité lorsque toute courbe est contenue
dans une région plane sans bord.

\begin{evidence}{core}{\lean{fox_general}, \lean{holGen_rotations}, \lean{commutator_motion}\,; supposé~: la présentation du groupe fondamental par les lacets nommés}
\begin{proposition}[Somme de Burgers avec bords et anses]\label{prop:burgers-any}
Supposons que les lacets $\lambda_i$ autour des sommets singuliers, $\beta_j$
autour des composantes de bord et $a_k, b_k$ autour des anses engendrent le
groupe fondamental de $\Sigma^\circ$. Pour tout mot $w$ en ces lacets,
\[
  \tau(w) = \sum_i \rho\Bigl(\frac{\partial w}{\partial x_i}\Bigr)(I - R_i)\,p_i
  + \sum_{g \in \{\beta_j, a_k, b_k\}} \rho\Bigl(\frac{\partial w}{\partial g}\Bigr)\tau(g).
\]
Pour un tour, le décalage est la somme de Burgers des sommets singuliers plus
un \emph{terme de périodes}, linéaire en les longueurs et non déterminé par
les positions développées des sommets singuliers. Une composante de bord sans
coin a pour holonomie la translation de longueur $L_j$ selon sa direction
développée\,; une anse contenue dans la région comprise entre deux tours
contribue le commutateur $(I - \rho(b))\tau(a) - (I - \rho(a))\tau(b)$, nul
lorsque ses deux lacets ont une rotation triviale.
\end{proposition}
\end{evidence}
\begin{proof}
Les deux membres vérifient $\phi(uv) = \phi(u) + \rho(u)\phi(v)$ et
coïncident sur les générateurs. Le commutateur de $x \mapsto Ax + a$ et de $x
\mapsto Bx + b$, à rotations qui commutent, est la translation de vecteur $(I
- B)a - (I - A)b$.
\end{proof}

Avec son terme de bord, la somme reproduit les $314$ côtés contenant un bord
sur lesquels elle a été testée, et sans lui aucun des $586$ côtés de ce type
de Suzanne et de Jerry (mesuré).

\begin{evidence}{core}{\lean{commutator_translation}, \lean{commutator_translations}\,; supposé~: le bord du disque découpé est le commutateur des deux courbes}
\begin{proposition}[Deux tours sécants sur un tore]\label{prop:torus}
Sur une surface fermée de genre un, supposons que les courbes de deux tours
soient simples et se croisent une fois. Découpée le long des deux, la surface
est un disque $D$ qui contient tous les sommets singuliers, et sa somme de
Burgers est nulle~: $\sum_i Q_i (I - R_i) p_i = 0$, dans l'ordre de première
visite, pour toutes longueurs équilibrées. Plus généralement, si $\kappa$ est
un lacet simple qui croise une fois la courbe d'un tour, de rotation $R^j$,
alors $(I - R^j)\,T_\gamma = \pm B_D$~: pour $j \neq 0$, la translation du
tour est fixée par les sommets singuliers\,; pour $j = 0$, c'est une période
libre et les sommets singuliers vérifient $B_D = 0$.
\end{proposition}
\end{evidence}
\begin{proof}
Le bord de $D$ est le commutateur des deux courbes, dont les holonomies sont
des translations et commutent\,; la proposition~\ref{prop:burgers}(a)
s'applique à $D$. Avec $\kappa$ de rotation $R^j$, le commutateur est la
translation de vecteur $(I - R^j)T_\gamma$.
\end{proof}

Sur Bob, de genre un, avec dix sommets de valence trois et dix de valence
cinq, la somme de Burgers est nulle pour les $304$ paires de tours simples qui
se croisent une fois, et pour aucun de $90$ ordres mélangés (mesuré). Avec une
seule paire 3--5, les deux points coniques se développeraient en un même
point\,; qu'un tel tore n'existe pas est connu~\cite{izmestiev2013torus}, et
c'est aussi une conséquence du théorème d'Abel
($c_5 - c_3$ n'est principal sur une courbe elliptique que si $c_5 = c_3$),
une obstruction à l'existence, non à la fermeture. Avec $j = 0$, la
proposition est l'argument du commutateur de~\cite[\S7]{izmestiev2013torus}.
Le conflit rencontré sur Bob est d'une autre nature~: $3\shift_{33} -
2\shift_{31}$, pour deux tours qui font respectivement trois et deux fois le
tour de la même anse, annule la période et est égal à deux longueurs d'arcs en
un sommet de valence trois, dont l'une est maintenue ouverte (mesuré).

La proposition~\ref{prop:linear}, le lemme~\ref{lem:gauss} pour les disques et
les couronnes, les propositions~\ref{prop:burgers} à~\ref{prop:checkerboard}
(pour les couronnes), \ref{prop:staircase}, \ref{prop:swap},
\ref{prop:winding} et~\ref{prop:commutator}, les
remarques~\ref{rem:definite} et~\ref{rem:exist}, le
théorème~\ref{thm:forced} et le corollaire~\ref{cor:chain} ne demandent pas de
sphère fermée\,; le lemme~\ref{lem:eight}, la proposition~\ref{prop:parity},
le corollaire~\ref{cor:parity}, le compte des anneaux indépendants, le cadre
de la proposition~\ref{prop:fox}, le corollaire~\ref{cor:checker} et la
question~\ref{conj:local} en demandent une\,; l'annexe~\ref{app:orders} traite
les maillages triangulaires.

\section{Les tours que toute quantification dans les bornes laisse décalés}\label{sec:forced}

Retournées, les obstructions minorent le nombre de tours que toute
quantification appartenant à un ensemble $P$ de longueurs laisse décalés.

\begin{evidence}{lean}{\lean{forced_helices}, \lean{card_ge_of_disjoint}, \lean{tauP_ge_of_disjoint}, \lean{TMesh.forced_of_certificates}}
\begin{theorem}[Décalages forcés]\label{thm:forced}
Soit $H$ un ensemble fini de tours d'un T-mesh et $P$ un ensemble non vide de
longueurs qui contient toutes les quantifications considérées~: par exemple
les longueurs équilibrées qui respectent des bornes de séparation données,
$\ell \ge b$, éventuellement avec un plafond $\frac1n \sum_a |\ell_a - t_a|
\le D$ sur l'écart moyen aux cibles $t$, ou avec un nombre de quads donné.
Soit
\[
  \tau_P(H) = \min\bigl\{\, |X| : X \subseteq H,\ \text{les tours de } H
  \setminus X \text{ peuvent être fermés ensemble dans } P \,\bigr\}.
\]
Alors, pour tout $\ell \in P$, entier ou non, au moins $\tau_P(H)$ tours de
$H$ ont $\shift_\gamma(\ell) \neq 0$. Si $H$ contient $k$ sous-ensembles deux
à deux disjoints dont aucun ne peut être fermé dans $P$, alors $\tau_P(H) \ge
k$.
\end{theorem}
\end{evidence}

\begin{proof}
En $\ell \in P$, les tours de décalage nul sont fermés par $\ell$ lui-même,
donc les tours décalés forment un ensemble $X$ comme dans la définition. Un
ensemble dont le retrait laisse un ensemble fermable rencontre chacun des $k$
sous-ensembles disjoints.
\end{proof}

Le théorème est élémentaire\,; son contenu est dans ce qui rend des
sous-ensembles infaisables. Il n'utilise ni convexité, ni linéarité, ni
longueurs entières. Les ensembles minimaux non fermables sont des
sous-systèmes infaisables irréductibles (IIS), $\tau_P(H)$ est la taille d'une
couverture minimale des IIS~\cite{parker1996iis,amaldi2003maxfs}, et le
théorème est la borne d'empilement de ce problème de couverture.

\begin{evidence}{core}{\lean{sign_change_between}, \lean{chain_of_dislocations}\,; supposé~: les identités entre décalages consécutifs}
\begin{corollary}[Chaînes de dislocations]\label{cor:chain}
Soient $\gamma_0, \gamma_1, \dots, \gamma_m$ des tours, et supposons que,
pour $i = 1, \dots, m$,
\[
  \shift_{\gamma_{i-1}}(\ell) - \shift_{\gamma_i}(\ell) = \varepsilon_i\,
  F_i(\ell) \quad \text{pour tout } \ell \text{ équilibré},
\]
avec $\varepsilon_i \in \{+1, -1\}$ et des formes $F_i$ strictement positives
aux points entiers de $P$\,; par exemple, $F_i$ la longueur d'un escalier
(proposition~\ref{prop:staircase}) entre un sommet de valence trois et un de
valence cinq situés entre $\gamma_{i-1}$ et $\gamma_i$
(propositions~\ref{prop:dislocation} et~\ref{prop:burgers}(b)), l'escalier
contenant un arc maintenu à une longueur d'au moins un. Soit $q$ le nombre
d'indices $i < m$ tels que $\varepsilon_{i+1} \neq \varepsilon_i$. Alors, pour
tout $\ell \in P$ entier, au plus $q + 1$ des tours $\gamma_i$ sont des
anneaux~: au moins $m - q$ sont des hélices. En particulier, $m + 1$ tours
parallèles emboîtés, séparés par $m$ dislocations dont les escaliers montent
tous dans le même sens, laissent au moins $m$ hélices.
\end{corollary}
\end{evidence}

\begin{proof}
Si $\gamma_i$ et $\gamma_j$, $i < j$, sont des anneaux, $\sum_{k=i+1}^{j}
\varepsilon_k F_k(\ell) = 0$ avec $F_k(\ell) > 0$, donc un changement de
signe se trouve entre $i$ et $j$\,; entre deux anneaux consécutifs se trouve
un changement de signe qui leur est propre.
\end{proof}

% evidence: none, elle dit ce que compte le théorème
\begin{remark}[Ce que compte la borne]
Le théorème compte les tours de décalage non nul\,; aux longueurs où un tour
n'est pas une bande du maillage~\cite[prop.~2]{gardent2026loops}, ou
lorsqu'une même bande en spirale porte plusieurs tours de $H$, ce n'est pas un
compte d'hélices.
\end{remark}

La borne est atteinte, par la correspondance classique entre couvertures
minimales des IIS et sous-systèmes faisables
maximaux~\cite{amaldi2003maxfs,parker1996iis}~:

\begin{evidence}{cited}{\cite{amaldi2003maxfs,parker1996iis}\,; (ii) aussi en Lean~: \lean{meets_infeasible}}
\begin{theorem}[Les décalages forcés comme nombre transversal\,; classique]\label{thm:transversal}
Disons que $Y \subseteq H$ est \emph{fermable} dans $P$ si un $\ell \in P$
vérifie $\shift_\gamma(\ell) = 0$ pour tout $\gamma \in Y$, et appelons
\emph{cœur} un ensemble non fermable dont tous les sous-ensembles stricts sont
fermables. Pour tout $P$~: (i) $H \setminus X$ est fermable si et seulement si
$X$ rencontre tout cœur\,; par conséquent, $\tau_P(H)$ est la plus petite
taille d'un ensemble qui rencontre tout cœur\,; (ii) pour tout $\ell \in P$,
les tours que $\ell$ laisse décalés rencontrent tout cœur\,; (iii) tout
ensemble de ce type de taille minimale est exactement l'ensemble des tours
laissés décalés par un certain $\ell \in P$\,; (iv) $P \subseteq P'$ implique
$\tau_P(H) \ge \tau_{P'}(H)$.
\end{theorem}
\end{evidence}

\begin{proof}
La fermabilité est héréditaire~: si $X$ évite un cœur, $H \setminus X$ le
contient, et un $H \setminus X$ non fermable contient un cœur. L'ensemble des
tours de décalage nul en $\ell$ est fermé par $\ell$. Pour un $X$ de taille
minimale, une longueur qui ferme $H \setminus X$ décale un sous-ensemble de
$X$ qui rencontre tout cœur, donc $X$ tout entier.
\end{proof}

Ainsi, $\tau_P(H)$ est compris entre le plus grand nombre de cœurs disjoints
et la taille du plus petit ensemble transversal des cœurs trouvés jusque-là,
et il est connu exactement dès que le complémentaire de cet ensemble est
fermable\,; la boucle d'\emph{ensemble transversal implicite}
(\emph{implicit hitting set}) de Parker et Ryan~\cite{parker1996iis} alterne
les deux\,; ici, les cœurs viennent des routines de filtrage de Chinneck et
Dravnieks~\cite{chinneck1991locating}, qui isolent un sous-système
irréalisable minimal.

\subsection{Longueurs entières et taille}\label{sec:whole}

\begin{evidence}{core}{\lean{closable_at_any_size}, \lean{smul_mem_system}, \lean{whole_solution_of_real}\,; supposé~: le nombre de quads est homogène de degré deux en les longueurs\,; la borne inférieure est prise sur papier}
\begin{proposition}[Longueurs entières et taille sans plafond]\label{prop:whole-size}
Si $P$ est donné par les équilibres et des bornes $\ell \ge b \ge 0$, un
ensemble fermable en longueurs réelles est fermable en longueurs entières, et
fermable en longueurs réelles à tout nombre de quads supérieur à $\inf_{P_Y}
Q$, et à $\inf_{P_Y} Q$ lui-même lorsque la borne inférieure est atteinte, où
$P_Y$ est l'ensemble des longueurs de $P$ qui ferment $Y$ et $Q(\ell)$ le
nombre de quads.
\end{proposition}
\end{evidence}

\begin{proof}
On multiplie un point rationnel de $P_Y$ par son dénominateur
(remarque~\ref{rem:exist}). Pour une taille $q$, prenons $\ell \in P_Y$ avec
$0 < Q(\ell) \le q$\,; $c\ell \in P_Y$ pour $c \ge 1$, et $Q(c\ell) = c^2
Q(\ell)$ atteint $q$.
\end{proof}

Avec un plafond d'écart, $P_Y$ est un polytope et $Q(P_Y)$ un intervalle~:
deux points exacts de part et d'autre d'une fenêtre de tailles certifient que
$Y$ se ferme à l'intérieur de celle-ci, tandis que la réciproque demande les
extrema de $Q$ sur un polytope. Le plafond est un budget partagé par tous les
tours~: sur le corps, les longueurs peuvent bouger de $980,7$ arêtes en tout,
alors que fermer les $109$ bandes de la variante~2 que les bornes ne forcent
pas coûte $1\,680,2$, et les cœurs deviennent des couvertures d'un sac à dos
(mesuré).

Sur une surface fermée où aucun patch ne se borde lui-même, un arc se trouve
sur un côté de chacun des deux patchs qu'il borde, si bien que sa colonne de
$B$ a deux coefficients $\pm1$~: $B$ est la matrice d'incidence d'un
\emph{graphe bidirigé} $G$ dont les sommets sont les équilibres et les arêtes
les arcs, un arc situé sur un bord donnant une demi-arête, le cadre
de~\cite{heistermann2023bimdf}, avec le rang de
Zaslavsky~\cite[\S8A]{zaslavsky1982signed}. Un cycle de $G$ est \emph{équilibré} si
des valeurs $\pm1$ sur ses arêtes équilibrent chacun de ses sommets, et
\emph{déséquilibré} sinon.

\begin{evidence}{core}{\lean{TMesh.handshake}, \lean{handshake}, \lean{no_whole_solution}\,; supposé~: une surface fermée}
\begin{lemma}[Poignées de main]\label{lem:handshake}
Soient $x \in \Z^n$ et $C$ une composante connexe de $G$. Alors $\sum_{r \in
C} (Bx)_r \equiv \sum_{a} x_a \pmod 2$, la seconde somme portant sur les arcs
de bord de $C$. Sur une surface fermée, $\sum_r (Bx)_r$ est pair.
\end{lemma}
\end{evidence}

\begin{proof}
Elle vaut $\sum_a x_a \sum_{r \in C} B_{ra}$, et chaque colonne a pour somme
$0$ ou $\pm2$, ou $\pm1$ pour un arc de bord.
\end{proof}

Ainsi, sur une surface fermée, un système $Bx = d$ avec $\sum_r d_r$ impair a
des solutions réelles mais aucune solution entière\,; cela n'affecte pas les
longueurs équilibrées~:

\begin{evidence}{core}{\lean{TMesh.free_arcs}, \lean{free_arcs}, \lean{free_arcs_of_cycle}\,; supposé~: le graphe est connexe (un arbre couvrant)}
\begin{proposition}[Arcs libres]\label{prop:free}
Supposons $G$ connexe et contenant un cycle déséquilibré. Soit $T$ l'ensemble
formé des arêtes d'un arbre couvrant de $G$ et d'une arête de plus qui ferme
un cycle déséquilibré, et $F$ l'ensemble des autres arêtes. Alors $B$ est de
rang plein en lignes, et tout $x_F \in \Z^F$ se prolonge en exactement un $x
\in \Z^n$ tel que $Bx = 0$. Les longueurs entières équilibrées forment donc un
réseau isomorphe à $\Z^F$, de dimension $n - m$ pour $m$ équilibres, qui a une
base de vecteurs $z_f$, $f \in F$, chacun supporté par $T \cup \{f\}$, avec
$z_f(f) = 1$ et des coefficients dans $\{0, \pm1, \pm2\}$.
\end{proposition}
\end{evidence}

\begin{proof}
$T$ existe~: le signe d'un cycle est un homomorphisme sur l'espace des cycles
sur $\Z/2$, qu'engendrent les cycles fondamentaux d'un arbre couvrant, si bien
que, s'ils étaient tous équilibrés, tout cycle le serait. $T$ a $m$ arêtes. Si
$B_T y = 0$, un sommet rencontré par une seule arête de $T$ force cette arête
à zéro\,; élaguer de telles arêtes laisse le cycle déséquilibré, sur lequel
chaque valeur fixe la suivante au signe près, et en faire le tour les force
toutes à zéro. Ainsi $B_T$ est inversible, $B$ est de rang plein en lignes,
et $x_F$ a exactement un prolongement réel. Pour $f \in F$, $T \cup \{f\}$ a
un espace de valeurs équilibrées de dimension un. Élagué, c'est soit trois
chemins entre deux sommets, dont les cycles sont le cycle déséquilibré de $T$,
un autre cycle déséquilibré et un cycle équilibré (le nombre de cycles
déséquilibrés d'un tel graphe est pair), le cycle équilibré contenant $f$,
puisque celui de $T$ ne le contient pas, et portant des valeurs $\pm1$\,; soit
deux cycles reliés par un chemin, éventuellement de longueur nulle, le cycle
qui passe par $f$ étant équilibré et portant $\pm1$, ou déséquilibré, et alors
$\pm1$ sur les deux cycles et $\pm2$ sur le chemin sont équilibrés. Dans les
deux cas, il existe un $z_f$ entier équilibré supporté par $T \cup \{f\}$ avec
$z_f(f) = \pm1$\,; les signes étant fixés, $x = \sum_{f \in F} x_f z_f$ est le
prolongement entier.
\end{proof}

Avec une forêt couvrante et, dans chaque composante qui a un cycle
déséquilibré, une arête qui en ferme un ou une demi-arête, la proposition
s'étend au-delà de la sphère fermée, là où le rang plein en lignes peut faire
défaut (sur part, $145$ sur $148$). Sur le corps masculin, $G$ est connexe
avec un cycle déséquilibré et le réseau des longueurs équilibrées est de
dimension $2\,069 - 1\,552 = 517$~: les équilibres seuls ne font jamais
obstacle aux longueurs entières, et l'élimination entière se termine sur une
seule ligne de plus grand commun diviseur $2$. Les décalages des tours retenus
laissent des réseaux de dimension $363$ à $469$, et dans les variantes 4, 5
et 6, l'élimination demande des étapes de l'algorithme d'Euclide~: les lignes
des tours peuvent détruire la structure d'arcs libres (mesuré). La seule
obstruction entière trouvée sur le corps est un coût d'arrondi en écart, non
une parité~: des anneaux qui se ferment en nombres réels à un écart moyen de
$0,399$ demandent au moins $0,536$ en nombres
entiers~\cite[\S7]{gardent2026loops}.

\subsection{Sur le corps masculin}

Dans le tableau~\ref{tab:forced}, $H$ est l'ensemble des tours retenus d'une
variante, ou celui des $111$ bandes du maillage du moteur pour la variante~2
inclinées de moins de $0,12$, et $P$ est $P_{\mathrm{eng}}$, puis
$P_{\mathrm{eng}}$ avec en plus le plafond du moteur $D = 0,474$ sur l'écart
moyen, puis avec en plus des longueurs entières et sa fenêtre de $10\,379$ à
$10\,801$ quads. Avec les seules bornes, la boucle d'ensemble transversal
implicite est exacte en au plus $66$ secondes par famille\,; avec le plafond
et la fenêtre, les encadrements viennent de sous-ensembles fermables, de
cœurs disjoints, de bornes du solveur et de points entiers vérifiés
exactement.

\begin{table}[t]
\begin{evidence}[unbreakable]{data}{la colonne $P_{\mathrm{eng}}$, exacte sur l\'export d\'où viennent ces nombres (\texttt{make lean-check})\,; les colonnes avec le plafond et la fenêtre sont mesurées, non vérifiées}
\centering
\small
\newcommand{\head}[1]{\begin{tabular}[b]{@{}r@{}}#1\end{tabular}}
\begin{tabular}{lrrrrr}
\toprule
variante & tours & bornes & \head{bornes,\\ plafond} & \head{plafond, taille,\\ entiers} &
\head{laissés décalés\\ par le moteur} \\
\midrule
0 & 48 & 1 & 1 & 1 & 13 \\
1 & 78 & 5 & 5 & 5--6 & 39 \\
2 & 52 & 1 & 1 & 1 & 16 \\
3 & 86 & 4 & 4 & 5--18 & 44 \\
4 & 84 & 1 & 1 & 1--12 & 49 \\
5 & 160 & 6 & 7--17 & 7--86 & 115 \\
6 & 88 & 1 & 1 & 1--17 & 44 \\
7 & 169 & 9 & 9--26 & 9--90 & 123 \\
\midrule
2, bandes du maillage & 111 & 2 & 3--11 & 3--43 & 81 \\
\bottomrule
\end{tabular}
\caption{Corps masculin~: tours de chaque famille que toute quantification de
$P$ laisse décalés (théorèmes~\ref{thm:forced} et~\ref{thm:transversal}\,;
mesuré), avec $P = P_{\mathrm{eng}}$, puis avec en plus le plafond du moteur
(longueurs réelles), puis avec en plus des longueurs entières et la fenêtre de
tailles du moteur. Un nombre seul est exact\,; un intervalle est certifié,
inférieurement par la borne prouvée du solveur ou par des cœurs disjoints,
supérieurement par un sous-ensemble fermable trouvé dans la limite de temps
(de $120$ à $1\,800$ secondes). Dernière colonne~: tours de la famille
décalés aux longueurs entières du moteur. Les tours sont comptés au signe
près.}
\label{tab:forced}
\end{evidence}
\end{table}

Avec les seules bornes, les tours forcés sont exactement $1$, $5$, $1$, $4$,
$1$, $6$, $1$ et $9$ (tableau~\ref{tab:forced})\,; dans les variantes 0, 2
et 4, le tour forcé est celui du poignet, dans la variante~6 l'un des tours $57$
et $58$, dont le conflit maintient l'arc $1241$ (le poignet y est retenu mais
peut être fermé)\,; dans la variante~7, neuf tours
sont nécessaires bien qu'au plus huit de ses $21$ cœurs soient disjoints, si
bien que la borne d'empilement n'est pas toujours atteinte. Avec le plafond
et la fenêtre, en longueurs entières, la valeur est encore exactement $1$ dans
les variantes 0 et 2~: des longueurs entières qui ferment tous les tours
retenus sauf celui du poignet existent à $10\,439$ et $10\,436$ quads, à des
écarts moyens de $0,442$ et $0,469$. Dans la variante~3, l'intégrité coûte au
moins un tour ($4$ en nombres réels, au moins $5$ en nombres entiers). Des
$16$ tours retenus de la variante~2 que le moteur laisse décalés, les bornes
en forcent $1$, et le plafond, la taille et l'intégrité aucun~: l'heuristique
est responsable des $15$ autres\,; dans les autres variantes, elle laisse au
moins $12$ à $37$ tours de plus que les meilleures longueurs trouvées. Le
plafond n'est celui du moteur qu'approximativement (il a refusé $0,4737$), et
des optimums qui ramènent d'autres arcs à zéro peuvent casser le maillage\,;
si l'on garde ouvert tout arc que le moteur garde ouvert, l'optimum pour la
variante~2 vaut $4$ (mesuré).

\section{Choisir les tours}\label{sec:choosing}

Appelons \RC{} le problème suivant~: étant donnés les équilibres $B$, des
bornes $b$, un ensemble $H$ de formes linéaires (les décalages) et un nombre
$k$, existe-t-il des longueurs dans $P = \{\ell : B\ell = 0,\ \ell \ge b\}$
qui ferment au moins $k$ formes de $H$\,? Le problème du sous-système
faisable maximum est NP-difficile et difficile à approcher en
général~\cite{amaldi1995complexity}, et le calcul d'une couverture minimale
des IIS est fortement NP-difficile déjà pour les flots dans les
réseaux~\cite{joormann2015flowiis}\,; la difficulté survit à la structure de
la quantification, un cas nouveau à notre connaissance.

\begin{evidence}{core}{la réduction~: \lean{closable_iff_independent}, \lean{tauP_eq_vertex_cover}\,; la NP-difficulté de Stable et de Couverture par sommets, et les bornes d'approximation, sont citées}
\begin{theorem}[Choisir les tours est NP-difficile]\label{thm:nphard}
\RC{} est NP-complet, en longueurs réelles ou entières, déjà lorsque $B$ est
la matrice d'incidence d'un graphe bidirigé, que les bornes valent $0$ ou $1$
et que les formes ont des coefficients $0$ ou $1$\,; il est fortement
NP-difficile, et reste NP-difficile avec un plafond d'écart, qui peut être
lâche. De plus, $\tau_P(H)$ peut être égal au nombre de couverture par
sommets (taille minimale d'une couverture par sommets) de n'importe quel
graphe, avec $|H|$ son nombre de sommets. Il ne peut donc pas être approché à
un facteur $\sqrt2 - \varepsilon$ près sauf si
P${}={}$NP~\cite{khot2017independent,dinur2018towards,dinur2018nonoptimally,khot2018grassmann},
ni à un facteur $2 - \varepsilon$ près si la conjecture des jeux uniques est
vraie~\cite{khot2008vertex}, et le plus grand sous-ensemble fermable de $H$ ne
peut pas être approché à un facteur $|H|^{1-\varepsilon}$ près sauf si
P${}={}$NP~\cite{hastad1999clique,zuckerman2007linear}.
\end{theorem}
\end{evidence}

\begin{proof}
Appartenance~: fermer un ensemble donné de $k$ formes est un programme
linéaire, et d'après la proposition~\ref{prop:whole-size}, la fermabilité en
réels et la fermabilité en entiers coïncident. Difficulté, à partir de
\textup{\textsc{Independent Set}} (ensemble indépendant)~: pour un graphe
$G$, prenons pour chaque arête $e = uv$ quatre arcs $y_{u,e}$, $y_{v,e}$,
$d_e$, $d'_e$ avec les équilibres
\[
  y_{u,e} + y_{v,e} - d_e = 0, \qquad y_{u,e} + y_{v,e} - d'_e = 0, \qquad
  d_e - d'_e = 0 ,
\]
dans lesquels chaque arc apparaît dans exactement deux lignes avec des
coefficients $\pm1$, les bornes $d_e, d'_e \ge 1$ et $y \ge 0$, et pour
chaque sommet $u$ la forme $\shift_u = \sum_{e \ni u} y_{u,e}$. Avec $y \ge
0$, $\shift_u = 0$ force $y_{u,e} = 0$ en $u$, si bien que fermer deux sommets
adjacents donne $d_e = 0 < 1$. Pour un ensemble indépendant $Y$, chaque arête
a au plus une extrémité dans $Y$~: mettons $y = 0$ de son côté et $1$ de
l'autre, ou, si aucune extrémité n'est dans $Y$, $1$ d'un côté quelconque et
$0$ de l'autre\,; alors $d_e = d'_e = 1$, et ce sont des longueurs entières de
$P$ qui ferment $Y$. Les ensembles fermables sont les ensembles indépendants
de $G$, et $\tau_P(H) = |V| - \alpha(G)$ est le nombre de couverture par
sommets. L'instance est linéaire en $G$, avec des nombres $0$ et $\pm1$\,; un
plafond assez grand ne contraint rien.
\end{proof}

La réduction a été vérifiée par ordinateur sur $40$ graphes aléatoires
(mesuré)\,; sa troisième équation d'équilibre met chaque arc dans deux
équilibres, comme dans un T-mesh. On ne sait pas qu'elle soit un T-mesh\,: ses
formes ont des coefficients positifs ou nuls et ne sont pas construites comme
les holonomies de marches droites\,; le théorème dit donc que l'algèbre de la
quantification permet des instances difficiles, non qu'une surface en
produise.

\begin{evidence}{open}{}
\begin{conjecture}[Instances réalisables]\label{conj:realize}
\RC{} reste NP-difficile lorsque $B$, $b$ et $H$ sont les équilibres, les
bornes de séparation et les tours d'un T-mesh d'une surface fermée de genre
zéro.
\end{conjecture}
\end{evidence}

Les indices en sa faveur sont que la couverture par sommets reste
NP-difficile sur les graphes planaires, et sur les graphes de petit
degré~\cite{garey1976simplified}, et que deux tours séparés par une
dislocation entrent en conflit comme les deux extrémités d'une arête.

Certaines parties sont faciles~: sans tours ni fenêtre, la quantification
entière est un problème de couplage polynomial~\cite{edmonds1970matching}\,;
avec des cœurs disjoints connus, $\tau_P(H)$ est leur nombre\,; avec des
formes de décalage de rang fixé, les ensembles fermables sont les plats d'un
arrangement, mais sur le corps les tours retenus d'une variante sont de rang
plein.

\begin{evidence}{open}{une esquisse}
\begin{remark}[La fenêtre de tailles\,; une esquisse]
La fenêtre sur le nombre de quads rend difficile, au sens faible, même le
fait de ne rien fermer. Prenons une réunion disjointe de tores, chacun formé
de deux patchs dont les côtés longs et courts $a_i$ et $c_i$ ont pour cibles
$N v_i$ et $1$, de sorte que $Q = \sum_i 2 a_i c_i$, et un plafond qui laisse les côtés courts $c_i$ grandir
de $T_0$ arêtes en tout. Pour $N$ grand, la fenêtre $Q = 2N(\sum_i v_i +
T_0)$ est respectée exactement lorsque $\sum_i v_i(c_i - 1) = T_0$ avec des
$c_i \ge 1$ entiers~: c'est le problème de la somme de sous-ensemble non
bornée (\emph{unbounded subset sum}, sac à dos entier), problème NP-complet
classique, avec
des nombres exponentiels en la taille de l'instance (les gadgets sont de genre
un et non connexes). On ne sait pas si la fenêtre seule est fortement
NP-difficile, avec $Q$ polynomial en la taille, ni quelle est la complexité de
la maximisation de $Q$ sur un polytope de longueurs réelles.
\end{remark}
\end{evidence}

\section{Ce qui est vérifié par machine}\label{sec:lean}

Le dossier \texttt{lean} à côté de cette note est un projet
Lean~4~\cite{demoura2021lean4} sur Mathlib~\cite{mathlib2020} (tous deux en
\texttt{v4.35.0}). Ses $480$ déclarations de théorèmes, dont beaucoup de
petits lemmes, ne dépendent que des axiomes standard de Lean\,; ce nombre
n'est pas un argument, et son \texttt{README.md} associe à chaque énoncé ses
théorèmes Lean. Lean vérifie le cœur algébrique et combinatoire de chaque
énoncé, les faits topologiques qu'il utilise (qu'un lacet se développe en un
déplacement donné, qu'une courbe borde un disque, que les lacets $\lambda_i$
présentent le groupe fondamental) étant des hypothèses explicites. Les fichiers partagent les
déplacements et le groupe des quarts de tour (\texttt{Basic.lean}). L'un
d'eux, \texttt{TMesh.lean}, définit un T-mesh combinatoire~: des patchs à
deux paires de côtés opposés, chaque côté la liste de ses arcs, chaque arc sur
exactement deux côtés avec un signe pour le recollement. La matrice des
équilibres $B$ et les décalages des tours en sont déduits, et sur lui la
proposition~\ref{prop:linear}, le lemme~\ref{lem:handshake}, la
proposition~\ref{prop:free}, la remarque~\ref{rem:definite}, le
théorème~\ref{thm:forced} et, par le maillage de quads qu'il définit aux
longueurs entières, la proposition~\ref{prop:parity} forment une seule
chaîne vérifiée. Les autres
énoncés gardent leur propre abstraction (un tour comme une liste de
déplacements, une région comme des ensembles finis avec des comptes)\,; la
surface et le champ de croix sous le T-mesh, son orientation et le lien entre
le décalage et l'holonomie $\hat h$ ne sont pas formalisés.

\ifdefined\plain
\emph{Vérifié}, énoncé par énoncé. Sur le T-mesh, les repères des patchs
étant ceux d'une orientation de la surface~: la
proposition~\ref{prop:linear} pour l'abscisse transverse à la marche,
c'est-à-dire les coefficients entiers de $\shift_\gamma$, indépendants de
$\ell$, le couloir découpé par des inégalités linéaires en l'abscisse et en
$\ell$, et la marche qui revient en $o + \shift_\gamma(\ell)$ et s'enroule en
hélice\,; que cette marche depuis un demi-entier soit une ligne duale du
maillage de quads est la proposition de complétude ci-dessous, dans son propre
modèle. Pour des déplacements dont les parties
linéaires sont dans un anneau quelconque~: la somme de Burgers, et le fait que
la récurrence des applications de recollement d'un arbre de coupures, une
hypothèse, se déroule dans l'ordre de première visite
(proposition~\ref{prop:burgers}(a))\,; la couronne découpée~(b)\,; les
échanges\,; l'étape de conjugaison de la proposition~\ref{prop:winding}. Sur
les réels~: les quatre ordres et orientations d'une dislocation et les
chaînes droites (proposition~\ref{prop:dislocation}(i), (ii))\,; les
escaliers\,; la positivité d'une forme à coefficients positifs ou nuls avec un
arc maintenu. Aux longueurs entières~: la règle du damier\,; la
proposition~\ref{prop:parity} sur le T-mesh, complète~: le graphe des quads
d'un T-mesh connexe à $P - A + N = 2$ est biparti, par algèbre linéaire sur
$\Z/2$ sur les chaînes de son maillage de quads, étant seulement donné qu'une
station suit le tour à toutes longueurs équilibrées strictement positives\,;
la caractéristique d'Euler du maillage de quads aux longueurs entières, $P -
A + N$ pour $P$ patchs, $A$ arcs et $N$ nœuds, indépendante des longueurs\,; la seconde
démonstration de la remarque~\ref{rem:borders}. La formule de Fox--Burgers
sur le groupe libre avec les enroulements tordus et l'insertion de la
relation, le cas d'un tour simple, la proposition~\ref{prop:commutator} et le
corollaire~\ref{cor:checker}. Le lemme~\ref{lem:handshake} sur une surface
fermée, les colonnes de $B$ de tout T-mesh fermé étant de somme paire\,; la
proposition~\ref{prop:free}, $B$ étant la matrice d'incidence d'un graphe
bidirigé quand aucun arc n'est sur deux côtés opposés d'un patch, et sous une forme plus forte
qu'énoncée (sur les arêtes $T$, l'image par $B$ des vecteurs entiers est le
réseau des vecteurs de somme paire)\,; les remarques~\ref{rem:definite}
et~\ref{rem:exist}, avec le lemme de Farkas sur un corps ordonné quelconque et
des multiplicateurs entiers\,; le théorème~\ref{thm:forced} et le
corollaire~\ref{cor:chain}. Pour des champs d'ordre quelconque~: la partie
rotation de Gauss--Bonnet et ses comptes sur les maillages triangulaires et
hexagonaux, les longueurs des dislocations et la congruence de
l'annexe~\ref{app:orders}. Ajoutés après la première version~: la formule de
Fox--Burgers pour des générateurs quelconques et le commutateur de deux
déplacements dont les rotations commutent (propositions~\ref{prop:burgers-any}
et~\ref{prop:torus}), les colliers, l'identité des peignes, la mise à l'échelle
de la proposition~\ref{prop:whole-size}, les conséquences de la parité
(corollaire~\ref{cor:parity}), l'homomorphisme de parité du
théorème~\ref{thm:parity-any} (sur les déplacements à rotation d'un quart de
tour et translation entière, nul sur les lacets autour des points coniques,
$L$ sur une composante de bord, $\alpha + s$ sur un tour), et la réduction du
théorème~\ref{thm:nphard}~: sur le gadget d'un graphe, les ensembles
fermables sont les stables et $\tau_P$ est le nombre de couverture par
sommets.

\emph{Vérifié sous une forme plus faible}~: les
lemmes~\ref{lem:gauss} et~\ref{lem:eight}, seulement le compte par la
caractéristique d'Euler sur la grille raffinée par deux et les totaux à partir
des nombres de coins convexes et concaves, le fait que les régions soient des
disques et la convexité de leurs coins étant des hypothèses\,; la complétude
des marches, proposition~2 de~\cite{gardent2026loops}, seulement dans un
modèle~: sur une surface fermée (côtés opposés d'un patch de même longueur,
deux arêtes unité jamais recollées sur la même) la marche des stations définie
par les données de recollement est une permutation, les bandes sont
\emph{définies} comme les cycles du successeur qu'elle donne sur les quads
au-dessus des stations, et elles sont en bijection avec les cycles de la
marche réelle depuis les demi-entiers $k + \frac12$, dans un modèle
unidimensionnel du recollement\,; que les côtés des patchs soient strictement
positifs, que la traversée d'un patch garde l'abscisse et que le modèle soit
le maillage de quads sont des hypothèses. 
\else
Ce que Lean vérifie de chaque énoncé, et ce qu'il suppose, est l'étiquette de
l'énoncé.
\fi
\emph{Non vérifié}~: la topologie prise comme hypothèse, les
bandes avec bords, les énoncés mesurés, et ceux qui ne sont démontrés que sur
papier (voir Assistance par IA). La formalisation n'a trouvé aucune erreur de
signe ou de convention\,; elle a précisé les hypothèses des
propositions~\ref{prop:parity}, \ref{prop:free} et~\ref{prop:staircase}, de
la remarque~\ref{rem:definite} et de la proposition de complétude, énoncées
ici sous leur forme précisée.

\section{Problèmes ouverts et limites}\label{sec:open}

Quand un ensemble de tours est-il fermable, en termes de configuration des
singularités\,? Tout certificat est une somme de peignes, ce qui ne dit rien
des singularités, et n'est pas toujours une somme d'escaliers, de calottes et
de séparatrices qui s'enroulent seuls\,; une réponse utile se trouve entre les
deux, par exemple des peignes portés par des ensembles disjoints de sommets,
chacun du signe du certificat. Calculer $\tau_P(H)$ à un nombre de quads donné
demande les extrema d'une fonction quadratique indéfinie sur un polytope.
\RC{} est-il NP-difficile pour un nombre fixé de sommets singuliers\,? Et
lorsque l'ensemble des tours dépend des longueurs, comme c'est le cas des
bandes du maillage, quel est au juste le problème\,? Quand existe-t-il, pour
une configuration de singularités donnée, une différentielle quartique dont
toutes les trajectoires sont finies~\cite{lei2020quartic}\,? Les tours en sont
la version discrète. Enfin, un modèle formel d'un complexe de rectangles à
côtés entiers, de ses tours et de son application développante ferait des
hypothèses topologiques des énoncés Lean des théorèmes.

\paragraph{Limites et reproductibilité.}
Les énoncés mesurés viennent d'un seul T-mesh d'un seul corps, exporté par un
moteur sans les passes de contraintes de l'article compagnon, et de quelques
modèles du banc avec une version du moteur non enregistrée\,; les identités
ont été vérifiées en virgule flottante sauf mention d'arithmétique exacte, et
les homotopies des propositions~\ref{prop:swap} et~\ref{prop:winding} n'ont
pas été vérifiées. Les comptes forcés sont relatifs aux bornes de séparation
du moteur et, approximativement, à son plafond, et les intervalles dépendent
des limites de temps du solveur sur une machine partagée. Les scripts et les
journaux des expériences exactes ne sont pas publiés~: les énoncés mesurés
peuvent être confrontés à la branche \texttt{rd-exact}, un point d'accroche
expérimental d'openclare, mais pas réexécutés tels quels. Les preuves Lean se
rejouent avec \texttt{make proofs NOTE=turns-of-a-t-mesh}. Les tests de
propriétés du paquet de données (\texttt{experiments/property}) fabriquent
$50$ T-mesh aléatoires sans le moteur à graine fixée, $35$ sphères et $15$
tores, fermés, et vérifient sur eux les définitions Lean du T-mesh et de ses
tours, chaque énoncé dont ils peuvent calculer l'hypothèse, et les certificats
de leurs conflits sous les bornes $\ell \ge 1$\,; les étiquettes disent où
ils ont tourné, et un test corrobore ou réfute mais ne démontre rien. Cette
version de la
note a été construite à partir du commit \texttt{\enginecommit}.

\section*{Assistance par IA}
Cette note, ses preuves Lean et le logiciel qu'elle décrit ont été faits par
un auteur humain travaillant avec un système d'IA, qui n'est pas un auteur.
L'auteur a choisi le problème et ses objectifs (des boucles d'arêtes qui se
ferment en anneaux et, en artiste, peu de sommets singuliers et peu de
spirales plutôt que des quads parfaitement carrés), a dirigé le travail, a
décidé quels résultats poursuivre, et les a jugés au travers des mesures et
des brouillons qui lui ont été présentés. Claude Opus 5.5 (Anthropic), utilisé
via Claude Code en octobre 2026 avec plusieurs sous-agents, a écrit les
changements du moteur et les scripts d'expérience, a mené les exports et les
calculs exacts, a fait la recherche bibliographique, a énoncé et démontré les
propositions, a écrit la formalisation Lean et a rédigé le texte. La note de
chaque entrée de la bibliographie dit quelle part de l'œuvre a été lue. Une œuvre dont seul le résumé, ou un compte rendu, a été lu n'est citée que pour ce que ce résumé ou ce compte rendu dit.
Au-delà des cœurs vérifiés par machine en Lean (section~\ref{sec:lean}), tout
a été écrit avec l'aide de l'IA et vérifié par des agents d'IA qui l'ont relu
et testé sur les données, non par un expert humain. En particulier, les
propositions~\ref{thm:combs} et~\ref{prop:pairing}, le
lemme~\ref{lem:gauss-border}, le théorème~\ref{thm:transversal}, le cas avec
bords du lemme~\ref{lem:handshake}, et les énoncés sur les rotations d'ordre
six autres que ceux qui sont vérifiés n'ont qu'une démonstration sur papier
écrite avec l'aide de l'IA\,; les propositions~\ref{prop:collar},
\ref{prop:burgers-any}, \ref{prop:torus} et~\ref{prop:whole-size}, le
lemme~\ref{lem:comb}, le corollaire~\ref{cor:parity}, le
théorème~\ref{thm:parity-any} et la réduction du théorème~\ref{thm:nphard}
ont eu leur cœur algébrique vérifié en Lean après la première version de
cette note. La relecture par IA a trouvé des erreurs
dans les propositions~\ref{thm:combs} et~\ref{prop:pairing} et le
théorème~\ref{thm:nphard}, ainsi que dans la proposition~\ref{prop:linear} et
la proposition de complétude de~\cite{gardent2026loops}, dont les
vérifications Lean sont plus faibles que les énoncés\,; toutes sont corrigées
ici. La note n'a pas été relue par des pairs\,; l'auteur, qui la publie, est
responsable de son contenu. La présente traduction française a elle aussi été
faite avec l'aide de l'IA\,; la version anglaise fait foi.

\bibliographystyle{plainnat}
\bibliography{refs}

\appendix

\section{Rotations d'autres ordres}\label{app:orders}

Sur les maillages triangulaires et hexagonaux, les rotations sont des
puissances de $\zeta = e^{i\pi/3}$, un sommet de valence $k$ a pour indice $6
- k$, et un tour selon $u$ revient translaté de $T = \alpha u + \shift\, \zeta
u$. La somme de Burgers, l'ordre de première visite et la formule de
Fox--Burgers n'utilisent pas l'ordre des rotations et valent sans changement,
de même que le théorème~\ref{thm:forced}, le corollaire~\ref{cor:chain} et le
lemme~\ref{lem:handshake}, qui n'utilisent aucune géométrie. Une paire de
valences cinq et sept a $B = \pm\zeta^{\pm1} d$, et les décalages de deux
tours diffèrent de $\pm d_1$ ou de $\pm(d_1 + d_2)$, avec $d = d_1 u + d_2
\zeta u$. Ce qui change, c'est la parité. Lorsque toute rotation est une
puissance $z^{j_i}$ d'une même rotation $z$, toute dérivée de Fox est
congrue à sa somme d'exposants modulo $1 - z$, et $1 - z^j$ à $j(1 - z)$
modulo $(1 - z)^2$, si bien que, dans tout anneau commutatif,
\[
  \tau(w) \equiv (1 - z) \sum_i \epsilon\Bigl(\frac{\partial w}{\partial
  x_i}\Bigr) j_i\, p_i \pmod{(1 - z)^2}.
\]
Pour les quarts de tour, $(1 - z)^2 = -2z$, et c'est le
corollaire~\ref{cor:checker}. Pour l'ordre six, $1 - z$ est inversible~: il
n'y a pas de règle de coloriage sur les maillages triangulaires ou
hexagonaux. Lorsque tout sommet singulier a une valence paire, les rotations
sont d'ordre trois, $(1 - \omega)^2 = -3\omega$ (comparer au lemme~11
de~\cite{izmestiev2013torus}), et sur une sphère fermée $\alpha \equiv \shift
\pmod 3$~: tout anneau a une longueur divisible par trois, comme sur les
octaèdres subdivisés, alors que, sur les icosaèdres subdivisés, des anneaux de
$5n$ arêtes prennent tous les résidus (mesuré)\,; une règle modulo trois avec
bords ou anses est une conjecture. Une dislocation d'ordre $n$ a $|B| =
2\sin(\pi/n)\,|d|$.

\end{document}
